% !TEX program = xelatex
\documentclass[11pt,a4paper]{report}

\usepackage{fontspec}
\usepackage[italian]{babel}
\usepackage[a4paper,margin=2.6cm,headheight=14pt]{geometry}
\usepackage{xcolor}
\usepackage{amsmath}
\usepackage{booktabs}
\usepackage{enumitem}
\usepackage{microtype}
\usepackage{fancyhdr}
\usepackage{listings}
\usepackage[most]{tcolorbox}
\usepackage{titlesec}
\usepackage{array}
\usepackage{longtable}
\newcolumntype{L}[1]{>{\raggedright\arraybackslash}p{#1}}
\setcounter{tocdepth}{1}

% ---------------------------------------------------------------------------
% Supporto ai caratteri CJK (usati solo da "Hello, 世界" della slide 8).
% Se nel sistema non c'è un font CJK noto, si ripiega su una sigla testuale,
% così il documento compila ovunque.
% ---------------------------------------------------------------------------
\newif\ifhascjk
\hascjkfalse
\def\cjkname{}
\newcommand{\trycjk}[1]{%
  \ifhascjk\else\IfFontExistsTF{#1}{\def\cjkname{#1}\hascjktrue}{}\fi}
\trycjk{Microsoft YaHei}
\trycjk{PingFang SC}
\trycjk{Noto Sans CJK SC}
\trycjk{Source Han Sans SC}
\trycjk{WenQuanYi Zen Hei}
\trycjk{SimSun}
\ifhascjk
  \usepackage{xeCJK}
  \setCJKmainfont{\cjkname}
  \setCJKsansfont{\cjkname}
  \setCJKmonofont{\cjkname}
\fi

% ---------------------------------------------------------------------------
% Colori
% ---------------------------------------------------------------------------
\definecolor{gokw}{HTML}{9B2393}
\definecolor{gotype}{HTML}{1C6E8C}
\definecolor{gostr}{HTML}{B3261E}
\definecolor{gocom}{HTML}{4F7A28}
\definecolor{gonum}{HTML}{8A5A00}
\definecolor{codebg}{HTML}{F7F7F9}
\definecolor{outbg}{HTML}{FFF8E1}
\definecolor{defcol}{HTML}{1F4E9C}
\definecolor{notecol}{HTML}{2E7D32}
\definecolor{warncol}{HTML}{B3261E}
\definecolor{cmpcol}{HTML}{6A3FA0}

% ---------------------------------------------------------------------------
% listings: linguaggio Go
% ---------------------------------------------------------------------------
\lstdefinelanguage{Go}{
  morekeywords={break,case,chan,const,continue,default,defer,else,fallthrough,
    for,func,go,goto,if,import,interface,map,package,range,return,select,
    struct,switch,type,var},
  morekeywords=[2]{bool,byte,complex128,error,float32,float64,int,int8,int16,
    int32,int64,rune,string,uint,uint8,uint16,uint32,uint64,uintptr,any},
  morekeywords=[3]{true,false,nil,iota},
  morekeywords=[4]{append,cap,close,copy,delete,len,make,new,panic,print,
    println,recover},
  sensitive=true,
  morecomment=[l]{//},
  morecomment=[s]{/*}{*/},
  morestring=[b]",
  morestring=[b]`,
  morestring=[b]'
}

\lstdefinestyle{gocode}{
  language=Go,
  basicstyle=\ttfamily\small,
  keywordstyle=\color{gokw}\bfseries,
  keywordstyle=[2]\color{gotype},
  keywordstyle=[3]\color{gonum},
  keywordstyle=[4]\color{gotype}\bfseries,
  commentstyle=\color{gocom}\itshape,
  stringstyle=\color{gostr},
  numbers=left,
  numberstyle=\tiny\color{gray},
  numbersep=7pt,
  backgroundcolor=\color{codebg},
  frame=single,
  rulecolor=\color{black!25},
  framesep=4pt,
  xleftmargin=14pt,
  framexleftmargin=14pt,
  breaklines=true,
  breakatwhitespace=false,
  postbreak=\mbox{\textcolor{gray}{$\hookrightarrow$}\space},
  columns=fullflexible,
  keepspaces=true,
  showstringspaces=false,
  tabsize=4,
  captionpos=t,
  aboveskip=1.2em,
  belowskip=0.8em,
  upquote=true
}

\lstdefinestyle{output}{
  basicstyle=\ttfamily\small,
  backgroundcolor=\color{outbg},
  frame=leftline,
  rulecolor=\color{orange!70!black},
  framesep=6pt,
  xleftmargin=6pt,
  breaklines=true,
  columns=fullflexible,
  keepspaces=true,
  showstringspaces=false,
  aboveskip=0.6em,
  belowskip=1em
}

\lstdefinestyle{shell}{
  basicstyle=\ttfamily\small,
  backgroundcolor=\color{black!5},
  frame=leftline,
  rulecolor=\color{black!50},
  framesep=6pt,
  xleftmargin=6pt,
  breaklines=true,
  columns=fullflexible,
  keepspaces=true,
  showstringspaces=false,
  aboveskip=0.6em,
  belowskip=1em
}

\renewcommand{\lstlistingname}{Listato}

% Codice in linea
\newcommand{\cod}[1]{\texttt{#1}}

% ---------------------------------------------------------------------------
% Box didattici
% ---------------------------------------------------------------------------
\tcbset{
  boxstyle/.style={enhanced, breakable, boxrule=0.6pt, arc=2pt,
    left=8pt,right=8pt,top=5pt,bottom=5pt,
    fonttitle=\bfseries\sffamily\small, before skip=0.9em, after skip=0.9em}
}
\newtcolorbox{definizione}[1]{boxstyle, colback=defcol!6, colframe=defcol,
  title={Definizione --- #1}}
\newtcolorbox{nota}[1][Osservazione]{boxstyle, colback=notecol!6, colframe=notecol,
  title={#1}}
\newtcolorbox{attenzione}[1][Errore comune]{boxstyle, colback=warncol!6, colframe=warncol,
  title={#1}}
\newtcolorbox{confronto}[1]{boxstyle, colback=cmpcol!6, colframe=cmpcol,
  title={Confronto --- #1}}

% ---------------------------------------------------------------------------
% Intestazioni e struttura
% ---------------------------------------------------------------------------
\titleformat{\chapter}[display]
  {\normalfont\sffamily\bfseries\Huge\color{defcol}}
  {\chaptertitlename\ \thechapter}{12pt}{\Huge}
\titlespacing*{\chapter}{0pt}{-20pt}{24pt}
\titleformat{\section}{\normalfont\sffamily\bfseries\Large\color{defcol!80!black}}{\thesection}{0.6em}{}
\titleformat{\subsection}{\normalfont\sffamily\bfseries\large}{\thesubsection}{0.6em}{}
\titleformat{\subsubsection}{\normalfont\sffamily\bfseries\normalsize\itshape}{}{0pt}{}

\pagestyle{fancy}
\fancyhf{}
\fancyhead[L]{\small\sffamily Guida a Go}
\fancyhead[R]{\small\sffamily\leftmark}
\fancyfoot[C]{\small\thepage}
\renewcommand{\headrulewidth}{0.3pt}
\fancypagestyle{plain}{\fancyhf{}\fancyfoot[C]{\small\thepage}\renewcommand{\headrulewidth}{0pt}}

% Riduce gli overfull nei paragrafi con molti nomi di file.
\emergencystretch=3em
\newcommand{\fn}[1]{\texttt{#1}}

\setlength{\parindent}{0pt}
\setlength{\parskip}{0.55em}
\setlist{itemsep=0.15em, topsep=0.3em}

\usepackage[hidelinks,pdfusetitle,bookmarksnumbered]{hyperref}
\hypersetup{pdftitle={Guida a Go},pdfsubject={Guida tecnica e didattica al linguaggio Go},pdflang={it}}

\begin{document}

% ===========================================================================
% FRONTESPIZIO
% ===========================================================================
\begin{titlepage}
\centering
\vspace*{2.2cm}
{\sffamily\bfseries\fontsize{40}{46}\selectfont\color{defcol} Guida a Go\par}
\vspace{0.6cm}
{\sffamily\Large Dal programma «Hello, World» alla concorrenza\par}
\vspace{1.2cm}
\rule{0.6\textwidth}{0.6pt}\par
\vspace{0.8cm}
{\large Guida tecnica e didattica al linguaggio Go\\[0.3em]
per studenti di Ingegneria Informatica\par}
\vspace{2.2cm}
\begin{minipage}{0.8\textwidth}
\centering
\small
Il contenuto è ricavato dalle slide \emph{Introduction to Go}\\
del corso di \emph{Sistemi Distribuiti e Cloud Computing}, A.A.\ 2025/26,\\
Laurea Magistrale in Ingegneria Informatica, Università di Roma Tor Vergata\\
(autrice delle slide: Valeria Cardellini).
\end{minipage}
\vfill
{\small Gli esempi di codice di questa guida si trovano nella cartella \texttt{examples/}\\
e sono stati compilati ed eseguiti con Go 1.27.3.\par}
\end{titlepage}

\pagenumbering{roman}
\tableofcontents
\clearpage
\pagenumbering{arabic}

% ===========================================================================
% INTRODUZIONE
% ===========================================================================
\chapter*{Introduzione e obiettivi}
\addcontentsline{toc}{chapter}{Introduzione e obiettivi}
\markboth{Introduzione}{Introduzione}

\section*{A chi si rivolge e che cosa contiene}

Questa guida trasforma le slide del corso in un testo organico per lo studio
individuale. Presuppone che il lettore conosca già la programmazione e abbia
familiarità con C, Python e Java: i concetti generali (variabili, cicli,
funzioni, ricorsione, \ldots) non vengono rispiegati. L'attenzione è invece sulle
scelte specifiche di Go e su come si collocano rispetto ai linguaggi già noti.

Gli argomenti sono esattamente quelli delle slide, in una successione
leggermente riordinata per rendere l'apprendimento più naturale
(per esempio i moduli sono trattati insieme alla toolchain, e la gestione degli
errori insieme alle interfacce). Sono stati omessi, perché non contribuiscono
all'apprendimento tecnico, i cenni storici, le classifiche di popolarità dei
linguaggi e i riferimenti bibliografici o web. Quando una slide accenna soltanto
a un argomento (per esempio RPC o il package \cod{net/http}), la guida lo
segnala ma non lo sviluppa.

\section*{Come è organizzata la guida}

\begin{center}
\small
\begin{tabular}{@{}rll@{}}
\toprule
\textbf{Cap.} & \textbf{Argomento} & \textbf{Cartella degli esempi}\\
\midrule
1 & Il linguaggio Go in sintesi & ---\\
2 & Struttura del programma, package, toolchain, moduli & \texttt{01\_basics}\\
3 & Variabili e tipi & \texttt{01\_basics}\\
4 & Funzioni, closure, funzioni variadiche & \texttt{02\_functions}\\
5 & Controllo di flusso e \cod{defer} & \texttt{03\_control\_flow}\\
6 & Puntatori e strutture dati & \texttt{04\_data\_structures}\\
7 & Metodi, interfacce, errori & \texttt{05\_methods\_interfaces}\\
8 & Concorrenza & \texttt{06\_concurrency}\\
9 & Test e benchmark & \texttt{07\_testing\_and\_tools}\\
\bottomrule
\end{tabular}
\end{center}

Le due appendici raccolgono la mappa tra slide, sezioni ed esempi (Appendice~A) e
le note della verifica incrociata, comprese le correzioni a errori trovati nelle
slide (Appendice~B).

\section*{Come usare gli esempi}

Ogni listato della guida riporta il nome del file corrispondente nella cartella
\texttt{examples/}. Il testo dei listati è incluso \emph{automaticamente} dai file
\texttt{.go} reali, quindi guida ed esempi non possono divergere. Anche gli output
mostrati sono stati ottenuti eseguendo i programmi; quando l'output non è
deterministico, il testo lo dice esplicitamente.

Ogni file è un programma indipendente (\cod{package main}) da eseguire dalla
cartella \texttt{examples/}, che contiene un file \texttt{go.mod} condiviso:

\begin{lstlisting}[style=shell]
cd examples
go run 01_basics/helloworld.go
go run 03_control_flow/echo.go Hello Go World
\end{lstlisting}

Poiché ogni cartella contiene più programmi con un proprio \cod{main}, non si
usa \cod{go build ./...} sull'intera cartella (il compilatore segnalerebbe
\cod{main redeclared}): si esegue un file alla volta. L'unica eccezione è la
cartella \texttt{07\_testing\_and\_tools}, che è un package ordinario testato con
\cod{go test} (Capitolo~\ref{ch:test}). Serve Go 1.22 o successivo (si veda la
nota sulle variabili di ciclo nel Capitolo~\ref{ch:conc}).

\section*{Convenzioni grafiche}

\begin{definizione}{esempio}
Un box blu introduce una definizione precisa di un concetto.
\end{definizione}
\begin{nota}
Un box verde raccoglie un'osservazione importante da ricordare.
\end{nota}
\begin{attenzione}
Un box rosso descrive un errore frequente, spiegando \emph{perché} avviene e come
evitarlo.
\end{attenzione}
\begin{confronto}{C, Python, Java}
Un box viola mette a confronto il costrutto Go con quelli dei linguaggi noti.
\end{confronto}

I blocchi con sfondo giallo e barra arancione riportano l'output reale dei programmi.

% ===========================================================================
\chapter{Il linguaggio Go in sintesi}\label{ch:go}
% ===========================================================================

Prima di scrivere codice conviene capire \emph{che tipo di linguaggio} è Go e
quali problemi i suoi progettisti hanno voluto risolvere. Molte scelte che a
prima vista sembrano restrizioni (niente eccezioni, niente classi, niente
aritmetica dei puntatori, regole di formattazione rigide) sono conseguenze
dirette di questi obiettivi.

\section{Obiettivi di progetto}

Go è definito come un linguaggio open source che rende facile costruire sistemi
\emph{semplici}, \emph{sicuri} e \emph{scalabili}. In termini più tecnici, gli
obiettivi dichiarati sono tre:

\begin{enumerate}
  \item essere \textbf{espressivo ed efficiente} sia in compilazione sia in
        esecuzione;
  \item favorire la scrittura di programmi \textbf{affidabili e robusti};
  \item offrire un linguaggio \textbf{staticamente tipato e compilato
        rapidamente}, ma che \emph{si usi} come un linguaggio a tipizzazione
        dinamica e interpretato.
\end{enumerate}

L'ultimo punto è il più interessante. Un programmatore Python scrive poche
righe e ottiene subito un risultato, ma scopre molti errori solo a runtime; un
programmatore C ottiene prestazioni elevate ma deve annotare tutto e gestire
manualmente la memoria. Go cerca un compromesso: ogni variabile ha un tipo
statico verificato dal compilatore, ma l'\emph{inferenza di tipo} (Capitolo~\ref{ch:tipi})
evita di ripeterlo continuamente, e la compilazione è abbastanza veloce da
rendere il ciclo modifica--esecuzione simile a quello di un linguaggio interpretato
(\cod{go run}, Sezione~\ref{sec:toolchain}).

\section{Che cosa Go eredita da C}

Go è spesso descritto come un linguaggio «C-like» o «il C per il ventunesimo
secolo». Dal C eredita:

\begin{itemize}
  \item la \textbf{sintassi delle espressioni};
  \item le \textbf{istruzioni di controllo del flusso} (\cod{for}, \cod{if}, \cod{switch});
  \item i \textbf{tipi di dato di base};
  \item il \textbf{passaggio dei parametri per valore} (\emph{call-by-value}):
        ogni argomento è copiato nel parametro formale;
  \item i \textbf{puntatori};
  \item l'\textbf{efficienza a runtime};
  \item la \textbf{tipizzazione statica}.
\end{itemize}

Per la concorrenza, invece, l'antenato dichiarato è il modello \textbf{CSP}
(\emph{Communicating Sequential Processes}, formulato da Tony Hoare): processi
sequenziali indipendenti che cooperano scambiandosi messaggi anziché condividendo
memoria. Questo modello è implementato da goroutine e channel
(Capitolo~\ref{ch:conc}).

\section{Le novità rispetto a C}

Le slide raggruppano le funzionalità moderne di Go in quattro aree.

\begin{description}
  \item[Concorrenza.] Strumenti nuovi ed efficienti per la programmazione
        concorrente, parte del linguaggio e non di una libreria esterna.
  \item[Astrazione dei dati e programmazione a oggetti.] Un approccio flessibile
        basato sulla \emph{composizione} piuttosto che sull'ereditarietà. Go non ha
        classi: il comportamento si associa a un tipo con i \emph{metodi}
        e si astrae con le \emph{interfacce} (Capitolo~\ref{ch:oop}).
  \item[Gestione della memoria.] Gestione automatica con \emph{garbage
        collection}: il programmatore non libera esplicitamente la memoria.
  \item[Leggibilità e usabilità.] Sintassi pulita e concisa, pensata per codice
        manutenibile e leggibile.
\end{description}

\begin{confronto}{gestione della memoria}
In \textbf{C} la memoria dinamica si alloca con \cod{malloc} e si libera con
\cod{free}; errori come il \emph{use-after-free} o i \emph{memory leak} sono a
carico del programmatore. \textbf{Java} e \textbf{Python} hanno anch'essi un
garbage collector (Python in CPython combina conteggio dei riferimenti e
rilevazione dei cicli). \textbf{Go} si colloca dalla parte di Java e Python:
il runtime libera automaticamente gli oggetti non più raggiungibili, eliminando
in particolare i problemi di use-after-free. A differenza di Java e Python, però,
Go espone i puntatori (senza aritmetica, Sezione~\ref{sec:puntatori}) e produce
codice nativo compilato, senza macchina virtuale né interprete.
\end{confronto}

\section{Go per i sistemi distribuiti e il cloud}

Le slide motivano Go nel contesto dei sistemi distribuiti: il linguaggio permette
di concentrarsi sui problemi del distribuito grazie a

\begin{itemize}
  \item un buon supporto alla \textbf{concorrenza};
  \item un buon supporto alle \textbf{RPC} (\emph{Remote Procedure Call});
  \item il \textbf{garbage collector}, che evita i problemi di \emph{use-after-free};
  \item la \textbf{sicurezza dei tipi} (\emph{type safety});
  \item la facilità di apprendimento e manutenzione.
\end{itemize}

Go è inoltre sempre più usato per le applicazioni \emph{cloud-native}. A titolo di
esempio, le slide citano la libreria \emph{Go Cloud}, che fornisce astrazioni sui
servizi dei diversi fornitori cloud: lo stesso codice può leggere un oggetto da
un \emph{blob storage} indipendentemente dal provider sottostante, cambiando solo
l'URL del bucket. Il frammento seguente riproduce quello della slide.

\begin{lstlisting}[style=gocode,caption={Frammento della slide sul Go Cloud (non eseguibile in isolamento)}]
ctx := context.Background()
bucket, err := blob.OpenBucket(ctx, "s3://my-bucket")
if err != nil {
	return err
}
defer bucket.Close()
blobReader, err := bucket.NewReader(ctx, "my-blob", nil)
if err != nil {
	return err
}
\end{lstlisting}

\begin{nota}[Perché questo frammento non ha un file in \texttt{examples/}]
Il codice richiede la libreria esterna \emph{Go Cloud}, credenziali del provider
e un bucket esistente: non è eseguibile né verificabile in modo autonomo, e
usarlo come esempio introdurrebbe argomenti esterni al programma. Il frammento
è utile soprattutto per riconoscere i pattern già visti: la
coppia risultato/errore (Sezione~\ref{sec:errori}), il controllo \cod{err != nil}
e l'istruzione \cod{defer} per la pulizia (Sezione~\ref{sec:defer}).
\end{nota}

\section{Ambiente di sviluppo}

Go si può scrivere con un qualsiasi editor di testo; gli ambienti indicati dalle
slide sono \textbf{GoLand} (IDE di JetBrains), \textbf{Visual Studio Code} e, per
chi usa vim o neovim, il plugin \textbf{vim-go}.

Dietro le quinte tutti questi strumenti si appoggiano a \textbf{gopls}, il
\emph{language server} ufficiale di Go. Un language server implementa il
\emph{Language Server Protocol} (LSP), uno standard con cui un editor
ottiene da un processo separato le funzioni avanzate di supporto al
linguaggio: completamento del codice, documentazione al passaggio del mouse,
ricerca dei simboli, diagnostica ed errori. Il vantaggio architetturale è che
l'intelligenza sul linguaggio è scritta una sola volta (in \cod{gopls}) e
riutilizzata da qualunque editor che parli LSP.

% ===========================================================================
\chapter{Struttura del programma, package e toolchain}\label{ch:struttura}
% ===========================================================================

\section{Il primo programma}

\subsection{Esempio pratico}

\begin{lstlisting}[style=gocode,caption={\texttt{01\_basics/helloworld.go}}]
// helloworld.go - Primo programma Go (slide 8-9).
//
// Ogni file inizia con la dichiarazione del package. Il package "main" è
// speciale: dice al toolchain di produrre un eseguibile, il cui punto di
// ingresso è la funzione main().
package main

// "fmt" fornisce l'I/O formattato (Println, Printf, ...).
import "fmt"

func main() {
	// Le stringhe Go sono UTF-8 nativamente: nessuna configurazione
	// particolare serve per stampare caratteri non ASCII.
	fmt.Println("Hello, 世界")
}
\end{lstlisting}
\begin{lstlisting}[style=output]
Hello, 世界
\end{lstlisting}

\subsection{Analisi del codice}

\begin{itemize}
  \item \textbf{Nessun punto e virgola} a fine istruzione o dichiarazione. In
        realtà il compilatore li inserisce automaticamente a fine riga in certi
        casi (si veda l'errore comune più avanti); il programmatore non li scrive.
  \item \textbf{Supporto Unicode nativo}: il sorgente è UTF-8 e le stringhe
        contengono testo Unicode, come mostra la stringa \cod{"Hello, 世界"}.
  \item \textbf{Struttura a package}: ogni programma Go è composto da
        \emph{package} (concetto analogo alle librerie C o ai moduli Python).
        Un package è costituito da uno o più file \cod{.go} nella stessa directory.
  \item \textbf{Struttura di un file sorgente}: inizia con la
        \emph{dichiarazione del package} (a quale package appartiene il file),
        seguita dagli \cod{import} degli altri package.
  \item \textbf{Punto di ingresso}: l'esecuzione parte dal package \cod{main} e
        dalla sua funzione \cod{main()}.
  \item Il package \cod{fmt} fornisce le funzioni per l'input e l'output
        formattato (\cod{Println}, \cod{Printf}, \ldots).
\end{itemize}

\begin{attenzione}[Errore comune: la graffa \cod{\{} a capo]
Poiché i punti e virgola sono inseriti dal compilatore a fine riga, la graffa
aperta deve stare sulla stessa riga della parola chiave. In particolare in un
\cod{if}\,/\,\cod{else} l'\cod{else} deve comparire subito dopo la \cod{\}} che
chiude il blocco precedente. Il codice
\begin{lstlisting}[style=gocode,numbers=none,frame=none,backgroundcolor={},aboveskip=0.3em,belowskip=0.3em]
if x > 0 {
	fmt.Println(1)
}
else {
	fmt.Println(2)
}
\end{lstlisting}
non compila: dopo la \cod{\}} della riga~3 il compilatore inserisce un
punto e virgola, l'istruzione \cod{if} risulta conclusa, e la riga successiva inizia
con un \cod{else} inatteso. Il messaggio reale è:
\texttt{syntax error: unexpected keyword else, expected \}}.
\end{attenzione}

\section{Package e visibilità}

Un programma Go è composto da package, e l'esecuzione parte dal package
\cod{main}. Un package può contenere dichiarazioni di tipi, funzioni, variabili
e costanti, e può essere molto piccolo o molto grande.

\begin{definizione}{regola di visibilità}
In Go la \textbf{visibilità} di un nome fuori dal package che lo dichiara è
determinata dalla \emph{maiuscola iniziale}. Un nome che inizia con una lettera
\textbf{maiuscola} è \emph{esportato} (pubblico); uno che inizia con una lettera
\textbf{minuscola} non è esportato (privato al package).
\end{definizione}

Così \cod{Foo} è esportato e \cod{foo} no. Un esempio concreto: \cod{math.Pi} è
visibile da altri package, mentre \cod{math.pi} non lo è, e usarlo produce un
errore di compilazione. Le slide riportano il messaggio
\texttt{./prog.go:9:19: undefined: math.pi}; il compilatore usato per
verificare questa guida (Go 1.27.3) lo completa con un suggerimento:

\begin{lstlisting}[style=output]
undefined: math.pi (but have Pi)
\end{lstlisting}

\begin{confronto}{C, Python, Java}
In \textbf{Java} la visibilità si controlla con le parole chiave
\cod{public}/\cod{private}/\cod{protected}; in \textbf{C} con \cod{static} (a
livello di file) e con ciò che si dichiara negli header; in \textbf{Python} vale
solo la convenzione del prefisso \cod{\_}. In \textbf{Go} la regola è
\emph{imposta dal compilatore} ed è codificata nel nome stesso: guardando un
identificatore si sa subito se fa parte dell'interfaccia pubblica di un package.
\end{confronto}

\section{Importare package}\label{sec:import}

L'istruzione \cod{import} include altri package. Gli import si possono
raggruppare in una forma «fattorizzata» tra parentesi, più leggibile della
ripetizione di tante istruzioni \cod{import}.

\begin{lstlisting}[style=gocode,caption={\texttt{01\_basics/imports.go}}]
// imports.go - Import "fattorizzato" tra parentesi (slide 12) e visibilità
// dei nomi (slide 11).
package main

import (
	"fmt"
	"math"
)

func main() {
	// math.Sqrt restituisce un float64; %g sceglie la forma compatta
	// più breve tra %e (esponenziale) e %f (decimale).
	fmt.Printf("Now you have %g problems.\n", math.Sqrt(7))

	// Un nome esportato (iniziale maiuscola) è accessibile dall'esterno:
	// math.Pi è visibile, mentre math.pi non esiste e l'uso di math.pi
	// produce l'errore di compilazione "undefined: math.pi".
	fmt.Println("math.Pi =", math.Pi)
}
\end{lstlisting}
\begin{lstlisting}[style=output]
Now you have 2.6457513110645907 problems.
math.Pi = 3.141592653589793
\end{lstlisting}

Il formato \cod{\%g} stampa un numero in virgola mobile nella forma compatta,
scegliendo automaticamente la più corta tra \cod{\%e} (esponenziale) e \cod{\%f}
(decimale). Si noti inoltre che l'uso di \cod{math.Sqrt(7)} richiede il prefisso
del package (\cod{math.}): i nomi importati non vengono mai «riversati» nello
spazio dei nomi locale.

\section{La toolchain di Go}\label{sec:toolchain}

Go è un linguaggio \textbf{compilato}. Il comando \cod{go} è la \emph{toolchain}
con cui si scaricano, compilano e installano package e comandi. Un punto
distintivo è la \textbf{configurazione zero}: non servono script di build esterni
né makefile; il compilatore ricava le dipendenze dagli stessi file sorgente.

\begin{description}
  \item[\texttt{go run}] compila ed esegue direttamente il programma.
  \item[\texttt{go build}] produce un eseguibile binario.
  \item[\texttt{go install}] compila e installa package o comandi.
\end{description}

\begin{lstlisting}[style=shell]
$ go run helloworld.go
Hello, 世界
$ go build helloworld.go
$ ls helloworld*
helloworld      helloworld.go
$ ./helloworld
Hello, 世界
\end{lstlisting}

Su Windows l'eseguibile prodotto da \cod{go build} si chiama
\texttt{helloworld.exe}. Il comando \cod{go run} compila in una directory
temporanea ed esegue il risultato: è comodo in fase di sviluppo, ma il binario non
viene conservato.

\begin{confronto}{Java, Python, C}
Rispetto a \textbf{Java}, il binario di Go è già codice macchina nativo: non
richiede una JVM per essere eseguito. Rispetto a \textbf{C}, non servono
\cod{Makefile} o \cod{CMake}: le dipendenze sono dichiarate dagli \cod{import}.
Rispetto a \textbf{Python}, la compilazione è un passo esplicito (\cod{go build})
ma \cod{go run} lo nasconde offrendo un ciclo modifica--esecuzione rapido.
\end{confronto}

\section{Moduli}\label{sec:moduli}

\begin{definizione}{modulo}
Un \textbf{modulo} è un insieme di package Go correlati, memorizzati in una
struttura di directory con un file \texttt{go.mod} nella radice.
\end{definizione}

Il file \texttt{go.mod} definisce:

\begin{itemize}
  \item il \textbf{percorso del modulo} (\emph{module path}), usato anche come
        percorso di import per la directory radice;
  \item la \textbf{versione minima di Go} richiesta dal modulo;
  \item le \textbf{dipendenze}: gli altri moduli necessari per una compilazione
        corretta, con la relativa versione minima.
\end{itemize}

Un esempio, come nelle slide:

\begin{lstlisting}[style=shell,caption={Esempio di \texttt{go.mod} (dalla slide)}]
module example.com/mymodule

go 1.17

require github.com/gorilla/mux v1.8.0
\end{lstlisting}

I comandi principali sono:

\begin{lstlisting}[style=shell]
$ go mod init <nome_modulo>   # crea il file go.mod
$ go mod tidy                 # aggiunge i requisiti mancanti e rimuove quelli inutilizzati
\end{lstlisting}

Nella cartella \texttt{examples/} di questa guida il file \texttt{go.mod} è
minimo, perché gli esempi usano solo la libreria standard:

\begin{lstlisting}[style=shell,caption={\texttt{examples/go.mod}}]
module goguide

go 1.22
\end{lstlisting}

\section{Regole del compilatore e strumenti di qualità}\label{sec:strumenti}

Go è «pignolo»: alcune pratiche che in altri linguaggi producono al massimo un
\emph{warning} qui sono \textbf{errori di compilazione}. Il motivo è tenere il
codice pulito per costruzione e rendere la lettura di sorgenti altrui più
prevedibile.

\begin{itemize}
  \item Una \textbf{variabile dichiarata e non usata} è un errore.
  \item Un \textbf{import non usato} è un errore.
  \item In un \cod{if}/\cod{else} la graffa \cod{\{} deve stare alla fine della
        stessa riga (\cod{\} else \{}, \cod{\} else if \ldots\ \{}); si veda
        l'errore comune sopra.
\end{itemize}

I messaggi reali del compilatore sono:

\begin{lstlisting}[style=output]
declared and not used: x
"os" imported and not used
\end{lstlisting}

Per le variabili di cui non interessa il valore (per esempio l'indice di un ciclo
quando serve solo l'elemento) si usa il \textbf{blank identifier} \cod{\_}, che
scarta il valore e soddisfa il compilatore (Sezione~\ref{sec:range}).

Gli strumenti da riga di comando consigliati dalle slide sono:

\begin{description}
  \item[\texttt{gofmt}] formatta il codice Go secondo lo stile standard:
        \cod{gofmt -w yourcode.go} riscrive il file formattato. Non esistono
        «guerre di stile»: il formato è uno solo.
  \item[\texttt{goimports}] aggiunge e rimuove automaticamente gli import
        (utile proprio per evitare l'errore di import inutilizzato).
  \item[\texttt{godoc}] permette di consultare la documentazione dei package.
\end{description}

\begin{nota}[Perché \texttt{gofmt} è importante]
Tutti gli esempi di questa guida sono stati formattati con \cod{gofmt}. In Go
l'indentazione con tabulazioni, la posizione delle graffe e lo spazio attorno
agli operatori non sono questioni di gusto personale: è la forma canonica che
gli editor applicano al salvataggio.
\end{nota}

% ===========================================================================
\chapter{Variabili e tipi}\label{ch:tipi}
% ===========================================================================

\section{Dichiarare variabili}

\subsection{Introduzione}

Go offre due modi per dichiarare una variabile: l'istruzione \cod{var} e la
\emph{dichiarazione breve} \cod{:=}. La prima si può usare sia a livello di
package sia dentro le funzioni; la seconda solo dentro le funzioni.

\subsection{Sintassi}

\begin{itemize}
  \item \cod{var} dichiara una o più variabili. Il \textbf{tipo viene per ultimo},
        dopo il nome: \cod{var i int}.
  \item È ammesso un \textbf{inizializzatore} per ciascuna variabile; se è
        presente, il tipo si può omettere e viene dedotto: \cod{var a, b = 1, "due"}.
  \item Una variabile dichiarata senza valore iniziale riceve il suo
        \textbf{valore zero}.
  \item La \textbf{dichiarazione breve} \cod{count := 10} dichiara e inizializza
        deducendo il tipo; è consentita solo all'interno delle funzioni.
\end{itemize}

\begin{definizione}{valore zero}
Ogni tipo di Go ha un \textbf{valore zero}, assegnato automaticamente a una
variabile non inizializzata esplicitamente: \cod{0} per i numeri, \cod{false}
per i booleani, la stringa vuota \cod{""} per le stringhe e \cod{nil} per
puntatori e altri tipi riferimento.
\end{definizione}

\subsection{Esempio pratico}

\begin{lstlisting}[style=gocode,caption={\texttt{01\_basics/variables.go}}]
// variables.go - Dichiarazione di variabili e valori zero (slide 17).
package main

import "fmt"

// Variabili a livello di package: qui ":=" NON è ammesso, serve "var".
// Il tipo segue il nome; tutte e tre ricevono il valore zero (false).
var c, python, java bool

func main() {
	// Variabile locale senza inizializzatore: valore zero (0 per gli int).
	var i int
	fmt.Println(i, c, python, java)

	// Con inizializzatori (uno per variabile) il tipo può essere omesso
	// e viene dedotto dal valore a destra.
	var a, b = 1, "due"
	fmt.Println(a, b)

	// Dichiarazione breve ":=" - solo dentro le funzioni.
	count := 10
	fmt.Println(count)

	// Valori zero dei tipi base: nessuna variabile Go è "non inizializzata".
	var (
		n int
		f float64
		s string
		t bool
		p *int
	)
	fmt.Printf("int=%d float64=%g string=%q bool=%t pointer=%v\n", n, f, s, t, p)
}
\end{lstlisting}
\begin{lstlisting}[style=output]
0 false false false
1 due
10
int=0 float64=0 string="" bool=false pointer=<nil>
\end{lstlisting}

\subsection{Analisi del codice}

Le variabili \cod{c}, \cod{python} e \cod{java} sono dichiarate a livello di
package con \cod{var} (qui \cod{:=} non sarebbe permesso). Un'unica
dichiarazione, \cod{var c, python, java bool}, assegna il tipo \cod{bool} a tutte
e tre; poiché nessuna ha un inizializzatore, valgono \cod{false}. Dentro
\cod{main}, \cod{var i int} vale \cod{0}. L'ultima parte stampa i valori zero di
tipi diversi (puntatore \cod{nil} compreso), dimostrando che \emph{in Go non
esistono variabili non inizializzate}.

\begin{confronto}{C, Python, Java}
In \textbf{C} una variabile locale non inizializzata ha un valore
\emph{indeterminato}: leggerla è un comportamento non definito e una classica
fonte di bug. In \textbf{Java} i campi hanno un valore di default ma le variabili
locali devono essere assegnate prima dell'uso (verificato dal compilatore). In
\textbf{Python} non si dichiara nulla: il nome nasce con la prima assegnazione.
In \textbf{Go} il valore zero è definito per tutti i tipi, quindi una variabile è
sempre in uno stato valido. Per l'inferenza, \cod{:=} ricorda \cod{var} di Java
(10+) per le variabili locali, ma in Go il tipo resta statico e fisso.
\end{confronto}

\begin{attenzione}[Errori comuni con \cod{:=} e \cod{var}]
\begin{enumerate}
  \item \textbf{Usare \cod{:=} fuori da una funzione.} A livello di package ogni
        istruzione deve iniziare con una parola chiave di dichiarazione
        (\cod{var}, \cod{func}, \ldots). Scrivere \cod{x := 3} fuori da una funzione
        dà \texttt{syntax error: non-declaration statement outside function body}.
  \item \textbf{Ridichiarare con \cod{:=}.} \cod{a := 1} seguito da \cod{a := 2}
        nello stesso blocco dà \texttt{no new variables on left side of :=}:
        \cod{:=} richiede almeno una variabile \emph{nuova}; per riassegnare si
        usa \cod{=}.
  \item \textbf{Variabili inutilizzate:} sono errori, non warning
        (Sezione~\ref{sec:strumenti}).
\end{enumerate}
\end{attenzione}

\section{Tipi di base, conversioni e inferenza}

I tipi di base comprendono \cod{bool}, \cod{string}, \cod{int}, \cod{uint},
\cod{float32}, \cod{float64} e altri.

\begin{definizione}{conversione esplicita}
Go richiede una \textbf{conversione esplicita} tra tipi diversi, anche
numerici: \cod{var f float64 = float64(i)}. Non esiste alcuna promozione
implicita (a differenza di C).
\end{definizione}

L'\textbf{inferenza di tipo} permette al compilatore di dedurre il tipo di una
variabile dal valore a destra: se \cod{var i int} è già dichiarata, allora in
\cod{j := i} la variabile \cod{j} è di tipo \cod{int}.

\begin{lstlisting}[style=gocode,caption={\texttt{01\_basics/types.go}}]
// types.go - Tipi base, conversioni esplicite e inferenza (slide 18).
package main

import "fmt"

func main() {
	// Tipi base.
	var (
		b  bool    = true
		s  string  = "go"
		i  int     = 42
		u  uint    = 7
		f2 float32 = 1.5
		f6 float64 = 2.25
	)
	fmt.Println(b, s, i, u, f2, f6)

	// Conversione esplicita: obbligatoria tra tipi diversi (a differenza
	// di C, dove int -> double avviene implicitamente).
	var n int = 42
	var f float64 = float64(n)
	fmt.Printf("%d -> %g\n", n, f)

	// Inferenza: j ha lo stesso tipo di k (int).
	var k int
	j := k
	fmt.Printf("tipo di j: %T\n", j)

	// Il tipo dedotto da una costante letterale dipende dalla sua forma.
	x := 3.14 // float64
	y := 'a'  // rune (alias di int32)
	fmt.Printf("%T %T\n", x, y)

	// Non si può sommare int e float64 senza conversione:
	//   fmt.Println(n + f)   // errore di compilazione: mismatched types
	fmt.Println(float64(n) + f)
}
\end{lstlisting}
\begin{lstlisting}[style=output]
true go 42 7 1.5 2.25
42 -> 42
tipo di j: int
float64 int32
84
\end{lstlisting}

\subsection{Analisi del codice}

\begin{itemize}
  \item \cod{float64(n)} converte il valore: omettere la conversione e scrivere
        \cod{n + f} con \cod{n} di tipo \cod{int} e \cod{f} di tipo \cod{float64}
        produce l'errore di compilazione
        \texttt{invalid operation: n + f (mismatched types int and float64)}.
  \item \cod{\%T} stampa il tipo dedotto: \cod{int} per \cod{j}, \cod{float64} per la
        costante \cod{3.14} e \cod{int32} per \cod{'a'}.
  \item In Go il letterale \cod{'a'} è una \emph{rune}, un alias di
        \cod{int32} che rappresenta un punto di codice Unicode. Per questo
        \cod{\%T} stampa \cod{int32}.
\end{itemize}

\begin{confronto}{C, Python, Java}
In \textbf{C} somme tra \cod{int} e \cod{double} funzionano grazie alle
conversioni aritmetiche implicite; in \textbf{Java} sono implicite le conversioni
\emph{di allargamento} (\cod{int}$\to$\cod{double}). In \textbf{Go} nessuna: la
scelta evita perdite di precisione o di segno non volute, al prezzo di un
\cod{float64(\ldots)} in più nel codice. In \textbf{Python} i tipi sono dinamici e
\cod{1 + 2.5} funziona perché è l'interprete a decidere a runtime.
\end{confronto}

% ===========================================================================
\chapter{Funzioni}\label{ch:funzioni}
% ===========================================================================

\section{Dichiarazione e parametri}

Una funzione può avere zero o più argomenti, e anche in questo caso il
\textbf{tipo segue il nome}: \cod{func add(x int, y int) int}. Quando più
parametri consecutivi hanno lo stesso tipo si può abbreviare:
\cod{func add(x, y int) int}.

\begin{lstlisting}[style=gocode,caption={\texttt{02\_functions/functions.go}}]
// functions.go - Funzioni: parametri, tipo dopo il nome, forma abbreviata
// (slide 13).
package main

import "fmt"

// Il tipo segue il nome del parametro; anche il tipo di ritorno è in coda.
func add(x int, y int) int {
	return x + y
}

// Forma abbreviata: parametri consecutivi dello stesso tipo.
func addShort(x, y int) int {
	return x + y
}

func main() {
	fmt.Println(add(3, 4))
	fmt.Println(addShort(10, 20))
}
\end{lstlisting}
\begin{lstlisting}[style=output]
7
30
\end{lstlisting}

Il tipo di ritorno compare dopo la lista dei parametri. Go è
\emph{call-by-value}: \cod{x} e \cod{y} sono copie degli argomenti (si veda la
Sezione~\ref{sec:puntatori} per come modificare l'originale con un puntatore).

\begin{confronto}{C, Python, Java}
Chi viene da C o Java è abituato a \cod{int add(int x, int y)}, con il tipo
\emph{prima} del nome. Go lo inverte (\cod{x int}) per due ragioni leggibili:
le dichiarazioni si leggono da sinistra a destra («\cod{x} è un \cod{int}») e
le forme complesse (puntatori a funzioni, slice di mappe, \ldots) restano
leggibili. In Python le annotazioni di tipo seguono il nome (\cod{x: int}) ma sono
facoltative e non verificate a runtime.
\end{confronto}

\section{Risultati multipli}

Una funzione può restituire \textbf{un numero qualsiasi di risultati}.

\begin{lstlisting}[style=gocode,caption={\texttt{02\_functions/swap.go}}]
// swap.go - Una funzione può restituire più risultati (slide 13-14).
package main

import "fmt"

func swap(x, y string) (string, string) {
	return y, x
}

func main() {
	// Assegnazione multipla: entrambi i risultati vengono "spacchettati".
	a, b := swap("hello", "world")
	fmt.Println(a, b)
}
\end{lstlisting}
\begin{lstlisting}[style=output]
world hello
\end{lstlisting}

L'assegnazione multipla \cod{a, b := swap("hello", "world")} spacchetta i due
valori restituiti. Il caso d'uso più importante è restituire un
\textbf{risultato e un errore}: l'errore è per convenzione l'\emph{ultimo} valore
restituito e \cod{nil} significa «nessun errore». La Sezione~\ref{sec:errori}
approfondisce questo schema; qui lo si vede con \cod{os.ReadFile}.

\begin{lstlisting}[style=gocode,caption={\texttt{02\_functions/result\_and\_error.go}}]
// result_and_error.go - Convenzione risultato + errore (slide 13, 15).
//
// La slide usa os.ReadFile("example.txt"). Qui il file viene prima creato in
// una directory temporanea, così l'esempio è autonomo e non dipende da file
// preesistenti; poi si prova anche un file inesistente per mostrare il ramo
// di errore. Il messaggio dell'errore dipende dal sistema operativo, quindi
// qui viene stampato solo il fatto che l'errore sia avvenuto.
package main

import (
	"fmt"
	"os"
	"path/filepath"
)

func readAndPrint(name string) {
	data, err := os.ReadFile(name)
	// Per convenzione l'errore è l'ultimo valore restituito e nil
	// significa "nessun errore". Va controllato subito.
	if err != nil {
		fmt.Println("Error: lettura fallita, err != nil")
		return
	}
	fmt.Println("File contents:", string(data))
}

func main() {
	dir, err := os.MkdirTemp("", "goguide")
	if err != nil {
		fmt.Println("Error:", err)
		return
	}
	defer os.RemoveAll(dir) // pulizia al termine di main

	path := filepath.Join(dir, "example.txt")
	if err := os.WriteFile(path, []byte("ciao dal file"), 0o644); err != nil {
		fmt.Println("Error:", err)
		return
	}

	readAndPrint(path)                         // caso di successo
	readAndPrint(filepath.Join(dir, "no.txt")) // caso di errore
}
\end{lstlisting}
\begin{lstlisting}[style=output]
File contents: ciao dal file
Error: lettura fallita, err != nil
\end{lstlisting}

\begin{nota}[Differenza rispetto alla slide]
La slide legge un file \texttt{example.txt} che si assume già esistente. Per
rendere l'esempio autonomo e riproducibile, il file viene prima creato in una
directory temporanea (con \cod{defer os.RemoveAll} per ripulirla), e
viene provato anche un file inesistente per mostrare il ramo d'errore. La parte
essenziale (\cod{data, err := os.ReadFile(...)} seguito da \cod{if err != nil})
è quella della slide. Il testo esatto dell'errore dipende dal sistema operativo,
quindi viene stampato un messaggio fisso.
\end{nota}

\begin{confronto}{C, Python, Java}
In \textbf{C} per restituire più valori si usano parametri di output
(puntatori) o struct; in \textbf{Python} si restituisce una tupla (meccanismo
simile a quello di Go); in \textbf{Java} servirebbe un oggetto contenitore.
Go ha il supporto nativo ai risultati multipli, ed è la base della gestione
degli errori (niente eccezioni).
\end{confronto}

\section{Valori di ritorno con nome}

Go permette di dare un nome ai valori di ritorno nella firma. Le variabili con
nome sono inizializzate al loro \textbf{valore zero} e un \cod{return} senza
argomenti (\emph{«naked return»}) restituisce i loro valori correnti.

\begin{lstlisting}[style=gocode,caption={\texttt{02\_functions/named\_returns.go}}]
// named_returns.go - Valori di ritorno con nome e "naked return" (slide 16).
package main

import "fmt"

// x e y sono variabili locali dichiarate nella firma, inizializzate al
// valore zero. Il "return" nudo restituisce i loro valori correnti.
func split(sum int) (x, y int) {
	x = sum * 4 / 9
	y = sum - x
	return // equivale a: return x, y
}

func main() {
	fmt.Println(split(17))
}
\end{lstlisting}
\begin{lstlisting}[style=output]
7 10
\end{lstlisting}

\subsection{Analisi del codice}

\cod{x} e \cod{y} sono variabili locali dichiarate nella firma
\cod{(x, y int)}. Con \cod{sum = 17}: \cod{x = 17 * 4 / 9 = 7} (divisione
intera) e \cod{y = 17 - 7 = 10}. Il \cod{return} nudo equivale a
\cod{return x, y}. L'output \cod{7 10} è quello stampato da
\cod{fmt.Println(split(17))}, in cui i due risultati diventano due argomenti.

\begin{attenzione}[Osservazione: naked return e leggibilità]
Il \emph{naked return} è comodo in funzioni brevi, ma in funzioni lunghe obbliga
il lettore a risalire alla firma per sapere che cosa viene restituito. I nomi
dei risultati hanno anche valore di documentazione: \cod{(x, y int)} dice più di
\cod{(int, int)}.
\end{attenzione}

\section{Funzioni variadiche}

Una \textbf{funzione variadica} accetta un numero variabile di argomenti.
\cod{fmt.Println} ne è un esempio. Si dichiara con il suffisso
\cod{...tipo} sull'ultimo parametro.

\begin{lstlisting}[style=gocode,caption={\texttt{02\_functions/variadic.go}}]
// variadic.go - Funzioni variadiche (slide 68).
package main

import "fmt"

// Dentro la funzione "nums" è una slice []int.
func sum(nums ...int) {
	fmt.Print(nums, " ")
	total := 0
	for _, num := range nums {
		total += num
	}
	fmt.Println(total)
}

func main() {
	sum(1, 2)
	sum(1, 2, 3)

	// Per passare una slice esistente si usa il suffisso "...".
	nums := []int{1, 2, 3, 4}
	sum(nums...)

	// Zero argomenti: nums è una slice nil e il totale è 0.
	sum()
}
\end{lstlisting}
\begin{lstlisting}[style=output]
[1 2] 3
[1 2 3] 6
[1 2 3 4] 10
[] 0
\end{lstlisting}

\subsection{Analisi del codice}

\begin{itemize}
  \item All'interno della funzione, \cod{nums} è una \textbf{slice} \cod{[]int}
        (Sezione~\ref{sec:slice}): si può iterare con \cod{range} e usare \cod{len}.
  \item Per passare una slice già esistente come lista di argomenti si usa il
        suffisso \cod{...} \emph{alla chiamata}: \cod{sum(nums...)}.
  \item Chiamata senza argomenti: \cod{nums} è una slice \cod{nil}, quindi
        stampa \cod{[]} e il totale è \cod{0}.
\end{itemize}

\begin{confronto}{C, Python, Java}
In \textbf{C} le funzioni variadiche usano \cod{stdarg.h} (\cod{va\_list}) e non
c'è controllo di tipo; in \textbf{Java} la sintassi \cod{int... nums} è
praticamente identica a quella di Go; in \textbf{Python} si usa \cod{*args}, e
l'espansione \cod{sum(*nums)} corrisponde a \cod{sum(nums...)} in Go.
\end{confronto}

\section{Funzioni anonime e closure}

Le funzioni Go possono essere \textbf{anonime}, cioè senza nome: si definiscono
\emph{in linea} e si usano subito o si assegnano a una variabile.

\begin{lstlisting}[style=gocode,caption={\texttt{02\_functions/anonymous.go}}]
// anonymous.go - Funzioni anonime (slide 37).
package main

import "fmt"

func main() {
	// Funzione anonima assegnata a una variabile.
	f := func(x, y int) int {
		return x + y
	}
	fmt.Println(f(3, 4)) // 7

	// Funzione anonima definita e invocata subito: le parentesi finali
	// () eseguono la chiamata.
	func(msg string) {
		fmt.Println("invocata subito:", msg)
	}("ciao")
}
\end{lstlisting}
\begin{lstlisting}[style=output]
7
invocata subito: ciao
\end{lstlisting}

\begin{definizione}{closure}
Una \textbf{closure} è una funzione anonima annidata che \emph{cattura le
variabili} dell'ambiente che la racchiude. Mantiene un proprio \textbf{stato
privato}, indipendente da quello di altre istanze.
\end{definizione}

\begin{lstlisting}[style=gocode,caption={\texttt{02\_functions/closure.go}}]
// closure.go - Closure con stato privato (slide 38).
package main

import "fmt"

// intSeq restituisce una funzione. La funzione anonima "cattura" la
// variabile i dell'ambiente che la racchiude: i sopravvive alle
// chiamate di intSeq ed è privata di quella particolare closure.
func intSeq() func() int {
	i := 0
	return func() int {
		i++
		return i
	}
}

func main() {
	nextInt := intSeq()
	fmt.Println(nextInt()) // 1
	fmt.Println(nextInt()) // 2
	fmt.Println(nextInt()) // 3

	// Una nuova closure ha una propria copia indipendente di i.
	newInts := intSeq()
	fmt.Println(newInts()) // 1
}
\end{lstlisting}
\begin{lstlisting}[style=output]
1
2
3
1
\end{lstlisting}

\subsection{Analisi del codice}

\cod{intSeq()} dichiara \cod{i := 0} e restituisce una funzione anonima che
incrementa e restituisce \cod{i}. Normalmente una variabile locale scompare quando
la funzione ritorna; qui \cod{i} \emph{sopravvive} perché la closure lo
referenzia (il runtime la sposta sullo heap, come per i puntatori a variabili
locali, Sezione~\ref{sec:puntatori}). Ogni chiamata di \cod{intSeq()} crea una
\emph{nuova} variabile \cod{i}: per questo \cod{newInts()} ricomincia da
\cod{1} senza influenzare \cod{nextInt}.

Un secondo esempio, elegante ed efficiente, genera i numeri di Fibonacci
tenendo lo stato \cod{(x, y)} nella closure:

\begin{lstlisting}[style=gocode,caption={\texttt{02\_functions/closure\_fib.go}}]
// closure_fib.go - Fibonacci con una closure (slide 39).
package main

import "fmt"

// fibonacci restituisce una funzione che, a ogni chiamata, restituisce
// il numero di Fibonacci successivo. Lo stato (x, y) vive nella closure.
func fibonacci() func() int {
	x, y := 0, 1
	return func() int {
		// Assegnazione parallela: il lato destro è valutato per intero
		// prima di assegnare, quindi non serve una variabile temporanea.
		x, y = y, x+y
		return x
	}
}

func main() {
	f := fibonacci()
	for i := 0; i < 10; i++ {
		fmt.Println(f())
	}
}
\end{lstlisting}
\begin{lstlisting}[style=output]
1
1
2
3
5
8
13
21
34
55
\end{lstlisting}

L'assegnazione parallela \cod{x, y = y, x+y} valuta \emph{prima} tutto il lato
destro e poi assegna, quindi non serve una variabile temporanea. A ogni
chiamata, \cod{x} avanza di un passo nella successione (1, 1, 2, 3, \ldots).

\begin{confronto}{C, Python, Java}
\textbf{C} non ha closure: si ottiene un effetto simile con un puntatore a
funzione più un puntatore \cod{void*} al «contesto», gestito a mano. In
\textbf{Python} le closure esistono (per \emph{modificare} una variabile
catturata serve \cod{nonlocal}). In \textbf{Java} le lambda catturano solo
variabili \emph{effectively final}: non possono modificarle, quindi serve di
solito un oggetto contenitore (es.\ un array di un elemento o un
\cod{AtomicInteger}). In \textbf{Go} la closure cattura la variabile stessa, non una
copia del suo valore, e può modificarla liberamente.
\end{confronto}

\subsection{Closure come middleware}\label{sec:funzioni-middleware}

Un impiego pratico delle closure è il \emph{middleware}: una closure può
avvolgere una funzione \emph{handler} per eseguire logica prima e/o dopo di essa.
La slide mostra \cod{timed}, che avvolge un handler HTTP per misurare il tempo
di elaborazione di una richiesta; con un server vero si registrerebbe
con \cod{http.HandleFunc("/hello", timed(hello))}.

\begin{lstlisting}[style=gocode,caption={\texttt{02\_functions/middleware.go}}]
// middleware.go - Closure come middleware (slide 40).
//
// La slide mostra timed(), che avvolge un handler HTTP per misurarne il tempo
// di elaborazione. Per rendere l'esempio eseguibile senza aprire porte di
// rete, la richiesta viene simulata con il package httptest.
package main

import (
	"fmt"
	"net/http"
	"net/http/httptest"
	"time"
)

// timed riceve un handler e ne restituisce uno nuovo che esegue codice
// "prima" e "dopo" la chiamata all'handler originale. f è catturato dalla
// closure.
func timed(f func(http.ResponseWriter, *http.Request)) func(http.ResponseWriter, *http.Request) {
	return func(w http.ResponseWriter, r *http.Request) {
		start := time.Now()
		f(w, r)
		// Il tempo esatto cambia a ogni esecuzione: lo si arrotonda per
		// rendere l'output (quasi sempre) stabile.
		fmt.Println("The request took", time.Since(start).Round(50*time.Millisecond))
	}
}

func hello(w http.ResponseWriter, r *http.Request) {
	time.Sleep(100 * time.Millisecond) // lavoro simulato
	fmt.Fprintln(w, "hello")
}

func main() {
	// Con un server vero si scriverebbe: http.HandleFunc("/hello", timed(hello))
	handler := timed(hello)

	req := httptest.NewRequest("GET", "/hello", nil)
	rec := httptest.NewRecorder()
	handler(rec, req)

	fmt.Print("risposta: ", rec.Body.String())
}
\end{lstlisting}
\begin{lstlisting}[style=output]
The request took 100ms
risposta: hello
\end{lstlisting}

\begin{nota}[Come è stato reso eseguibile]
La slide mostra solo \cod{timed}. Per verificarlo senza aprire una porta di rete,
l'esempio simula la richiesta con il package \cod{httptest} della libreria standard.
Il tempo viene arrotondato a multipli di 50\,ms per rendere l'output quasi
sempre stabile; il valore esatto, essendo una misura di tempo, non è
deterministico: l'handler dorme 100\,ms, quindi si legge circa \cod{100ms}.
\end{nota}

\subsection{Funzioni anonime e goroutine: una nota sulle variabili di ciclo}

Le closure sono usate molto spesso nel codice concorrente (Capitolo~\ref{ch:conc}),
dove una funzione anonima è lanciata come goroutine all'interno di un ciclo. Il
punto delicato è se la closure veda una \emph{variabile di ciclo condivisa} da
tutte le iterazioni oppure una nuova per ciascuna. Dal Go~1.22, ogni iterazione
di un ciclo \cod{for i := \ldots} ha la propria copia di \cod{i}, quindi
\cod{go func() \{ use(i) \}()} funziona come ci si aspetta; con versioni
precedenti (o con \texttt{go.mod} che dichiara una versione anteriore) le closure
condividevano un'unica \cod{i}, causando bug sottili. Gli esempi di questa guida
dichiarano \cod{go 1.22} in \texttt{go.mod}.

% ===========================================================================
\chapter{Controllo di flusso}\label{ch:flusso}
% ===========================================================================

Le istruzioni di controllo del flusso di Go sono \cod{for}, \cod{if} (con
\cod{else}), \cod{switch} e \cod{defer}. Hanno in comune due regole di stile
che le distinguono da C e Java: \textbf{nessuna parentesi} attorno alle
condizioni e \textbf{graffe sempre obbligatorie}, anche per un solo statement.

\section{Il ciclo \texttt{for}}\label{sec:for}

Go ha \textbf{un solo costrutto di ciclo}: \cod{for}. Nella sua forma completa ha
tre componenti: l'istruzione \emph{init}, l'espressione \emph{condizione} e
l'istruzione \emph{post}.

\begin{lstlisting}[style=gocode,numbers=none,caption={Forma generale}]
for init; condizione; post {
	// corpo
}
\end{lstlisting}

Le istruzioni \emph{init} e \emph{post} sono facoltative: se omesse, il
\cod{for} diventa il «while» di Go. Se si omette anche la condizione, il ciclo è
\textbf{infinito} (\cod{for \{ \}}).

\begin{lstlisting}[style=gocode,caption={\texttt{03\_control\_flow/for\_loops.go}}]
// for_loops.go - L'unico costrutto di ciclo di Go: for (slide 20-21).
package main

import "fmt"

func main() {
	// Forma completa: init; condizione; post. Niente parentesi attorno
	// ai tre componenti, graffe sempre obbligatorie.
	sum := 0
	for i := 0; i < 10; i++ {
		sum += i
	}
	fmt.Println("somma 0..9 =", sum)

	// Senza init e post il for si comporta come un "while".
	sum = 1
	for sum < 1000 {
		sum += sum
	}
	fmt.Println("primo valore >= 1000:", sum)

	// Senza condizione il ciclo è infinito; si esce con break (o return).
	n := 0
	for {
		n++
		if n == 5 {
			break
		}
	}
	fmt.Println("uscito dal ciclo infinito con n =", n)
}
\end{lstlisting}
\begin{lstlisting}[style=output]
somma 0..9 = 45
primo valore >= 1000: 1024
uscito dal ciclo infinito con n = 5
\end{lstlisting}

\subsection{Analisi del codice}

\begin{itemize}
  \item Il primo ciclo somma $0+1+\dots+9=45$.
  \item Il secondo è un \emph{while}: \cod{sum} raddoppia a ogni iterazione
        (1, 2, 4, \ldots) finché non raggiunge o supera 1000, cioè
        $2^{10}=1024$.
  \item Il terzo è infinito e viene interrotto con \cod{break} quando
        \cod{n} vale 5.
\end{itemize}

\begin{confronto}{C, Python, Java}
In \textbf{C} e \textbf{Java} esistono \cod{for}, \cod{while} e \cod{do-while};
in \textbf{Python} \cod{for \ldots\ in} e \cod{while}. Go li unifica in un solo
\cod{for}. L'iterazione su una collezione si ottiene con la clausola
\cod{range} (Sezione~\ref{sec:range}), che ricorda il \cod{for-in} di Python.
\end{confronto}

\subsection{Esempio: \texttt{echo}}

Il programma \cod{echo} stampa gli argomenti della riga di comando.
\cod{os.Args} è una \emph{slice} di stringhe (Sezione~\ref{sec:slice}); l'elemento
di indice 0 è il nome del programma, quindi il ciclo parte da 1.

\begin{lstlisting}[style=gocode,caption={\texttt{03\_control\_flow/echo.go}}]
// Echo prints its command-line arguments. (slide 22)
//
// Esecuzione: go run echo.go Hello Go World
package main

import (
	"fmt"
	"os"
)

func main() {
	// s e sep sono inizializzate implicitamente alla stringa vuota.
	var s, sep string
	// os.Args è una slice di stringhe; os.Args[0] è il nome del programma,
	// quindi gli argomenti utili partono da 1.
	for i := 1; i < len(os.Args); i++ {
		// Al primo giro sep è "" (nessuno spazio iniziale); poi diventa " ".
		s += sep + os.Args[i]
		sep = " "
	}
	fmt.Println(s)
}
\end{lstlisting}
\begin{lstlisting}[style=output]
Hello Go World
\end{lstlisting}

Nell'esempio \cod{s} e \cod{sep} sono implicitamente inizializzate alla stringa
vuota (valore zero). Al primo giro \cod{sep} è vuoto e non si inserisce
alcuno spazio iniziale; dal secondo giro \cod{sep} vale \cod{" "}. Il comando
eseguito è \cod{go run echo.go Hello Go World}.

\section{Istruzioni condizionali: \texttt{if}}\label{sec:if}

L'istruzione \cod{if} (con \cod{else}) ha forma simile a quella del \cod{for}:
nessuna parentesi attorno alla condizione, graffe sempre richieste, e si possono
concatenare più condizioni con \cod{if\ldots else if\ldots else}. Come nel
\cod{for}, c'è un'istruzione \emph{opzionale} prima della condizione, separata
da \cod{;}:

\begin{lstlisting}[style=gocode,numbers=none,caption={Forma generale}]
if istruzioneOpzionale1; espressioneBooleana1 {
	blocco1
} else if istruzioneOpzionale2; espressioneBooleana2 {
	blocco2
} else {
	blocco3
}
\end{lstlisting}

\begin{lstlisting}[style=gocode,caption={\texttt{03\_control\_flow/if\_else.go}}]
// if_else.go - if con istruzione di inizializzazione (slide 23-24).
package main

import (
	"fmt"
	"math"
)

// pow restituisce x^n se è inferiore a limit, altrimenti stampa un
// messaggio e restituisce limit.
func pow(x, n, limit float64) float64 {
	// v è dichiarata nell'istruzione iniziale dell'if: è visibile nel blocco
	// if e nel relativo else, ma non dopo.
	if v := math.Pow(x, n); v < limit {
		return v
	} else {
		// Nota: "} else {" deve stare sulla stessa riga.
		fmt.Printf("%g >= %g\n", v, limit)
	}
	// Qui "v" non esiste più (errore di compilazione se la si usasse).
	return limit
}

func main() {
	fmt.Println(pow(3, 2, 10), pow(3, 3, 20))

	// Catena if / else if / else, ciascuno con la propria istruzione
	// opzionale.
	x := 7
	if r := x % 2; r == 0 {
		fmt.Println("pari")
	} else if r == 1 {
		fmt.Println("dispari")
	} else {
		fmt.Println("negativo e dispari")
	}
}
\end{lstlisting}
\begin{lstlisting}[style=output]
27 >= 20
9 20
dispari
\end{lstlisting}

\subsection{Analisi del codice}

\begin{itemize}
  \item L'istruzione opzionale è usata spesso per una breve dichiarazione di
        variabile. In \cod{if v := math.Pow(x, n); v < limit}, la variabile
        \cod{v} esiste \emph{solo} nel blocco \cod{if} e nel relativo
        \cod{else} (\emph{scope}): dopo l'\cod{if} non è più visibile.
  \item \cod{pow(3, 2, 10)} calcola $3^2=9<10$ e restituisce 9. \cod{pow(3, 3,
        20)} calcola $27\ge 20$: stampa \cod{27 >= 20} (messaggio del ramo
        \cod{else}) e restituisce il limite, 20. Poiché gli argomenti di
        \cod{Println} sono valutati prima, la riga \cod{27 >= 20} compare
        \emph{prima} di \cod{9 20}.
  \item L'\cod{else} deve stare sulla stessa riga della graffa di chiusura
        (Sezione~\ref{sec:strumenti}).
  \item L'ultimo blocco mostra una catena \cod{else if} con l'istruzione
        \cod{r := x \% 2}.
\end{itemize}

\section{Istruzioni condizionali: \texttt{switch}}\label{sec:switch}

L'istruzione \cod{switch} sceglie \emph{uno} tra molti casi da eseguire. I
\cod{case} sono valutati dall'alto verso il basso e ci si ferma al primo che
corrisponde.

Le differenze rispetto a C sono sostanziali:

\begin{itemize}
  \item viene eseguito \textbf{solo il caso selezionato}, non tutti quelli
        seguenti: il \cod{break} di C è fornito \emph{automaticamente};
  \item i \cod{case} \textbf{non devono essere costanti};
  \item i valori possono essere \textbf{di qualunque tipo}, non solo interi
        (le slide usano stringhe).
\end{itemize}

\begin{lstlisting}[style=gocode,caption={\texttt{03\_control\_flow/switch.go}}]
// switch.go - switch senza fallthrough implicito (slide 25).
package main

import "fmt"

func describe(day string) {
	switch day {
	case "Monday":
		fmt.Println("Start of the week")
	case "Friday":
		fmt.Println("Almost weekend!")
	case "Saturday", "Sunday": // più valori nello stesso case
		fmt.Println("Weekend!")
	default:
		fmt.Println("Giorno qualunque")
	}
}

func main() {
	for _, d := range []string{"Monday", "Friday", "Sunday", "Wednesday"} {
		describe(d)
	}

	// I case non devono essere costanti: possono essere espressioni.
	// Senza espressione dopo "switch" equivale a "switch true": viene
	// eseguito il primo case la cui condizione è vera.
	n := 15
	switch {
	case n < 10:
		fmt.Println("piccolo")
	case n < 100:
		fmt.Println("medio")
	default:
		fmt.Println("grande")
	}
}
\end{lstlisting}
\begin{lstlisting}[style=output]
Start of the week
Almost weekend!
Weekend!
Giorno qualunque
medio
\end{lstlisting}

\subsection{Analisi del codice}

\cod{describe} confronta \cod{day} con le stringhe dei \cod{case}.
\cod{case "Saturday", "Sunday"} raggruppa più valori nello stesso ramo.
Per \cod{"Wednesday"} nessun caso corrisponde e viene eseguito \cod{default}.
Nell'ultima parte, \cod{switch} \emph{senza espressione} equivale a
\cod{switch true}: i \cod{case} sono condizioni booleane arbitrarie e viene
eseguito il primo che risulta vero (qui \cod{n < 100} perché \cod{n} vale 15).

\begin{confronto}{C, Python, Java}
In \textbf{C} (e nel \cod{switch} classico di \textbf{Java}) senza \cod{break}
l'esecuzione «cade» nel caso successivo: è una fonte tipica di errori. In
\textbf{Go} il comportamento predefinito è l'opposto, quindi un caso non
interferisce con il successivo. In C le etichette devono essere costanti
intere; in Go possono essere espressioni di tipo arbitrario. In
\textbf{Python} (3.10+) il costrutto equivalente è \cod{match}.
\end{confronto}

\section{L'istruzione \texttt{defer}}\label{sec:defer}

\subsection{Introduzione}

\cod{defer} \textbf{rinvia l'esecuzione} di una chiamata di funzione \emph{fino al
momento in cui la funzione che la contiene ritorna}. Il problema che risolve è la
pulizia delle risorse (chiusura di file e connessioni, rilascio di lock): si
scrive l'azione di rilascio \emph{subito dopo} l'acquisizione, e il
linguaggio garantisce che venga eseguita comunque la funzione termini
(anche con un \cod{return} anticipato).

\subsection{Funzionamento}

\begin{itemize}
  \item Gli \textbf{argomenti} della chiamata differita sono valutati
        \emph{immediatamente}, ma la chiamata effettiva avviene solo
        quando la funzione circostante è terminata.
  \item Le chiamate differite sono impilate in uno \textbf{stack} ed eseguite in
        ordine \textbf{LIFO}: l'ultima registrata è la prima eseguita.
\end{itemize}

\begin{lstlisting}[style=gocode,caption={\texttt{03\_control\_flow/defer.go}}]
// defer.go - Esecuzione differita (slide 26).
package main

import "fmt"

func main() {
	// La chiamata viene registrata ora (argomenti valutati ora) ma eseguita
	// solo quando main sta per terminare.
	defer fmt.Println("world")
	fmt.Println("hello")
}
\end{lstlisting}
\begin{lstlisting}[style=output]
hello
world
\end{lstlisting}

Il secondo esempio mostra le due proprietà contemporaneamente: l'ordine LIFO e la
valutazione immediata degli argomenti.

\begin{lstlisting}[style=gocode,caption={\texttt{03\_control\_flow/defer\_lifo.go}}]
// defer_lifo.go - Ordine LIFO e valutazione immediata degli argomenti
// (slide 26).
package main

import "fmt"

func countdown() {
	// Le chiamate differite sono impilate: l'ultima registrata è la prima
	// eseguita. Ogni defer cattura il valore corrente di i (argomento
	// valutato subito).
	for i := 1; i <= 3; i++ {
		defer fmt.Println("defer", i)
	}
	fmt.Println("fine del corpo di countdown")
}

func main() {
	x := 1
	// L'argomento x viene valutato qui: la stampa mostrerà 1, non 2.
	defer fmt.Println("x al momento del defer =", x)
	x = 2
	fmt.Println("x ora =", x)

	countdown()
}
\end{lstlisting}
\begin{lstlisting}[style=output]
x ora = 2
fine del corpo di countdown
defer 3
defer 2
defer 1
x al momento del defer = 1
\end{lstlisting}

\subsection{Analisi del codice}

Dentro \cod{countdown}, il ciclo registra tre chiamate differite, con
\cod{i} uguale a 1, 2, 3 \emph{valutato al momento del defer}. Il corpo stampa
«fine del corpo di countdown»; solo \emph{allo scadere della funzione} le
chiamate sono eseguite in ordine inverso (3, 2, 1). In \cod{main}, il
\cod{defer} che stampa \cod{x} è registrato quando \cod{x} vale 1; anche se poi
\cod{x} diventa 2, la riga differita stampa 1, perché l'argomento era già stato
valutato. Essendo la prima registrata in \cod{main}, è l'ultima a essere
eseguita: compare in fondo all'output.

\begin{confronto}{C, Python, Java}
In \textbf{Java} la pulizia si fa con \cod{try}/\cod{finally} (o
\cod{try-with-resources}), in \textbf{Python} con \cod{try}/\cod{finally} o
\cod{with}. In \textbf{C} tipicamente con un'etichetta di \cod{goto cleanup}.
\cod{defer} offre il vantaggio di mettere il rilascio vicino all'acquisizione
senza annidare blocchi. Si noti che \cod{defer} si lega alla \emph{funzione}, non
al blocco: in un ciclo, i defer si accumulano fino alla fine della funzione.
\end{confronto}

\begin{attenzione}[Errore comune: defer dentro un ciclo]
Poiché \cod{defer} è eseguito solo all'uscita della \emph{funzione}, scrivere
\cod{defer f.Close()} dentro un ciclo che apre molti file li tiene tutti aperti
fino alla fine della funzione. In quei casi conviene racchiudere il corpo del
ciclo in una funzione ausiliaria, in modo che il defer scatti a ogni iterazione.
\end{attenzione}

% ===========================================================================
\chapter{Puntatori e strutture dati}\label{ch:dati}
% ===========================================================================

\section{Puntatori}\label{sec:puntatori}

\subsection{Introduzione}

\begin{definizione}{puntatore}
Un \textbf{puntatore} è un valore che memorizza l'\emph{indirizzo} di una
variabile. Gli operatori sono quelli usuali: \cod{\&} produce l'indirizzo di una
variabile; \cod{*} \emph{dereferenzia}, cioè recupera il valore a cui il
puntatore punta.
\end{definizione}

Poiché Go è \emph{call-by-value}, ogni argomento viene copiato: per
permettere a una funzione di modificare una variabile del chiamante occorre
passarne l'indirizzo.

\subsection{Funzionamento e differenze rispetto a C}

\begin{itemize}
  \item \textbf{Nessuna aritmetica dei puntatori}: \cod{p++} o \cod{p+1} non
        compilano (errore reale: \texttt{invalid operation: p++ (non-numeric type *int)}).
        Questo elimina un'intera classe di errori di accesso alla memoria.
  \item Il \textbf{valore zero} di un puntatore è \cod{nil}.
  \item È \textbf{sicuro restituire l'indirizzo di una variabile locale}: il
        runtime di Go fa sì che la variabile sopravviva alla funzione (se
        \emph{«scappa»} dallo scope, viene allocata sullo \emph{heap} anziché
        sullo stack).
\end{itemize}

\begin{lstlisting}[style=gocode,caption={\texttt{04\_data\_structures/pointers.go}}]
// pointers.go - Puntatori (slide 27).
package main

import "fmt"

// Go è call-by-value: la funzione riceve una copia del puntatore, che però
// punta alla stessa variabile del chiamante.
func incrementByValue(n int) { n++ }
func incrementByPointer(n *int) {
	*n++ // dereferenziazione: modifica la variabile originale
}

// È sicuro restituire l'indirizzo di una variabile locale: il runtime la
// alloca sullo heap (escape) e il garbage collector la libera quando
// nessuno la referenzia più. In C sarebbe un dangling pointer.
func newCounter() *int {
	c := 100
	return &c
}

func main() {
	var p *int
	fmt.Println("valore zero del puntatore:", p) // <nil>

	i := 1
	p = &i          // p punta a i
	fmt.Println(*p) // 1
	*p = 2          // equivale a i = 2
	fmt.Println(i)  // 2

	incrementByValue(i)
	fmt.Println("dopo incrementByValue:", i) // invariato
	incrementByPointer(&i)
	fmt.Println("dopo incrementByPointer:", i)

	q := newCounter()
	*q++
	fmt.Println("contatore sullo heap:", *q)

	// Nessuna aritmetica dei puntatori: p++ o p+1 non compilano.
}
\end{lstlisting}
\begin{lstlisting}[style=output]
valore zero del puntatore: <nil>
1
2
dopo incrementByValue: 2
dopo incrementByPointer: 3
contatore sullo heap: 101
\end{lstlisting}

\subsection{Analisi del codice}

\begin{itemize}
  \item \cod{p = \&i} fa puntare \cod{p} a \cod{i}; \cod{*p = 2} equivale
        a \cod{i = 2}.
  \item \cod{incrementByValue(i)} riceve una \emph{copia} di \cod{i} e la
        incrementa: il chiamante non vede alcun cambiamento (resta~2).
        \cod{incrementByPointer(\&i)} riceve una copia del \emph{puntatore},
        che però punta alla stessa variabile: \cod{*n++} la modifica (da 2 a 3).
  \item \cod{newCounter} restituisce \cod{\&c} dove \cod{c} è locale: è lecito
        perché \cod{c} viene portata sullo heap. In C lo stesso codice
        restituirebbe un \emph{dangling pointer} verso un frame di stack
        già distrutto.
\end{itemize}

\begin{attenzione}[Errore comune: dereferenziare un puntatore \cod{nil}]
Un puntatore non inizializzato vale \cod{nil}; dereferenziarlo causa un
\emph{panic} a runtime: \texttt{panic: runtime error: invalid memory address or
nil pointer dereference}. Come in C, il controllo \cod{if p != nil} va fatto
dove \cod{nil} è un valore possibile; a differenza di C il programma non prosegue
con memoria corrotta ma si arresta in modo controllato.
\end{attenzione}

\begin{confronto}{C, Python, Java}
Rispetto a \textbf{C}: stessa notazione \cod{\&} e \cod{*}, ma senza
aritmetica e con garbage collector (niente \cod{free}). Rispetto a \textbf{Java}
e \textbf{Python}, dove gli oggetti sono sempre referenziati implicitamente e
la distinzione tra valore e riferimento è decisa dal tipo, in Go è
\emph{il programmatore a scegliere}: un \cod{Vertex} passato per valore viene
copiato, un \cod{*Vertex} no.
\end{confronto}

\section{Strutture dati aggregate}

\subsection{Struct}

Una \textbf{struct} è una collezione tipizzata di \emph{campi}, di dimensione
fissa, con sintassi simile a quella del C. I campi si accedono con la notazione
\cod{.}; si possono accedere anche tramite un puntatore a struct
(\cod{p := \&v}), nel qual caso Go permette di scrivere \cod{p.X} al posto di
\cod{(*p).X}.

\begin{lstlisting}[style=gocode,caption={\texttt{04\_data\_structures/structs.go}}]
// structs.go - Struct (slide 28).
package main

import "fmt"

type Vertex struct {
	X int
	Y int
}

func main() {
	// Letterale con campi nominati; i campi omessi valgono zero.
	v := Vertex{X: 1, Y: 2}
	fmt.Println(v.X, v)

	// Un puntatore a struct: Go consente p.X al posto di (*p).X.
	p := &v
	p.X = 10
	fmt.Println(v.X) // 10: p punta a v

	// Le struct sono valori: l'assegnazione copia tutti i campi.
	w := v
	w.Y = 99
	fmt.Println(v, w) // v non cambia

	var zero Vertex
	fmt.Println(zero) // {0 0}
}
\end{lstlisting}
\begin{lstlisting}[style=output]
1 {1 2}
10
{10 2} {10 99}
{0 0}
\end{lstlisting}

\subsubsection*{Analisi del codice}

\cod{Vertex\{X: 1, Y: 2\}} è un \emph{letterale composito} con campi nominati;
quelli omessi valgono zero (\cod{\{0 0\}} per \cod{zero}). Dopo \cod{p := \&v}
l'assegnazione \cod{p.X = 10} modifica \cod{v}, perché \cod{p} punta a essa. Al
contrario \cod{w := v} \textbf{copia} tutti i campi: modificare \cod{w.Y} lascia
\cod{v} invariata.

\begin{confronto}{C, Python, Java}
La struct di Go è simile a quella di \textbf{C}: un valore, copiato
nelle assegnazioni. A differenza di una \emph{classe} Java o Python, la struct
non è un riferimento, non ha ereditarietà e i metodi si associano
separatamente (Capitolo~\ref{ch:oop}).
\end{confronto}

\subsection{Array}

Un \textbf{array} \cod{[n]T} è un array di \cod{n} elementi di tipo \cod{T}, di
\textbf{dimensione fissa} (non ridimensionabile). La lunghezza fa parte del tipo:
\cod{[2]string} e \cod{[3]string} sono tipi \emph{diversi}.

\begin{lstlisting}[style=gocode,caption={\texttt{04\_data\_structures/arrays.go}}]
// arrays.go - Array a dimensione fissa (slide 28).
package main

import "fmt"

func main() {
	// La lunghezza fa parte del tipo: [2]string e [3]string sono tipi
	// diversi.
	var a [2]string
	a[0] = "Hello"
	a[1] = "World"
	fmt.Println(a[0], a[1])
	fmt.Println(a)

	primes := [6]int{2, 3, 5, 7, 11, 13}
	fmt.Println(primes, len(primes))

	// Un array è un valore: l'assegnazione ne copia tutti gli elementi
	// (diverso da C, dove un array decade in puntatore).
	b := primes
	b[0] = 99
	fmt.Println(primes[0], b[0]) // 2 99
}
\end{lstlisting}
\begin{lstlisting}[style=output]
Hello World
[Hello World]
[2 3 5 7 11 13] 6
2 99
\end{lstlisting}

\begin{confronto}{C, Python, Java}
In \textbf{C} un array nelle espressioni «decade» in un puntatore al primo
elemento: passarlo a una funzione non lo copia. In \textbf{Go} un array è un
\emph{valore}: assegnarlo o passarlo copia tutti gli elementi (come mostra
\cod{b := primes}). In \textbf{Java} gli array sono oggetti referenziati.
Per questo in Go gli array si usano di rado direttamente: si usano le
\emph{slice}.
\end{confronto}

\section{Slice}\label{sec:slice}

\subsection{Introduzione}

La \textbf{slice} è il tipo di dato più importante di Go e molto più flessibile
dell'array. \cod{[]T} è una slice con elementi di tipo \cod{T}: una \emph{vista},
di dimensione dinamica, su un \textbf{array sottostante}.

\subsection{Creare una slice da un array}

Si crea una slice \emph{affettando} un array (o un'altra slice) con due indici,
\cod{s[i:j]}: il limite inferiore è \textbf{incluso}, quello superiore
\textbf{escluso}. Una slice \emph{riferisce} una sezione dell'array: modificare un
elemento della slice modifica l'elemento corrispondente dell'array.

\begin{lstlisting}[style=gocode,caption={\texttt{04\_data\_structures/slices.go}}]
// slices.go - Slice come vista su un array sottostante (slide 29-30).
package main

import "fmt"

func main() {
	primes := [6]int{2, 3, 5, 7, 11, 13}

	// s[i:j] include l'elemento i ed esclude j.
	var s []int = primes[1:4]
	fmt.Println(s, "len =", len(s), "cap =", cap(s)) // [3 5 7] len=3 cap=5

	// La slice NON copia i dati: modificare s modifica l'array.
	s[0] = 100
	fmt.Println("array dopo s[0]=100:", primes)

	// Due slice sulla stessa porzione si influenzano a vicenda.
	t := primes[:2]
	fmt.Println("t =", t) // [2 100]

	// La capacità arriva fino alla fine dell'array sottostante.
	u := s[:5]
	fmt.Println(u, len(u), cap(u))

	// Controllo dei limiti a runtime: fuori dalla capacità -> panic.
	defer func() {
		fmt.Println("recuperato:", recover())
	}()
	_ = s[:6] // slice bounds out of range
}
\end{lstlisting}
\begin{lstlisting}[style=output]
[3 5 7] len = 3 cap = 5
array dopo s[0]=100: [2 100 5 7 11 13]
t = [2 100]
[100 5 7 11 13] 5 5
recuperato: runtime error: slice bounds out of range [:6] with capacity 5
\end{lstlisting}

\subsection{Analisi del codice}

\begin{itemize}
  \item \cod{primes[1:4]} contiene gli indici 1, 2, 3, cioè \cod{[3 5 7]}
        (come nella slide): \cod{len = 3}, \cod{cap = 5}.
  \item \cod{s[0] = 100} modifica \cod{primes[1]}: la slice non contiene una
        copia dei dati. La slice \cod{t := primes[:2]}, definita dopo,
        vede il nuovo valore (\cod{[2 100]}).
  \item \cod{s[:5]} è consentito perché la \emph{capacità} di \cod{s} è 5
        (fino alla fine dell'array sottostante), anche se la lunghezza è 3.
  \item \cod{s[:6]} eccede la capacità: Go controlla i limiti e causa un
        panic, qui intercettato con \cod{recover} dentro un \cod{defer}
        solo per mostrare il messaggio
        \texttt{slice bounds out of range [:6] with capacity 5}.
\end{itemize}

\begin{nota}[Come è fatta una slice]
Una slice è un piccolo descrittore con tre campi: un \textbf{puntatore}
all'array sottostante, la \textbf{lunghezza} (\cod{len}) e la \textbf{capacità}
(\cod{cap}). Copiare una slice copia questo descrittore, non gli elementi. Ecco
perché due slice sullo stesso array si influenzano a vicenda e perché passare una
slice a una funzione è economico.
\end{nota}

\begin{confronto}{C, Python, Java}
In \textbf{Python} \cod{a[1:4]} crea una \emph{copia} della sezione; in Go la
slice \emph{condivide} la memoria. In \textbf{Java} il costrutto più simile a
una slice è la vista \cod{List.subList}, mentre un \cod{ArrayList} somiglia alla
slice per la crescita dinamica. In \textbf{C} si ottiene qualcosa di simile con la
coppia (puntatore, lunghezza), ma a mano e senza controlli.
\end{confronto}

\subsection{Lunghezza, capacità e \texttt{make}}

Una slice ha sia una \textbf{lunghezza} sia una \textbf{capacità}.
\cod{len(s)} è il numero di elementi che contiene; \cod{cap(s)} è il numero di
elementi dell'array sottostante a partire dall'inizio della slice. Go
effettua il \emph{bounds checking}: l'accesso fuori dai limiti causa un errore di
compilazione (se riconoscibile staticamente) o di runtime.

Una slice si può creare anche con \cod{make}, che permette di specificare
lunghezza e (facoltativamente) capacità: \cod{s := make([]int, 5, 10)}. L'array
sottostante viene creato automaticamente, ma non è direttamente visibile.

\begin{lstlisting}[style=gocode,caption={\texttt{04\_data\_structures/making\_slice.go}}]
// making_slice.go - Creare slice con make (slide 30-31).
package main

import "fmt"

func main() {
	a := make([]int, 0, 5) // len(s)=0, cap(s)=5
	printSlice("a", a)

	b := make([]int, 5) // len=5, cap=5: cinque zeri
	printSlice("b", b)

	c := make([]int, 5, 10)
	printSlice("c", c)

	d := c[2:4] // len=2, cap=8 (si parte dall'indice 2 dell'array)
	printSlice("d", d)

	var nilSlice []int // slice nil: len=0, cap=0, ma utilizzabile con append
	printSlice("nil", nilSlice)
	fmt.Println("nilSlice == nil?", nilSlice == nil)
}

func printSlice(s string, x []int) {
	fmt.Printf("%s len=%d cap=%d %v\n", s, len(x), cap(x), x)
}
\end{lstlisting}
\begin{lstlisting}[style=output]
a len=0 cap=5 []
b len=5 cap=5 [0 0 0 0 0]
c len=5 cap=10 [0 0 0 0 0]
d len=2 cap=8 [0 0]
nil len=0 cap=0 []
nilSlice == nil? true
\end{lstlisting}

\subsubsection*{Analisi del codice}

\cod{make([]int, 0, 5)} crea una slice vuota (\cod{len=0}) con spazio per 5
elementi (\cod{cap=5}): è il primo output, identico a quello della slide
(\cod{a len=0 cap=5 []}). Con \cod{make([]int, 5)} la capacità coincide con la
lunghezza e gli elementi sono cinque zeri. \cod{d := c[2:4]} ha \cod{len=2} e
\cod{cap=8}: la capacità si conta dall'inizio della slice (indice~2) fino alla
fine dell'array (10). Infine una slice dichiarata e non inizializzata è
\cod{nil}: ha \cod{len=0}, \cod{cap=0} ed è comunque utilizzabile con
\cod{append}.

\begin{attenzione}[Errore comune: confondere array e slice]
\cod{[3]int} è un array (tipo con lunghezza, \emph{valore}); \cod{[]int}
è una slice (tipo senza lunghezza, \emph{vista}). Un array assegnato si copia,
una slice assegnata condivide l'array sottostante. Un indice costante fuori
limite su un array è rifiutato dal compilatore (\texttt{invalid argument: index 5
out of bounds [0:3]}), mentre un indice variabile fuori limite dà un panic a
runtime (\texttt{index out of range [5] with length 3}).
\end{attenzione}

\subsection{L'operazione \texttt{append}}

Si aggiungono elementi a una slice con \cod{append}, la cui firma concettuale è
\cod{func append(slice []T, elems ...T) []T}. Quando si aggiunge a una slice, la
slice \emph{può essere ingrandita} se necessario.

\begin{itemize}
  \item Se \cod{len(s) < cap(s)}: i nuovi elementi sono aggiunti
        all'\textbf{array sottostante esistente}, senza crearne uno nuovo.
  \item Se \cod{len(s) == cap(s)}: Go \textbf{alloca un nuovo array più grande}
        (di norma il doppio della capacità precedente), \textbf{copia} gli
        elementi esistenti, e \textbf{restituisce una slice} che punta al nuovo
        array. Il vecchio array resta invariato se altre slice vi puntano ancora.
\end{itemize}

\begin{lstlisting}[style=gocode,caption={\texttt{04\_data\_structures/append.go}}]
// append.go - append e crescita della capacità (slide 32-33).
package main

import "fmt"

func main() {
	var s []int
	printSlice(s)

	s = append(s, 0) // funziona su slice nil
	printSlice(s)

	s = append(s, 1) // la slice cresce quando serve
	printSlice(s)

	s = append(s, 2, 3, 4) // più elementi in una volta
	printSlice(s)

	// Capacità piena: viene allocato un nuovo array e gli elementi sono
	// copiati.
	t := make([]int, 3, 3)
	t = append(t, 4)
	fmt.Println(t) // [0 0 0 4]
	printSlice(t)

	// Capacità disponibile: nessuna nuova allocazione, e l'array
	// sottostante è condiviso.
	base := make([]int, 3, 10)
	x := append(base, 1)
	y := append(base, 2)    // sovrascrive la stessa posizione di x!
	fmt.Println(x[3], y[3]) // 2 2
}

func printSlice(s []int) {
	fmt.Printf("len=%d cap=%d %v\n", len(s), cap(s), s)
}
\end{lstlisting}
\begin{lstlisting}[style=output]
len=0 cap=0 []
len=1 cap=1 [0]
len=2 cap=2 [0 1]
len=5 cap=6 [0 1 2 3 4]
[0 0 0 4]
len=4 cap=6 [0 0 0 4]
2 2
\end{lstlisting}

\subsection{Analisi del codice}

\begin{itemize}
  \item \cod{append} funziona anche su una slice \cod{nil}, come mostrano le
        prime righe.
  \item \cod{s = append(s, 2, 3, 4)} aggiunge tre elementi in una chiamata.
        Dopo l'aggiunta la lunghezza è 5, ma la capacità è \textbf{6}, non 5.
        La slide parla di «di norma il doppio»; la regola esatta è
        un dettaglio dell'implementazione: il runtime calcola la nuova
        capacità e poi la \emph{arrotonda} alla dimensione dei blocchi
        che l'allocatore fornisce. Per questo il valore osservato (6) può non
        coincidere con una semplice regola «raddoppia»: qui servivano 5
        posti, più del doppio della capacità precedente (2$\times$2), e l'arrotondamento
        ha portato a 6. Ciò che resta valido è la proprietà qualitativa: la
        capacità cresce a salti, non a ogni elemento, e ogni salto copia
        l'array. La regola esatta è un dettaglio dell'implementazione e non va
        data per scontata.
  \item Il caso della slide, \cod{make([]int, 3, 3)} seguito da
        \cod{append(s, 4)}, dà \cod{[0 0 0 4]}: la capacità era piena, quindi è
        stato creato un nuovo array.
  \item Nell'ultimo blocco, \cod{base} ha capacità 10 e lunghezza 3: entrambe le
        chiamate \cod{append(base, \ldots)} scrivono nella \emph{stessa} posizione
        dell'array condiviso (indice 3), quindi \cod{x[3]} e \cod{y[3]} valgono
        entrambi 2 --- la seconda chiamata ha sovrascritto la prima.
\end{itemize}

\begin{attenzione}[Errore comune: dimenticare di assegnare il risultato di \cod{append}]
\cod{append} può restituire una slice che punta a un array \emph{nuovo}, quindi il
risultato \textbf{va sempre riassegnato}. Scrivere \cod{append(s, 4)} senza
riassegnare è rifiutato dal compilatore (\texttt{append(s, 2) (value of type
[]int) is not used}). Un errore più subdolo si ha passando la slice a una
funzione che fa \cod{append}: la funzione riceve una \emph{copia} del
descrittore, quindi il chiamante non vede né la nuova lunghezza né il nuovo
array.
\end{attenzione}

\begin{lstlisting}[style=gocode,caption={\texttt{04\_data\_structures/slice\_pitfalls.go}}]
// slice_pitfalls.go - Errori comuni con array e slice.
package main

import "fmt"

// Errato: append può restituire una slice con un nuovo array; il
// risultato va riassegnato. Qui la modifica si perde per il chiamante,
// perché s è una copia dell'header (puntatore, len, cap).
func addWrong(s []int) {
	s = append(s, 99)
}

// Corretto: restituire la slice aggiornata.
func addRight(s []int) []int {
	return append(s, 99)
}

func main() {
	s := []int{1, 2, 3}
	addWrong(s)
	fmt.Println("dopo addWrong:", s) // [1 2 3]
	s = addRight(s)
	fmt.Println("dopo addRight:", s) // [1 2 3 99]

	// append senza riassegnare: il compilatore stesso lo rifiuta
	// ("append(s, 4) (value of type []int) is not used"). Va scritto
	// s = append(s, 4).

	// Una funzione può però modificare gli ELEMENTI esistenti, perché la
	// slice condivide l'array sottostante.
	double(s)
	fmt.Println("dopo double:", s)

	// Array e slice sono diversi: l'array si copia per valore.
	arr := [3]int{1, 2, 3}
	arrCopy := arr
	arrCopy[0] = 100
	fmt.Println(arr, arrCopy)
}

func double(s []int) {
	for i := range s {
		s[i] *= 2
	}
}
\end{lstlisting}
\begin{lstlisting}[style=output]
dopo addWrong: [1 2 3]
dopo addRight: [1 2 3 99]
dopo double: [2 4 6 198]
[1 2 3] [100 2 3]
\end{lstlisting}

\subsubsection*{Perché avviene}

\cod{addWrong} riceve una copia di \cod{(puntatore, len, cap)} e ne aggiorna la
\emph{copia}: la \cod{len} del chiamante resta 3. \cod{addRight} restituisce la
slice aggiornata, che il chiamante riassegna. La funzione \cod{double} mostra
invece che modificare gli \emph{elementi} funziona, perché tutte le copie del
descrittore puntano allo stesso array.

\section{Map}\label{sec:map}

\subsection{Introduzione e funzionamento}

Una \textbf{map} è una collezione di coppie chiave--valore. Il tipo
\cod{map[K]V} è un \emph{riferimento} a una tabella hash in cui \cod{K} e
\cod{V} sono i tipi di chiavi e valori. Si crea con \cod{make}:
\cod{m = make(map[string]Vertex)}.

Le operazioni principali sono:

\begin{lstlisting}[style=gocode,numbers=none]
m[key] = element      // inserimento o aggiornamento
elem := m[key]        // lettura
delete(m, key)        // cancellazione
elem, ok = m[key]     // test: ok è true se la chiave esiste
\end{lstlisting}

\begin{lstlisting}[style=gocode,caption={\texttt{04\_data\_structures/maps.go}}]
// maps.go - Map: tabelle hash chiave-valore (slide 34).
package main

import "fmt"

type Vertex struct {
	Lat, Long float64
}

func main() {
	// Una map va creata con make (o con un letterale): la map nil
	// permette letture ma non scritture.
	m := make(map[string]Vertex)
	m["Bell Labs"] = Vertex{40.68433, -74.39967} // inserimento
	fmt.Println(m["Bell Labs"])                  // lettura

	m["Bell Labs"] = Vertex{1, 2} // aggiornamento
	fmt.Println(m)

	// Chiave inesistente: si ottiene il valore zero, senza errore.
	fmt.Println("chiave assente:", m["Google"])

	// Idioma "comma ok": distingue "assente" da "presente con valore zero".
	elem, ok := m["Google"]
	fmt.Println(elem, ok) // {0 0} false
	elem, ok = m["Bell Labs"]
	fmt.Println(elem, ok) // {1 2} true

	delete(m, "Bell Labs") // cancellazione (nessun errore se assente)
	fmt.Println(len(m))

	// Errore comune: scrivere su una map nil provoca panic.
	var nilMap map[string]int
	fmt.Println(nilMap["x"]) // lettura ok: 0
	defer func() { fmt.Println("recuperato:", recover()) }()
	nilMap["x"] = 1 // panic: assignment to entry in nil map
}
\end{lstlisting}
\begin{lstlisting}[style=output]
{40.68433 -74.39967}
map[Bell Labs:{1 2}]
chiave assente: {0 0}
{0 0} false
{1 2} true
0
0
recuperato: assignment to entry in nil map
\end{lstlisting}

\subsection{Analisi del codice}

\begin{itemize}
  \item Leggere una chiave \emph{inesistente} \textbf{non è un errore}: si ottiene il
        \emph{valore zero} del tipo dei valori (qui \cod{\{0 0\}}).
  \item Per distinguere «chiave assente» da «chiave presente con valore zero» si
        usa l'idioma \emph{comma ok}: \cod{elem, ok := m[key]}.
  \item \cod{delete} su una chiave assente non fa nulla.
  \item Scrivere su una map \cod{nil} causa un panic
        (\texttt{assignment to entry in nil map}); la lettura da una map \cod{nil}
        è invece lecita e restituisce il valore zero.
\end{itemize}

\begin{attenzione}[Errore comune: accesso a una chiave inesistente]
Poiché la lettura di una chiave assente restituisce silenziosamente il valore
zero, un contatore \cod{m["x"]} sembra funzionare ma una mappa di booleani
non distingue «falso» da «assente»: serve \cod{v, ok := m[k]}. Ricordare inoltre di
\emph{creare} la map con \cod{make} (o con un letterale) prima di scrivere.
\end{attenzione}

\begin{confronto}{C, Python, Java}
In \textbf{Python} una chiave assente solleva \cod{KeyError} (o si usa
\cod{.get}); in \textbf{Java} \cod{HashMap.get} restituisce \cod{null}; in
\textbf{Go} si ottiene il valore zero e si usa \cod{ok} per sapere se la chiave
c'era. In \textbf{C} non esiste una map nella libreria standard. Una differenza
da non trascurare: l'ordine di iterazione di una map Go \emph{non è
specificato} (è anzi deliberatamente casuale), mentre un \cod{dict} Python
mantiene l'ordine di inserimento.
\end{confronto}

\section{L'istruzione \texttt{range}}\label{sec:range}

\cod{range} itera sugli elementi di varie strutture dati (array, slice,
stringhe, map e channel); il comportamento \emph{dipende dal tipo}:

\begin{itemize}
  \item array, slice, stringhe: \cod{range} fornisce \textbf{indice e valore} di
        ogni elemento;
  \item map: fornisce \textbf{chiave e valore};
  \item channel: fornisce i valori ricevuti finché il channel non è chiuso
        (Sezione~\ref{sec:chanrange}).
\end{itemize}

\begin{lstlisting}[style=gocode,caption={\texttt{04\_data\_structures/range.go}}]
// range.go - Iterare con range su una slice (slide 35).
package main

import "fmt"

var pow = []int{1, 2, 4, 8, 16, 32, 64, 128}

func main() {
	// Per array e slice, range produce indice e copia del valore.
	for i, v := range pow {
		fmt.Printf("2**%d = %d\n", i, v)
	}
}
\end{lstlisting}
\begin{lstlisting}[style=output]
2**0 = 1
2**1 = 2
2**2 = 4
2**3 = 8
2**4 = 16
2**5 = 32
2**6 = 64
2**7 = 128
\end{lstlisting}

L'esempio successivo mostra come \textbf{ignorare} l'indice o il valore con il
\emph{blank identifier} \cod{\_}, e come iterare su una map.

\begin{lstlisting}[style=gocode,caption={\texttt{04\_data\_structures/range2.go}}]
// range2.go - range: ignorare indice o valore, iterare su map (slide 36).
package main

import "fmt"

func main() {
	nums := []int{2, 3, 4}
	sum := 0
	// "_" (blank identifier) scarta l'indice: Go vieta variabili inutilizzate.
	for _, num := range nums {
		sum += num
	}
	fmt.Println("sum:", sum)

	for i, num := range nums {
		if num == 3 {
			fmt.Println("index:", i)
		}
	}

	// Su una map range produce chiave e valore. L'ordine di iterazione NON
	// è specificato (in Go è volutamente casuale): un ciclo diretto
	// "for k, v := range kvs" stampa le due coppie in ordine variabile.
	// Per avere un output stabile si itera su un elenco di chiavi
	// ordinato a mano e si consulta la map.
	kvs := map[string]string{"a": "apple", "b": "banana"}
	for _, k := range []string{"a", "b"} {
		fmt.Printf("%s -> %s\n", k, kvs[k])
	}

	// Con la sola chiave si omette il secondo valore.
	n := 0
	for k := range kvs {
		_ = k
		n++
	}
	fmt.Println("chiavi iterate:", n)

	// Anche le stringhe: range restituisce indice di byte e rune (Unicode).
	for i, r := range "gò" {
		fmt.Printf("%d:%c ", i, r)
	}
	fmt.Println()
}
\end{lstlisting}
\begin{lstlisting}[style=output]
sum: 9
index: 1
a -> apple
b -> banana
chiavi iterate: 2
0:g 1:ò 
\end{lstlisting}

\subsection{Analisi del codice}

\begin{itemize}
  \item \cod{for \_, num := range nums} scarta l'indice; Go vieta le variabili
        dichiarate e inutilizzate, quindi \cod{\_} è la forma corretta per
        dire «non mi serve».
  \item \cod{for k := range kvs} itera sulle sole chiavi (valore omesso).
  \item \textbf{Ordine delle map:} la slide mostra l'output
        \cod{a -> apple}, \cod{b -> banana}; ma con un ciclo
        \cod{for k, v := range kvs} diretto le due coppie possono comparire in
        \emph{qualunque ordine}, anche diverso da un'esecuzione all'altra. Per questo
        l'esempio della guida scorre un elenco di chiavi ordinato a mano
        e consulta la map: l'output è deterministico.
  \item Su una stringa \cod{range} restituisce la posizione in \emph{byte} e la
        \emph{rune} (punto di codice Unicode): per \cod{"gò"} gli indici sono 0 e
        1 e il carattere \cod{ò} occupa due byte UTF-8.
\end{itemize}

% ===========================================================================
\chapter{Metodi, interfacce ed errori}\label{ch:oop}
% ===========================================================================

Go \textbf{non ha classi}. Il comportamento si associa ai tipi con i
\emph{metodi}, e le astrazioni si esprimono con le \emph{interfacce}, che
descrivono ciò che un tipo \emph{sa fare} e non ciò che \emph{contiene}.

\section{Metodi}

\subsection{Introduzione e sintassi}

Go supporta i \textbf{metodi} sui tipi concreti, comprese le struct.

\begin{definizione}{metodo e receiver}
Un \textbf{metodo} è una funzione con un argomento speciale, il \emph{receiver}
(ricevitore): un parametro aggiuntivo scritto \emph{tra} la parola chiave
\cod{func} e il nome del metodo. Associa la funzione a un tipo specifico
(per esempio una struct) e permette di definire il comportamento del tipo.
\end{definizione}

\begin{lstlisting}[style=gocode,numbers=none,caption={Dalla slide: un metodo \cod{Abs} su \cod{Vertex}}]
type Vertex struct {
	X, Y float64
}

func (v Vertex) Abs() float64 {
	return math.Sqrt(v.X*v.X + v.Y*v.Y)
}
\end{lstlisting}

\subsection{Funzionamento: receiver per valore e per puntatore}

Il receiver è, a tutti gli effetti, un \emph{parametro}: poiché Go passa tutto
\emph{per valore} (Capitolo~\ref{ch:funzioni}), un receiver di tipo \cod{Vertex}
riceve una \textbf{copia} della struct. Se invece il receiver è un puntatore
(\cod{v *Vertex}) il metodo lavora sull'originale e può modificarlo. Le slide
usano la forma a puntatore negli esempi sul \emph{mutex} (\cod{func (l *Lock)
Lock()}, Capitolo~\ref{ch:conc}), dove è indispensabile.

\begin{lstlisting}[style=gocode,caption={\texttt{05\_methods\_interfaces/methods.go}}]
// methods.go - Metodi su tipi concreti (slide 41).
package main

import (
	"fmt"
	"math"
)

type Vertex struct {
	X, Y float64
}

// Un metodo è una funzione con un receiver, scritto tra "func" e il nome.
// Con receiver per valore (v Vertex) il metodo lavora su una COPIA.
func (v Vertex) Abs() float64 {
	return math.Sqrt(v.X*v.X + v.Y*v.Y)
}

// Con receiver puntatore (v *Vertex) il metodo può modificare l'originale.
func (v *Vertex) Scale(f float64) {
	v.X *= f
	v.Y *= f
}

// Errore comune: receiver per valore, la modifica agisce sulla copia
// e va persa.
func (v Vertex) ScaleWrong(f float64) {
	v.X *= f
	v.Y *= f
}

func main() {
	v := Vertex{3, 4}
	fmt.Println(v.Abs()) // 5

	v.ScaleWrong(10)
	fmt.Println("dopo ScaleWrong:", v) // invariato

	// v è una variabile indirizzabile: Go traduce v.Scale(10) in
	// (&v).Scale(10).
	v.Scale(10)
	fmt.Println("dopo Scale:", v, v.Abs())
}
\end{lstlisting}
\begin{lstlisting}[style=output]
5
dopo ScaleWrong: {3 4}
dopo Scale: {30 40} 50
\end{lstlisting}

\subsection{Analisi del codice}

\cod{Abs} usa un receiver per valore: legge \cod{X} e \cod{Y} senza modificarli
e restituisce $\sqrt{3^2+4^2}=5$. \cod{ScaleWrong} ha lo stesso receiver per
valore ma tenta di modificare i campi: poiché agisce su una copia, \cod{v} nel
\cod{main} rimane \cod{\{3 4\}}. \cod{Scale}, con receiver \cod{*Vertex},
modifica l'originale: dopo \cod{v.Scale(10)} il vettore vale \cod{\{30 40\}} e
\cod{Abs()} dà 50. Si noti che nel \cod{main} si scrive semplicemente
\cod{v.Scale(10)}: poiché \cod{v} è una variabile, Go la traduce
automaticamente in \cod{(\&v).Scale(10)}.

\begin{attenzione}[Errore comune: modificare con un receiver per valore]
Dimenticare l'asterisco nel receiver di un metodo che deve cambiare lo stato è
un errore classico, spesso silenzioso: il codice compila e funziona, ma le
modifiche vanno perse con la copia. La regola pratica: se il metodo deve
\emph{modificare} il valore, serve un receiver puntatore.
\end{attenzione}

\begin{confronto}{C, Python, Java}
In \textbf{Java} e \textbf{Python} un metodo riceve il riferimento all'oggetto
(\cod{this}/\cod{self}) e può sempre modificarne lo stato. In \textbf{Go} la
scelta tra receiver per valore e per puntatore è esplicita e cambia la
semantica. Rispetto a \textbf{C}, un metodo Go è «zucchero sintattico» per una
funzione che prende la struct come primo parametro: \cod{v.Abs()} corrisponde
grosso modo a \cod{Vertex\_Abs(v)}. Il receiver è dichiarato \emph{fuori} dalla
definizione del tipo: i metodi non vivono dentro un blocco \cod{class}.
\end{confronto}

\section{Interfacce}\label{sec:interfacce}

\subsection{Introduzione}

\begin{definizione}{interfaccia}
Un'\textbf{interfaccia} è una collezione \emph{nominata} di firme di metodi.
\end{definizione}

Il punto più importante è l'\textbf{implementazione implicita}: in Go un tipo
\emph{soddisfa} un'interfaccia semplicemente implementandone \emph{tutti} i
metodi. Non serve alcuna dichiarazione esplicita (nessun \cod{implements}): il
tipo soddisfa automaticamente l'interfaccia.

L'obiettivo è progettare astrazioni basate sul \textbf{comportamento} (le
azioni) e non sui dati. Le interfacce abilitano il \textbf{polimorfismo} in Go.

\subsection{Esempio pratico: le forme geometriche}

\begin{lstlisting}[style=gocode,caption={\texttt{05\_methods\_interfaces/interface.go}}]
// interface.go - Interfacce implicite e polimorfismo (slide 42-45).
package main

import (
	"fmt"
	"math"
)

// Define the interface for geometric shapes
type geometry interface {
	area() float64
	perim() float64
}

// Implement the interface on rectangular and circle types
type rect struct {
	width, height float64
}
type circle struct {
	radius float64
}

// Per implementare un'interfaccia in Go basta implementare tutti i suoi
// metodi: nessuna dichiarazione "implements".

// Implement area and perim methods on rect
func (r rect) area() float64 {
	return r.width * r.height
}
func (r rect) perim() float64 {
	return 2*r.width + 2*r.height
}

// Implement area and perim methods on circle
func (c circle) area() float64 {
	return math.Pi * c.radius * c.radius
}
func (c circle) perim() float64 {
	return 2 * math.Pi * c.radius
}

// Se una variabile ha tipo interfaccia si possono chiamare i metodi
// dell'interfaccia: measure funziona con qualsiasi tipo che implementi
// geometry.
func measure(g geometry) {
	fmt.Println(g)
	fmt.Println(g.area())
	fmt.Println(g.perim())
}

func main() {
	r := rect{width: 3, height: 4}
	c := circle{radius: 5}

	// rect e circle soddisfano entrambe geometry, quindi possono essere
	// passate a measure.
	measure(r)
	measure(c)
}
\end{lstlisting}
\begin{lstlisting}[style=output]
{3 4}
12
14
{5}
78.53981633974483
31.41592653589793
\end{lstlisting}

\subsection{Analisi del codice}

\begin{itemize}
  \item \cod{geometry} dichiara due metodi: \cod{area()} e \cod{perim()}.
  \item \cod{rect} e \cod{circle} sono struct che definiscono entrambi i
        metodi con i giusti receiver: \emph{nient'altro} le collega a
        \cod{geometry}.
  \item \cod{measure(g geometry)} accetta qualunque tipo che implementi
        \cod{geometry}; dentro la funzione si possono chiamare solo i metodi
        dell'interfaccia. La chiamata \cod{g.area()} è risolta a runtime sul tipo
        concreto contenuto in \cod{g}.
  \item \cod{fmt.Println(g)} stampa il valore concreto: \cod{\{3 4\}} per
        \cod{rect} e \cod{\{5\}} per \cod{circle}.
  \item Output per \cod{r} (3$\times$4): area 12, perimetro 14; per \cod{c}
        (raggio 5): area $25\pi\approx78.5398$, perimetro $10\pi\approx31.4159$.
\end{itemize}

Un secondo esempio mostra il polimorfismo in una collezione eterogenea:

\begin{lstlisting}[style=gocode,caption={\texttt{05\_methods\_interfaces/interface\_slice.go}}]
// interface_slice.go - Polimorfismo: una collezione eterogenea di forme.
//
// Esempio aggiuntivo (non presente nelle slide) che riusa solo concetti già
// visti: interfacce implicite, metodi e range su slice.
package main

import (
	"fmt"
	"math"
)

type shape interface {
	area() float64
}

type square struct{ side float64 }
type circle struct{ radius float64 }

func (s square) area() float64 { return s.side * s.side }
func (c circle) area() float64 { return math.Pi * c.radius * c.radius }

// totalArea somma le aree di tipi diversi: a runtime ogni chiamata
// s.area() viene risolta sul tipo concreto contenuto nell'interfaccia.
func totalArea(shapes []shape) float64 {
	total := 0.0
	for _, s := range shapes {
		total += s.area()
	}
	return total
}

func main() {
	shapes := []shape{square{2}, circle{1}, square{3}}
	fmt.Printf("area totale = %.4f\n", totalArea(shapes))
}
\end{lstlisting}
\begin{lstlisting}[style=output]
area totale = 16.1416
\end{lstlisting}

La somma è $4+\pi+9\approx16.1416$: \cod{[]shape} contiene tipi diversi, e a
ogni iterazione \cod{s.area()} invoca l'implementazione giusta.

\begin{attenzione}[Errore comune: un metodo mancante]
Se un tipo non implementa \emph{tutti} i metodi, il compilatore rifiuta
l'assegnazione all'interfaccia, spiegando quale metodo manca. Esempio con un
tipo \cod{T} che implementa solo \cod{A()} assegnato a un'interfaccia che richiede
\cod{A()} e \cod{B()}:
\begin{lstlisting}[style=output,aboveskip=0.3em,belowskip=0.3em]
cannot use T{} (value of struct type T) as I value in variable declaration:
T does not implement I (missing method B)
\end{lstlisting}
Il controllo è fatto \emph{a compile time}, quindi l'implementazione implicita
non sacrifica la sicurezza dei tipi.
\end{attenzione}

\begin{confronto}{C, Python, Java}
In \textbf{Java} un tipo deve dichiarare esplicitamente \cod{implements
Interfaccia}, e quindi l'autore della classe deve conoscere l'interfaccia in
anticipo. In \textbf{Go} la relazione è \emph{strutturale}: si può definire
un'interfaccia \emph{dopo}, anche in un altro package, e i tipi esistenti la
soddisfano se hanno i metodi giusti. In \textbf{Python} il \emph{duck typing}
ha un'idea simile («se cammina come un'anatra\ldots»), ma il controllo avviene a
runtime; in Go è statico. In \textbf{C} un'astrazione simile richiede tabelle
di puntatori a funzione costruite a mano.
\end{confronto}

\section{Gestione degli errori}\label{sec:errori}

\subsection{Introduzione}

Go usa \textbf{valori di errore} per indicare stati anomali, non eccezioni. Gli
errori sono comunicati con un valore di ritorno \emph{separato}. Come ha scritto
R.~Pike (citato nelle slide): «\emph{Error handling [in Go] does not obscure the
flow of control}», cioè la gestione degli errori non nasconde il flusso di
controllo.

\subsection{Funzionamento}

\begin{itemize}
  \item L'errore è, per \textbf{convenzione}, l'\emph{ultimo} valore restituito.
  \item Un errore \cod{nil} significa che \textbf{non si è verificato nessun errore}.
  \item Il chiamante controlla subito:
\end{itemize}

\begin{lstlisting}[style=gocode,numbers=none]
result, err := SomeFunction()
if err != nil {
	// gestisci l'errore
}
\end{lstlisting}

Il tipo \cod{error} è un'\textbf{interfaccia predefinita}:

\begin{lstlisting}[style=gocode,numbers=none]
type error interface {
	Error() string
}
\end{lstlisting}

Qualunque tipo con un metodo \cod{Error() string} soddisfa \cod{error}, per la
regola dell'implementazione implicita appena vista. \cod{errors.New} costruisce
un errore base con un messaggio.

\begin{lstlisting}[style=gocode,caption={\texttt{05\_methods\_interfaces/errors.go}}]
// errors.go - Gestione degli errori e interfaccia error (slide 69).
package main

import (
	"errors"
	"fmt"
)

// L'interfaccia predefinita error è:
//
//	type error interface { Error() string }
//
// Qualunque tipo con un metodo Error() string la soddisfa.

// divide segnala la condizione anomala con un errore invece che con
// un'eccezione. L'errore è l'ultimo valore restituito.
func divide(a, b float64) (float64, error) {
	if b == 0 {
		// errors.New costruisce un errore base con un messaggio.
		return 0, errors.New("divisione per zero")
	}
	return a / b, nil // nil = nessun errore
}

// Un tipo errore personalizzato: porta con sé dati strutturati.
type rangeError struct {
	value, max int
}

func (e rangeError) Error() string {
	return fmt.Sprintf("valore %d oltre il massimo %d", e.value, e.max)
}

func check(v int) error {
	if v > 10 {
		return rangeError{v, 10}
	}
	return nil
}

func main() {
	result, err := divide(10, 4)
	if err != nil {
		fmt.Println("errore:", err)
		return
	}
	fmt.Println("risultato:", result)

	if _, err := divide(1, 0); err != nil {
		fmt.Println("errore:", err)
	}

	for _, v := range []int{5, 42} {
		if err := check(v); err != nil {
			fmt.Println("check fallito:", err)
		} else {
			fmt.Println("check ok:", v)
		}
	}
}
\end{lstlisting}
\begin{lstlisting}[style=output]
risultato: 2.5
errore: divisione per zero
check ok: 5
check fallito: valore 42 oltre il massimo 10
\end{lstlisting}

\subsection{Analisi del codice}

\begin{itemize}
  \item \cod{divide} restituisce \cod{(float64, error)}: in caso di divisore
        nullo restituisce un errore creato con \cod{errors.New}, altrimenti
        \cod{nil}. Dividere 10 per 4 dà 2,5; dividere per zero produce
        l'errore, che il chiamante stampa.
  \item \cod{rangeError} è un tipo \emph{definito dal programmatore} che
        implementa \cod{error} con il metodo \cod{Error()}: così l'errore può
        trasportare \emph{dati strutturati} (qui il valore e il massimo).
  \item Il pattern \cod{if err := check(v); err != nil \{ \ldots\ \}} usa
        l'istruzione di inizializzazione dell'\cod{if} (Sezione~\ref{sec:if}).
\end{itemize}

\begin{attenzione}[Errore comune: ignorare l'errore]
Scrivere \cod{data, \_ := os.ReadFile(...)} scarta l'errore e prosegue con un
valore probabilmente vuoto. Poiché il linguaggio non obbliga a gestirlo
(non esistono \cod{try}/\cod{catch}), la disciplina del \cod{if err != nil}
subito dopo la chiamata è la regola d'oro del codice Go.
\end{attenzione}

\begin{confronto}{C, Python, Java}
In \textbf{C} gli errori si segnalano con codici di ritorno o con \cod{errno}.
In \textbf{Java} e \textbf{Python} con eccezioni (\cod{throw}/\cod{raise},
\cod{try}/\cod{catch}/\cod{except}), che interrompono il flusso e risalgono lo
stack finché qualcuno le cattura. Go torna ai valori di ritorno del C, ma
con un tipo dedicato (\cod{error}) e risultati multipli: il flusso di
controllo resta visibile nel codice. Il costo è un po' di verbosità
(\cod{if err != nil} ripetuto); il beneficio è che ogni punto che può fallire è
evidente.
\end{confronto}

% ===========================================================================
\chapter{Concorrenza}\label{ch:conc}
% ===========================================================================

\section{La concorrenza in Go}

Go fornisce la \textbf{concorrenza come parte del linguaggio di base}, non come
libreria aggiuntiva. Gli strumenti sono due: le \textbf{goroutine}, che eseguono
codice concorrentemente, e i \textbf{channel}, con cui le goroutine comunicano.
Insieme implementano il modello \textbf{CSP} (Capitolo~\ref{ch:go}) e permettono di
realizzare diversi \emph{pattern} di concorrenza.

Il principio guida è: invece di far condividere memoria a più flussi di
esecuzione proteggendola con lock, si fa in modo che i flussi \emph{si scambino
valori attraverso channel}. Go fornisce comunque anche i lock tradizionali
(\cod{sync.Mutex}), trattati nella Sezione~\ref{sec:mutex}.

\begin{attenzione}[Concorrenza e determinismo]
I programmi concorrenti di questo capitolo sono stati eseguiti più volte. Dove
l'output \emph{può cambiare} da un'esecuzione all'altra (ordine di stampa
delle goroutine, scelta casuale di \cod{select}, risultato di una data race) il
testo lo dichiara esplicitamente e l'output mostrato è \emph{uno dei
possibili}. In tutti gli altri casi l'output è deterministico.
\end{attenzione}

\section{Goroutine}\label{sec:goroutine}

\subsection{Introduzione}

\begin{definizione}{goroutine}
Una \textbf{goroutine} è un \emph{thread leggero} gestito dal runtime di Go.
\end{definizione}

Si avvia anteponendo \cod{go} a una chiamata di funzione:

\begin{lstlisting}[style=gocode,numbers=none]
go f(x, y, z)    // avvia una nuova goroutine che esegue f(x, y, z)
\end{lstlisting}

Gli argomenti sono valutati nella goroutine \emph{chiamante}; l'esecuzione di
\cod{f} avviene nella nuova goroutine, e il chiamante prosegue senza attendere.

\subsection{Funzionamento}

\begin{itemize}
  \item Le goroutine \textbf{condividono lo stesso spazio di indirizzi}: per
        accedere a memoria condivisa serve sincronizzazione.
  \item \textbf{Attenzione}: quando \cod{main} ritorna, il programma termina
        \emph{immediatamente}, senza attendere il completamento delle altre
        goroutine.
\end{itemize}

\subsubsection{Goroutine contro thread}

Le goroutine \emph{non sono thread} del sistema operativo, ma astrazioni
leggere costruite sopra di essi.

\begin{itemize}
  \item Lo \textbf{scheduler di Go} gestisce le goroutine, assegnandole ai
        thread del sistema operativo: \emph{un thread OS può eseguire molte
        goroutine}.
  \item Le goroutine sono molto più economiche da creare e distruggere dei
        thread OS e consumano molta meno memoria. Le slide indicano come ordine
        di grandezza circa $5\times$ in tempo di creazione/distruzione e circa
        $500\times$ in memoria; sono valori \emph{riportati dalle slide}, non
        misurati in questa guida, e dipendono da piattaforma e configurazione.
\end{itemize}

\subsubsection{Ordine di esecuzione e priorità}

Le goroutine \textbf{non hanno garanzie sull'ordine di esecuzione} rispetto
all'ordine in cui sono state dichiarate, hanno tutte la \textbf{stessa priorità}
e \textbf{non si può controllare} l'ordine con cui vengono eseguite.

\subsubsection{Controllare le prestazioni: \texttt{GOMAXPROCS}}

Le slide indicano che si può impostare \cod{GOMAXPROCS} per controllare quanti
thread del sistema operativo Go usa \emph{simultaneamente}, tipicamente il
numero di core virtuali della CPU. Si veda l'Appendice~B per una precisazione sul
nome: \cod{GOMAXPROCS} è una \emph{variabile d'ambiente}, mentre
\cod{runtime.GOMAXPROCS(n)} è la \emph{funzione} del package \cod{runtime} che
fa la stessa cosa da codice (con \cod{n = 0} restituisce solo il valore corrente).

\subsection{Esempio: \texttt{main} termina prima delle goroutine}

\begin{lstlisting}[style=gocode,caption={\texttt{06\_concurrency/goroutine\_termination.go}}]
// goroutine_termination.go - ERRORE COMUNE: main termina prima delle
// goroutine (slide 47).
//
// Quando main ritorna, il programma termina immediatamente senza attendere
// le altre goroutine. L'output di questo programma è NON deterministico:
// dopo la riga di main possono comparire da 0 a 3 righe "goroutine N"
// (anche prima della riga di main, in ordine qualsiasi). Nelle prove fatte
// sono comparse quantità diverse da un'esecuzione all'altra: non c'è
// alcuna garanzia.
package main

import "fmt"

func say(id int) {
	fmt.Println("goroutine", id)
}

func main() {
	for i := 1; i <= 3; i++ {
		go say(i) // avvia la goroutine e prosegue subito
	}
	fmt.Println("main termina senza attendere")
	// Nessuna sincronizzazione: il programma esce qui.
	// Soluzione: vedi goroutine_wait.go (sync.WaitGroup).
}
\end{lstlisting}

Questo programma ha un output \textbf{non deterministico}. In sei esecuzioni
si sono osservati esiti diversi; tre di essi sono:

\begin{lstlisting}[style=output]
# esecuzione A
main termina senza attendere

# esecuzione B
goroutine 3
goroutine 1
goroutine 2
main termina senza attendere

# esecuzione C
goroutine 1
main termina senza attendere
goroutine 2
\end{lstlisting}

\subsubsection*{Perché avviene}

L'istruzione \cod{go say(i)} \emph{non attende} che \cod{say} venga eseguita. Dopo
aver avviato le tre goroutine, \cod{main} stampa il proprio messaggio e ritorna: a
quel punto il processo termina, e le goroutine che non hanno ancora avuto
il tempo di stampare vengono abbandonate. Quante ne riescano a stampare
dipende da come lo scheduler le assegna ai thread: da zero a tre, in ordine
qualsiasi. Il programma è quindi \emph{scorretto}, anche se in alcune
esecuzioni sembra funzionare.

\subsection{La correzione: attendere con \texttt{sync.WaitGroup}}

Il rimedio è sincronizzare \cod{main} con le goroutine. Lo strumento è
\cod{sync.WaitGroup} (Sezione~\ref{sec:waitgroup}).

\begin{lstlisting}[style=gocode,caption={\texttt{06\_concurrency/goroutine\_wait.go}}]
// goroutine_wait.go - Versione corretta di goroutine_termination.go:
// main attende le goroutine con sync.WaitGroup (slide 47, 64).
//
// L'ORDINE con cui le goroutine stampano non è garantito: tutte hanno la
// stessa priorità e lo scheduler decide. Quindi l'output varia da
// un'esecuzione all'altra (ma contiene sempre le tre righe).
package main

import (
	"fmt"
	"sync"
)

func say(id int, wg *sync.WaitGroup) {
	defer wg.Done() // decrementa il contatore all'uscita, anche in caso di return anticipato
	fmt.Println("goroutine", id)
}

func main() {
	var wg sync.WaitGroup
	for i := 1; i <= 3; i++ {
		wg.Add(1) // PRIMA di avviare la goroutine, non al suo interno
		go say(i, &wg)
	}
	wg.Wait() // blocca finché il contatore torna a zero
	fmt.Println("main: tutte le goroutine sono terminate")
}
\end{lstlisting}

Ogni esecuzione stampa sempre le tre righe delle goroutine \emph{prima} della
riga finale di \cod{main}, ma l'ordine delle tre goroutine varia. Due
esecuzioni:

\begin{lstlisting}[style=output]
goroutine 3
goroutine 1
goroutine 2
main: tutte le goroutine sono terminate

goroutine 2
goroutine 1
goroutine 3
main: tutte le goroutine sono terminate
\end{lstlisting}

\subsubsection*{Analisi del codice}

\begin{itemize}
  \item \cod{wg.Add(1)} incrementa il contatore \emph{prima} di avviare la
        goroutine. Se \cod{Add} fosse chiamata \emph{dentro} la goroutine,
        \cod{wg.Wait()} potrebbe eseguirsi quando il contatore è ancora zero
        e ritornare subito.
  \item \cod{defer wg.Done()} decrementa il contatore all'uscita di \cod{say},
        anche in caso di \cod{return} anticipato.
  \item \cod{wg.Wait()} blocca \cod{main} finché il contatore torna a zero.
  \item Il \cod{*sync.WaitGroup} è passato \emph{per puntatore}: se fosse
        passato per valore, ogni goroutine decrementerebbe una \emph{copia} e
        \cod{Wait} non terminerebbe mai.
\end{itemize}

\subsection{Molte goroutine, pochi thread}

Il programma seguente avvia 10\,000 goroutine, ben più dei thread OS disponibili.
Lo scheduler le distribuisce sui thread. Ogni goroutine scrive in una posizione
\emph{diversa} della slice, quindi non ci sono accessi concorrenti allo stesso
elemento. Il valore di \cod{GOMAXPROCS} dipende dalla macchina: per rendere
l'output identico ovunque si stampa solo un controllo di coerenza.

\begin{lstlisting}[style=gocode,caption={\texttt{06\_concurrency/gomaxprocs.go}}]
// gomaxprocs.go - Goroutine leggere e thread del sistema operativo
// (slide 47-48).
//
// runtime.GOMAXPROCS(0) restituisce il numero massimo di thread OS che
// possono eseguire codice Go contemporaneamente, senza modificarlo. Il
// valore dipende dalla macchina, quindi viene stampato solo un controllo
// di coerenza per mantenere l'output identico ovunque.
//
// Qui si avviano 10000 goroutine, molte più dei thread disponibili: lo
// scheduler di Go le multiplexa sui thread OS. Con 10000 thread OS reali
// il costo in memoria e tempo di creazione sarebbe molto maggiore.
package main

import (
	"fmt"
	"runtime"
	"sync"
)

func main() {
	fmt.Println("GOMAXPROCS(0) >= 1:", runtime.GOMAXPROCS(0) >= 1)

	const n = 10000
	finished := make([]bool, n) // una posizione per goroutine: nessuna race
	var wg sync.WaitGroup
	for i := 0; i < n; i++ {
		wg.Add(1)
		go func() {
			defer wg.Done()
			finished[i] = true
		}()
	}
	wg.Wait()

	count := 0
	for _, f := range finished {
		if f {
			count++
		}
	}
	fmt.Println("goroutine completate:", count)
}
\end{lstlisting}
\begin{lstlisting}[style=output]
GOMAXPROCS(0) >= 1: true
goroutine completate: 10000
\end{lstlisting}

\begin{confronto}{C, Python, Java}
Un thread POSIX (\cod{pthread}) o un thread Java \emph{tradizionale} corrisponde
a un thread del kernel, con uno stack di dimensione tipicamente dell'ordine del
megabyte: crearne decine di migliaia è costoso. Una goroutine parte con uno
stack molto piccolo e cresciuto dinamicamente, ed è schedulata dal
\emph{runtime} in spazio utente, quindi se ne possono avere centinaia di migliaia.
In \textbf{Python} (CPython) i thread sono thread del sistema operativo, ma
il \emph{Global Interpreter Lock} impedisce l'esecuzione \emph{parallela} di
bytecode Python in più thread; le goroutine invece possono girare in parallelo
su più core. In Go non c'è un oggetto \cod{Thread} da creare e gestire: si
scrive \cod{go f()}.
\end{confronto}

\section{Channel}\label{sec:channel}

\subsection{Introduzione}

I \textbf{channel} permettono una comunicazione \emph{efficiente e sicura}
tra goroutine tramite l'invio di valori. Un channel è una \textbf{coda
thread-safe} gestita dal runtime di Go; nasconde la complessità della
comunicazione tra thread (le slide precisano che internamente usa
meccanismi di sincronizzazione come mutex e semafori). Più goroutine possono inviare
valori nello stesso channel: è ideale per notifiche, multiplexing, \ldots\ e la
comunicazione è thread-safe \emph{«out of the box»}.

\begin{definizione}{channel}
Un \textbf{channel} è un \emph{condotto tipizzato} attraverso il quale una
goroutine può inviare e ricevere valori con l'operatore \cod{<-}. I dati fluiscono
\emph{nella direzione della freccia}.
\end{definizione}

\subsection{Sintassi}

\begin{lstlisting}[style=gocode,numbers=none]
ch <- v        // invia v sul channel ch
v := <-ch      // riceve da ch e assegna a v
ch := make(chan int)   // crea un channel di int
\end{lstlisting}

\begin{itemize}
  \item Un channel trasporta valori di un tipo specifico (\cod{int}, struct,
        \ldots). Per impostazione predefinita è \textbf{bidirezionale}.
  \item Va creato con \cod{make} prima dell'uso. Il valore zero di un channel è
        \cod{nil}.
  \item \textbf{Invio e ricezione bloccano} finché l'altro lato non è pronto
        (per i channel senza buffer, vedi Sezione~\ref{sec:buffered}). Le
        goroutine possono quindi \textbf{sincronizzarsi} senza lock espliciti né
        variabili di condizione.
\end{itemize}

\subsection{Un channel come sincronizzazione}

Il programma seguente (esempio aggiuntivo) mostra l'uso più semplice di un
channel come meccanismo di sincronizzazione: \cod{main} attende la fine di una
goroutine ricevendo da un channel \cod{done}.

\begin{lstlisting}[style=gocode,caption={\texttt{06\_concurrency/channel\_sync.go}}]
// channel_sync.go - Un channel come meccanismo di sincronizzazione
// (slide 50).
//
// Esempio aggiuntivo: main attende la fine di una goroutine ricevendo da un
// channel "done". Non servono lock né condition variable.
package main

import (
	"fmt"
	"time"
)

func worker(done chan bool) {
	fmt.Println("worker: inizio lavoro...")
	time.Sleep(100 * time.Millisecond) // lavoro simulato
	fmt.Println("worker: fatto")
	done <- true // segnala il completamento
}

func main() {
	done := make(chan bool)
	go worker(done)

	// La ricezione blocca main finché il worker non invia.
	<-done
	fmt.Println("main: il worker ha terminato")
}
\end{lstlisting}
\begin{lstlisting}[style=output]
worker: inizio lavoro...
worker: fatto
main: il worker ha terminato
\end{lstlisting}

La ricezione \cod{<-done} blocca \cod{main} finché il worker non invia
\cod{true}; quindi «main: il worker ha terminato» compare sempre dopo
«worker: fatto».

\subsection{Esempio: somma distribuita}

La somma di una slice è distribuita tra due goroutine: è il pattern \textbf{SPMD}
(\emph{Single Program Multiple Data}), in cui ogni goroutine esegue lo stesso
programma su una porzione diversa dei dati.

\begin{lstlisting}[style=gocode,caption={\texttt{06\_concurrency/channel\_sum.go}}]
// channel_sum.go - Somma distribuita tra due goroutine (slide 51).
//
// Pattern SPMD (Single Program Multiple Data): ogni goroutine esegue lo
// stesso codice su una porzione diversa dei dati.
package main

import "fmt"

func sum(s []int, c chan int) {
	sum := 0
	for _, v := range s {
		sum += v
	}
	c <- sum // send sum to c
}

func main() {
	s := []int{7, 2, 8, -9, 4, 0}
	c := make(chan int)
	go sum(s[:len(s)/2], c) // prima metà: 7+2+8 = 17
	go sum(s[len(s)/2:], c) // seconda metà: -9+4+0 = -5

	// Le due ricezioni bloccano finché le goroutine non inviano.
	// Non si sa quale somma arrivi per prima, ma il totale è lo stesso
	// (l'addizione è commutativa).
	sum1, sum2 := <-c, <-c
	fmt.Println("Total sum:", sum1+sum2)
}
\end{lstlisting}
\begin{lstlisting}[style=output]
Total sum: 12
\end{lstlisting}

\subsubsection*{Analisi del codice e flusso di esecuzione}

\begin{enumerate}
  \item \cod{main} crea il channel \cod{c} (senza buffer) e lancia due goroutine
        \cod{sum}, sulla prima e sulla seconda metà di \cod{s}:
        $7+2+8=17$ e $-9+4+0=-5$.
  \item \cod{sum1, sum2 := <-c, <-c} esegue due ricezioni. Ciascuna \emph{blocca}
        \cod{main} finché una goroutine non invia il proprio risultato.
  \item Non si sa quale somma arrivi per prima, ma il totale
        $17+(-5)=12$ è lo stesso perché l'addizione è commutativa: l'output
        è deterministico anche se l'ordine di arrivo non lo è.
\end{enumerate}

\subsection{Esempio: generatore di Fibonacci}

\begin{lstlisting}[style=gocode,caption={\texttt{06\_concurrency/channel\_fib.go}}]
// channel_fib.go - Generatore di Fibonacci con goroutine e channel
// (slide 52).
package main

import "fmt"

// fib è un produttore: invia numeri di Fibonacci sul channel c, per
// sempre. Ogni invio blocca finché il main non riceve (channel senza
// buffer), quindi il produttore non va "avanti" rispetto al consumatore.
func fib(c chan int) {
	x, y := 1, 1
	for {
		c <- x
		x, y = y, x+y
	}
}

func main() {
	c := make(chan int)
	go fib(c)
	for i := 0; i < 10; i++ {
		fmt.Println(<-c)
	}
	// Quando main termina, la goroutine fib (bloccata su un invio)
	// viene semplicemente abbandonata: si veda select per una chiusura
	// ordinata.
}
\end{lstlisting}
\begin{lstlisting}[style=output]
1
1
2
3
5
8
13
21
34
55
\end{lstlisting}

\cod{fib} è un \emph{produttore} che invia per sempre sul channel; \cod{main} è il
\emph{consumatore} e riceve solo 10 valori. Poiché il channel non ha buffer, ogni
invio blocca finché \cod{main} non riceve: il produttore non può «correre
avanti». Dopo la decima ricezione, \cod{main} termina e la goroutine \cod{fib}
(bloccata su un invio) viene semplicemente abbandonata: è accettabile qui
perché il programma finisce, ma in un programma di lunga durata una goroutine
bloccata per sempre è una \emph{perdita} (\emph{goroutine leak}). La versione
con \cod{select} e un channel \cod{quit} (Sezione~\ref{sec:select}) la ferma in
modo ordinato.

\begin{confronto}{C, Python, Java}
Un channel somiglia a una coda bloccante: in \textbf{Java}
\cod{SynchronousQueue} (senza buffer) o \cod{ArrayBlockingQueue} (con buffer), in
\textbf{Python} \cod{queue.Queue}, in \textbf{C} un \emph{pipe} o una coda
protetta da mutex e condition variable scritta a mano. Le differenze
caratteristiche di Go sono che il channel è un tipo \emph{di primo livello} del
linguaggio, ha tipo statico, direzione e operazioni di \cod{close}, e si
combina con \cod{select}.
\end{confronto}

\section{Channel con e senza buffer}\label{sec:buffered}

\begin{description}
  \item[Channel senza buffer (\emph{unbuffered}).] Per impostazione
        predefinita le operazioni su un channel bloccano. Secondo la specifica del
        linguaggio: se la capacità è zero o assente, il channel è
        senza buffer e la comunicazione riesce \emph{solo quando mittente e
        ricevente sono entrambi pronti}; in un channel senza buffer il mittente
        resta bloccato finché il ricevente non ha ricevuto il valore.
  \item[Channel con buffer (\emph{buffered}).] Non bloccano se non sono
        pieni (per l'invio) o non sono vuoti (per la ricezione). Si specifica la
        capacità alla creazione: \cod{ch := make(chan int, 100)}. L'invio
        blocca \emph{solo} se il buffer è pieno; la ricezione blocca \emph{solo}
        se il buffer è vuoto. È utile per \emph{disaccoppiare} le goroutine:
        possono procedere senza sincronizzarsi a ogni passo.
\end{description}

\begin{lstlisting}[style=gocode,caption={\texttt{06\_concurrency/channel\_buffered.go}}]
// channel_buffered.go - Channel con e senza buffer (slide 53).
package main

import "fmt"

func main() {
	// Con buffer di capacità 2: i primi due invii NON bloccano anche se
	// nessuno sta ancora ricevendo.
	ch := make(chan int, 2)
	ch <- 1
	ch <- 2
	fmt.Println("len =", len(ch), "cap =", cap(ch)) // len=2 cap=2

	// Un terzo invio bloccherebbe (buffer pieno) e, non esistendo altre
	// goroutine che ricevano, causerebbe un deadlock. Le ricezioni
	// seguono l'ordine FIFO.
	fmt.Println(<-ch) // 1
	fmt.Println(<-ch) // 2

	// Un channel senza buffer (capacità 0) richiede sender e receiver
	// pronti contemporaneamente: l'invio avviene in una goroutine perché
	// altrimenti main si bloccherebbe da solo.
	unbuf := make(chan string)
	go func() { unbuf <- "ping" }()
	fmt.Println(<-unbuf)
}
\end{lstlisting}
\begin{lstlisting}[style=output]
len = 2 cap = 2
1
2
ping
\end{lstlisting}

\subsection{Analisi del codice}

\begin{itemize}
  \item Con \cod{make(chan int, 2)} i primi due invii non bloccano anche se
        nessuno riceve; \cod{len} e \cod{cap} dicono quanti elementi ci sono e
        quanti ne entrano. Le ricezioni seguono l'ordine \textbf{FIFO}.
  \item Un \emph{terzo} invio senza ricezioni nel frattempo bloccherebbe per
        sempre: nel \cod{main} da solo, il runtime lo rileverebbe come deadlock
        (Sezione~\ref{sec:deadlock}).
  \item Nel channel senza buffer, l'invio è fatto in una goroutine apposita:
        eseguito direttamente da \cod{main} bloccherebbe \cod{main} stesso, prima
        che arrivi la ricezione sulla riga successiva.
\end{itemize}

\begin{nota}[Una precisazione sulle slide]
Una slide afferma che «invio e ricezione bloccano finché l'altro lato è pronto»;
questa descrizione è esatta per i channel \emph{senza buffer}. Con un buffer
(slide successive) i blocchi si verificano solo a buffer pieno (invio) o vuoto
(ricezione). Le due affermazioni sono coerenti se lette in quest'ordine.
\end{nota}

\section{Chiudere un channel e iterare con \texttt{range}}\label{sec:close}

\subsection{\texttt{close}}

Un channel si chiude con \cod{close(ch)}. Le regole sono:

\begin{itemize}
  \item il/i \textbf{mittente/i} devono smettere di inviare dopo la chiusura;
  \item il/i \textbf{ricevente/i} possono ancora ricevere i \emph{valori già nel
        buffer};
  \item i riceventi possono verificare se il channel è chiuso con la forma a due
        valori \cod{v, ok := <-ch}: \cod{ok} è \cod{false} se il channel è chiuso
        \emph{e} non ci sono più valori da ricevere;
  \item \textbf{solo il mittente} (non il ricevente) dovrebbe chiudere il
        channel;
  \item \textbf{inviare su un channel chiuso} causa un \textbf{panic}:
        \texttt{panic: send on closed channel};
  \item ricevere da un channel chiuso e vuoto restituisce subito il \emph{valore
        zero} (con \cod{ok == false}).
\end{itemize}

\begin{lstlisting}[style=gocode,caption={\texttt{06\_concurrency/channel\_close.go}}]
// channel_close.go - close e rilevamento della chiusura (slide 54).
package main

import "fmt"

func main() {
	jobs := make(chan int, 5)
	done := make(chan bool)

	// Il consumatore legge finché il channel è chiuso E vuoto.
	go func() {
		for {
			j, more := <-jobs // more == false: chiuso e senza altri valori
			if more {
				fmt.Println("received job", j)
			} else {
				fmt.Println("received all jobs")
				done <- true
				return
			}
		}
	}()

	// Il produttore (il sender) invia e poi chiude: solo il sender
	// dovrebbe chiudere un channel.
	for j := 1; j <= 3; j++ {
		jobs <- j
		fmt.Println("sent job", j)
	}
	close(jobs)
	fmt.Println("sent all jobs")

	<-done // attende il consumatore

	// Ricevere da un channel chiuso e vuoto restituisce subito il valore
	// zero e ok == false.
	v, ok := <-jobs
	fmt.Println(v, ok) // 0 false
}
\end{lstlisting}
\begin{lstlisting}[style=output]
sent job 1
sent job 2
sent job 3
sent all jobs
received job 1
received job 2
received job 3
received all jobs
0 false
\end{lstlisting}

\subsubsection*{Analisi del codice}

Il consumatore (goroutine anonima) legge in un ciclo con
\cod{j, more := <-jobs}: finché \cod{more} è vero stampa il job, appena diventa
falso (channel chiuso \emph{e} vuoto) segnala la fine su \cod{done}. Il
produttore (\cod{main}) invia tre job, chiude il channel, e attende \cod{done}. In
fondo la ricezione da un channel già chiuso e vuoto dà \cod{0 false}.

\begin{nota}
Nell'output l'\emph{intercalarsi} delle righe «sent» e «received» non è
garantito: dipende dallo scheduler. Qui il buffer di 5 posti permette a \cod{main}
di inviare tutti e tre i job prima che il consumatore parta, ma in
un'altra esecuzione le righe potrebbero alternarsi. Ciò che è garantito è
l'ordine interno (job 1, 2, 3) e che «received all jobs» segua tutti i
«received job».
\end{nota}

\begin{attenzione}[Errore comune: inviare su un channel chiuso]
L'invio su un channel chiuso non restituisce un errore ma provoca un
\emph{panic}: si arresta l'intero programma. Per questo la chiusura spetta al
mittente, che sa quando non ha più nulla da inviare; un ricevente non ha modo
di sapere se altri mittenti stanno ancora inviando. L'esempio seguente termina
\emph{volutamente} con un panic.

\begin{lstlisting}[style=gocode,caption={\texttt{06\_concurrency/pitfall\_send\_closed.go}}]
// pitfall_send_closed.go - ERRORE COMUNE: invio su un channel chiuso
// (slide 54).
//
// ATTENZIONE: questo programma TERMINA VOLUTAMENTE con un panic:
//
//	panic: send on closed channel
//
// Si noti che la RICEZIONE da un channel chiuso è invece lecita.
package main

import "fmt"

func main() {
	ch := make(chan int, 2)
	ch <- 1
	close(ch)

	v, ok := <-ch // ricezione dal buffer residuo: 1 true
	fmt.Println(v, ok)
	v, ok = <-ch // channel chiuso e vuoto: 0 false
	fmt.Println(v, ok)

	fmt.Println("invio su channel chiuso...")
	ch <- 2 // PANIC
}
\end{lstlisting}
\begin{lstlisting}[style=output]
1 true
0 false
invio su channel chiuso...
panic: send on closed channel

goroutine 1 [running]:
main.main()
        .../pitfall_send_closed.go:24 +0x1a5
exit status 2
\end{lstlisting}
La ricezione dopo \cod{close} è lecita (prima il valore nel buffer, \cod{1 true},
poi il valore zero, \cod{0 false}); l'invio no.
\end{attenzione}

\subsection{\texttt{range} su un channel}\label{sec:chanrange}

Con \cod{range} si ricevono ripetutamente valori da un channel \emph{finché
non viene chiuso}:

\begin{lstlisting}[style=gocode,numbers=none]
for elem := range ch {
	fmt.Println(elem)
}
\end{lstlisting}

\begin{lstlisting}[style=gocode,caption={\texttt{06\_concurrency/channel\_range.go}}]
// channel_range.go - range su un channel (slide 55).
package main

import "fmt"

func main() {
	ch := make(chan int, 3)

	go func() {
		for i := 1; i <= 5; i++ {
			ch <- i
		}
		// Senza close, il range sotto non terminerebbe mai e main
		// andrebbe in deadlock dopo aver letto i cinque valori.
		close(ch)
	}()

	// Riceve ripetutamente finché il channel è chiuso e svuotato.
	for elem := range ch {
		fmt.Println(elem)
	}
	fmt.Println("channel chiuso, ciclo terminato")
}
\end{lstlisting}
\begin{lstlisting}[style=output]
1
2
3
4
5
channel chiuso, ciclo terminato
\end{lstlisting}

\begin{attenzione}[Errore comune: dimenticare \texttt{close} con \texttt{range}]
Il ciclo \cod{for elem := range ch} termina \emph{solo} quando \cod{ch} è chiuso
e svuotato. Se il mittente dimentica \cod{close(ch)}, dopo l'ultimo valore il
ciclo si blocca in attesa di altri: se non esistono altre goroutine in grado di
inviare, il runtime segnala un deadlock.
\end{attenzione}

\section{\texttt{select}}\label{sec:select}

\subsection{Introduzione e sintassi}

\cod{select} permette a una goroutine di \textbf{attendere su più channel
contemporaneamente}. Blocca finché uno dei suoi \cod{case} può procedere, poi
esegue quel caso.

\begin{itemize}
  \item Se \emph{più} case sono pronti, ne viene scelto \textbf{uno a caso}.
  \item Se nessun case è pronto e c'è un ramo \cod{default}, viene eseguito
        quello (comportamento non bloccante).
\end{itemize}

\begin{lstlisting}[style=gocode,numbers=none]
select {
case msg1 := <-ch1:    // ricezione su ch1
	// gestisci msg1
case msg2 := <-ch2:    // ricezione su ch2
	// gestisci msg2
}
\end{lstlisting}

È utile per il \emph{multiplexing} e per gestire i \emph{timeout}.

\begin{nota}[Un refuso nelle slide]
Nel frammento della slide la variabile della prima ricezione è scritta
\cod{mgs1} (e poi commentata come \cod{msg1}). È un errore di battitura:
il nome corretto, usato nei listati di questa guida, è \cod{msg1}.
\end{nota}

\subsection{Esempio: Fibonacci con \texttt{select} e un channel \texttt{quit}}

Versione del generatore con due channel: \cod{ch} per inviare i valori e
\cod{quit} per segnalare quando fermarsi.

\begin{lstlisting}[style=gocode,caption={\texttt{06\_concurrency/channel\_fib\_select.go}}]
// channel_fib_select.go - Fibonacci con select e channel quit (slide 57-58).
package main

import "fmt"

func fibonacci(ch, quit chan int) {
	x, y := 1, 1
	for {
		select {
		case ch <- x: // invia il numero di Fibonacci corrente
			x, y = y, x+y
		case <-quit: // riceve il segnale di terminazione
			fmt.Println("quit")
			return
		}
	}
}

func main() {
	ch := make(chan int) // channel senza buffer
	quit := make(chan int)

	go func() { // goroutine anonima: il consumatore
		for i := 0; i < 10; i++ {
			fmt.Println(<-ch) // riceve un valore di Fibonacci
		}
		quit <- 0 // dice a fibonacci di fermarsi
	}()

	fibonacci(ch, quit) // eseguita nella goroutine principale
}
\end{lstlisting}
\begin{lstlisting}[style=output]
1
1
2
3
5
8
13
21
34
55
quit
\end{lstlisting}

\subsubsection*{Analisi del codice e flusso di esecuzione}

\begin{enumerate}
  \item \cod{main} avvia una goroutine anonima (il consumatore) e poi esegue
        \cod{fibonacci(ch, quit)} \emph{nella goroutine principale}.
  \item Dentro \cod{fibonacci}, il \cod{select} offre due possibilità: inviare
        \cod{x} su \cod{ch} (pronto solo quando il consumatore sta ricevendo) o
        ricevere da \cod{quit} (pronto solo quando qualcuno ha inviato).
  \item Il consumatore riceve 10 valori, poi invia su \cod{quit}. Da quel
        momento il ramo \cod{quit} del \cod{select} diventa pronto
        mentre il ramo di invio non lo è più (nessuno riceve): \cod{fibonacci}
        stampa «quit» e ritorna, terminando in modo \emph{pulito}.
\end{enumerate}

\begin{attenzione}[Errore nella slide: variabile non definita]
Nelle slide \cod{main} dichiara \cod{ch := make(chan int)} ma poi chiama
\cod{fibonacci(c, quit)}: la variabile \cod{c} \emph{non esiste} e il
programma non compila (errore reale: \texttt{undefined: c}). Si tratta di un
\textbf{errore di trascrizione} nella slide, non una scelta tecnica: il
listato di questa guida usa il nome corretto \cod{ch}. Si veda anche
l'Appendice~B.
\end{attenzione}

\subsection{Scelta casuale tra case pronti}

Se entrambi i case sono pronti, la scelta è casuale. Il programma seguente ha
due channel già pieni:

\begin{lstlisting}[style=gocode,caption={\texttt{06\_concurrency/select\_multiplex.go}}]
// select_multiplex.go - select su più channel e scelta casuale (slide 56).
package main

import "fmt"

func main() {
	ch1 := make(chan string, 1)
	ch2 := make(chan string, 1)
	ch1 <- "da ch1"
	ch2 <- "da ch2"

	// Entrambi i case sono pronti: select ne sceglie UNO a caso. L'ordine
	// delle due righe stampate può quindi variare da esecuzione a
	// esecuzione, ma entrambi i messaggi vengono sempre ricevuti.
	for i := 0; i < 2; i++ {
		select {
		case msg1 := <-ch1:
			fmt.Println("ricevuto", msg1)
		case msg2 := <-ch2:
			fmt.Println("ricevuto", msg2)
		}
	}
}
\end{lstlisting}

L'output è \textbf{non deterministico}: l'ordine delle due righe varia. In sei
esecuzioni consecutive si sono osservate entrambe le sequenze:

\begin{lstlisting}[style=output]
ricevuto da ch1
ricevuto da ch2

ricevuto da ch2
ricevuto da ch1
\end{lstlisting}

La scelta casuale evita che un channel «affami» gli altri.

\subsection{\texttt{select} non bloccante con \texttt{default}}

Con un ramo \cod{default} il \cod{select} \emph{non blocca} se i channel non sono
pronti. I case possono essere anche invii:

\begin{lstlisting}[style=gocode,caption={\texttt{06\_concurrency/select\_default.go}}]
// select_default.go - select non bloccante con default (slide 56, 59).
package main

import "fmt"

func main() {
	ch1 := make(chan string)
	ch3 := make(chan string, 1)

	// Ricezione non bloccante: nessun valore pronto -> scatta default.
	select {
	case msg := <-ch1:
		fmt.Println("ricevuto", msg)
	default:
		fmt.Println("nessun messaggio pronto")
	}

	// Invio non bloccante: ch3 ha un posto libero nel buffer.
	select {
	case ch3 <- "ciao":
		fmt.Println("inviato su ch3")
	default:
		fmt.Println("ch3 pieno")
	}

	// Secondo invio: il buffer è ora pieno, quindi si va in default.
	select {
	case ch3 <- "altro":
		fmt.Println("inviato su ch3")
	default:
		fmt.Println("ch3 pieno")
	}

	// Ora c'è un valore pronto su ch3: viene scelto il case di ricezione.
	select {
	case msg := <-ch3:
		fmt.Println("ricevuto", msg)
	default:
		fmt.Println("nessun messaggio pronto")
	}
}
\end{lstlisting}
\begin{lstlisting}[style=output]
nessun messaggio pronto
inviato su ch3
ch3 pieno
ricevuto ciao
\end{lstlisting}

\subsubsection*{Analisi del codice}

Nel primo \cod{select} nessuno ha inviato su \cod{ch1}: scatta \cod{default}.
\cod{ch3} ha buffer 1: il primo invio riesce (spazio libero), il secondo trova il
buffer pieno e va in \cod{default}; infine la ricezione trova un valore
pronto. Senza \cod{default}, ognuna di queste operazioni bloccherebbe.

\section{Timer}\label{sec:timer}

Un \textbf{timer} fornisce un channel (il campo \cod{C}) su cui riceve un valore
quando la durata specificata è trascorsa. Si usa per attendere un tempo
preciso:

\begin{lstlisting}[style=gocode,numbers=none]
timer := time.NewTimer(time.Second * 2)   // attende 2 secondi
<-timer.C                                  // blocca finché il timer scatta
\end{lstlisting}

Un timer può essere \textbf{annullato} prima che scatti con \cod{Stop()}. È
utile per implementare:

\begin{itemize}
  \item \textbf{timeout}: attendere un evento o il completamento di
        un'operazione prima di procedere;
  \item \textbf{ritardi}: attendere un periodo prima di eseguire un compito.
\end{itemize}

\begin{lstlisting}[style=gocode,caption={\texttt{06\_concurrency/timers.go}}]
// timers.go - Timer e timeout con select (slide 56, 60).
package main

import (
	"fmt"
	"time"
)

func main() {
	// Attesa di una durata: il timer espone un channel C su cui riceve
	// un valore quando scade.
	start := time.Now()
	timer := time.NewTimer(200 * time.Millisecond)
	<-timer.C // blocca fino allo scadere
	fmt.Println("timer 1 scaduto, trascorsi almeno 200ms:",
		time.Since(start) >= 200*time.Millisecond)

	// Un timer si può annullare con Stop() prima che scatti.
	timer2 := time.NewTimer(time.Second)
	stopped := timer2.Stop()
	fmt.Println("timer 2 annullato:", stopped)

	// Timeout: select attende un risultato oppure il timer, a seconda di
	// quale arrivi prima.
	result := make(chan string)
	go func() {
		time.Sleep(500 * time.Millisecond) // operazione lenta
		result <- "risultato"
	}()

	// Anche il timeout è un timer: il suo channel C si "attiva" allo
	// scadere dei 100ms. Vince il case che diventa pronto per primo.
	timeout := time.NewTimer(100 * time.Millisecond)
	select {
	case r := <-result:
		fmt.Println(r)
		timeout.Stop()
	case <-timeout.C:
		fmt.Println("timeout: l'operazione è troppo lenta")
	}
}
\end{lstlisting}
\begin{lstlisting}[style=output]
timer 1 scaduto, trascorsi almeno 200ms: true
timer 2 annullato: true
timeout: l'operazione è troppo lenta
\end{lstlisting}

\subsubsection*{Analisi del codice}

Il primo timer blocca \cod{main} per 200\,ms; il programma verifica solo che sia
trascorso \emph{almeno} quel tempo (la durata esatta non è deterministica, ma non può
essere inferiore). \cod{Stop()} annulla il secondo timer: restituisce \cod{true}
perché il timer non era ancora scattato. L'ultimo blocco realizza un
\emph{timeout} con \cod{select}: l'operazione lenta impiega 500\,ms, il timer da
100\,ms scatta per primo e vince il relativo \cod{case}. Quella goroutine resterà
in attesa di inviare su un channel senza ricevente, ma qui il programma termina
subito.

\section{Mutua esclusione}\label{sec:mutex}

\subsection{Esercizio: un mutex realizzato con un channel}

Go fornisce i \emph{mutex} per garantire l'accesso sicuro ai dati condivisi tra
goroutine; \cod{sync.Mutex} è lo strumento usuale per bloccare e sbloccare.
Le slide propongono anche un esercizio: realizzare un mutex \emph{usando un
channel}, completando lo scheletro

\begin{lstlisting}[style=gocode,numbers=none]
type Lock struct { /* ? */ }
func NewLock() Lock       { /* ? */ }
func (l *Lock) Lock()     { /* ? */ }
func (l *Lock) Unlock()   { /* ? */ }
\end{lstlisting}

\textbf{Idea della soluzione.} Un channel con buffer di capacità 1 contiene al
più un «gettone». Chi possiede il gettone detiene il lock; chi non lo trova
si blocca sulla ricezione:

\begin{itemize}
  \item \cod{NewLock} crea il channel con buffer 1 e vi inserisce un valore:
        il lock parte \emph{libero} (gettone presente);
  \item \cod{Lock()} \emph{riceve} dal channel: se il gettone c'è lo prende (il
        lock è ora occupato), altrimenti la goroutine si blocca finché qualcuno
        lo restituisce;
  \item \cod{Unlock()} \emph{invia} sul channel: restituisce il gettone,
        liberando il lock e svegliando eventualmente una goroutine in
        attesa.
\end{itemize}

L'invariante è: gettoni nel channel $=1$ se il lock è libero, $0$ se occupato.

\begin{lstlisting}[style=gocode,caption={\texttt{06\_concurrency/mutex\_using\_channel.go}}]
// mutex_using_channel.go - Esercizio: un mutex realizzato con un channel
// (slide 61-62).
//
// Idea: un channel con buffer di capacità 1 contiene al più un "gettone".
// Chi possiede il gettone detiene il lock; chi lo trova assente si blocca
// sulla ricezione.
package main

import (
	"fmt"
	"sync"
)

// Nota: "go vet" segnala "return copies lock value" su NewLock perché il
// tipo ha metodi Lock/Unlock e vet applica l'euristica dei lock non
// copiabili (come sync.Mutex). Qui la copia è innocua: la struct contiene
// solo un riferimento al channel, quindi tutte le copie condividono lo
// stesso "gettone". Il codice è quello della slide.
type Lock struct {
	ch chan bool
}

func NewLock() Lock {
	lock := Lock{make(chan bool, 1)} // channel con buffer
	lock.ch <- true                  // gettone iniziale: il lock parte libero
	return lock
}

// Lock prende il gettone; se non c'è, blocca finché qualcuno lo restituisce.
func (l *Lock) Lock() {
	<-l.ch // receive
}

// Unlock restituisce il gettone.
func (l *Lock) Unlock() {
	l.ch <- true // send
}

func main() {
	lock := NewLock()
	counter := 0

	var wg sync.WaitGroup
	for i := 0; i < 100; i++ {
		wg.Add(1)
		go func() {
			defer wg.Done()
			for j := 0; j < 1000; j++ {
				lock.Lock()
				counter++ // sezione critica
				lock.Unlock()
			}
		}()
	}
	wg.Wait()
	fmt.Println("counter =", counter) // sempre 100000
}
\end{lstlisting}
\begin{lstlisting}[style=output]
counter = 100000
\end{lstlisting}

Cento goroutine incrementano mille volte la stessa variabile \cod{counter}
proteggendola con il lock: il risultato è sempre $100\times1000=100000$.

\begin{nota}[\texttt{go vet} e la copia del lock]
Il codice è quello della slide. L'analizzatore \cod{go vet} segnala
\texttt{return copies lock value} su \cod{NewLock}, perché il tipo ha metodi
\cod{Lock}/\cod{Unlock} e vet applica l'euristica dei «lock non copiabili»
(come \cod{sync.Mutex}, che \emph{non} va copiato). Nel nostro caso la copia
è innocua: la struct contiene solo il \emph{riferimento} a un channel, e tutte
le copie condividono lo stesso gettone. Le slide stesse notano che, nella
pratica, \cod{sync.Mutex} è più idiomatico ed efficiente.
\end{nota}

\subsection{\texttt{sync.Mutex}}

\textbf{Mutex} (\cod{sync.Mutex}): garantisce la \emph{mutua esclusione} e
previene le \emph{race condition}; solo una goroutine alla volta può accedere
alla \emph{sezione critica}. L'uso è:

\begin{lstlisting}[style=gocode,numbers=none]
mu.Lock()
counter++
mu.Unlock()
\end{lstlisting}

\begin{lstlisting}[style=gocode,caption={\texttt{06\_concurrency/mutex.go}}]
// mutex.go - sync.Mutex per la mutua esclusione (slide 63).
package main

import (
	"fmt"
	"sync"
)

// Counter protegge il campo n con un mutex.
type Counter struct {
	mu sync.Mutex
	n  int
}

func (c *Counter) Inc() {
	c.mu.Lock()
	// defer garantisce lo Unlock anche in caso di return anticipato o panic.
	defer c.mu.Unlock()
	c.n++
}

func main() {
	var c Counter
	var wg sync.WaitGroup

	for i := 0; i < 100; i++ {
		wg.Add(1)
		go func() {
			defer wg.Done()
			for j := 0; j < 1000; j++ {
				c.Inc()
			}
		}()
	}
	wg.Wait()
	fmt.Println("counter =", c.n) // sempre 100000
}
\end{lstlisting}
\begin{lstlisting}[style=output]
counter = 100000
\end{lstlisting}

Si noti \cod{defer c.mu.Unlock()} subito dopo il \cod{Lock}: garantisce il
rilascio anche in caso di \cod{return} anticipato o di panic, ed è l'applicazione
tipica di \cod{defer} (Sezione~\ref{sec:defer}).

\begin{attenzione}[Errore comune: mancata sincronizzazione (data race)]
Se si omette il mutex, l'operazione \cod{counter++} \emph{non è atomica}: si
compone di lettura, incremento e scrittura. Due goroutine possono leggere lo
stesso valore, incrementarlo e riscrivere \emph{lo stesso} risultato,
perdendo un aggiornamento. L'esempio seguente ha un output \textbf{non
deterministico} (di solito inferiore a 100000).

\begin{lstlisting}[style=gocode,caption={\texttt{06\_concurrency/race\_counter.go}}]
// race_counter.go - ERRORE COMUNE: data race su una variabile condivisa
// (slide 47, 63).
//
// counter++ NON è atomico (legge, incrementa, scrive). Con più goroutine
// non sincronizzate alcuni aggiornamenti si perdono. Il risultato è
// NON deterministico e quasi sempre inferiore a 100000. (In rari casi su
// macchine con un solo thread potrebbe risultare esatto.)
package main

import (
	"fmt"
	"sync"
)

func main() {
	counter := 0
	var wg sync.WaitGroup

	for i := 0; i < 100; i++ {
		wg.Add(1)
		go func() {
			defer wg.Done()
			for j := 0; j < 1000; j++ {
				counter++ // DATA RACE: nessuna mutua esclusione
			}
		}()
	}
	wg.Wait()
	fmt.Println("atteso 100000, ottenuto", counter)
	// Correzione: vedi mutex.go (sync.Mutex).
}
\end{lstlisting}
\begin{lstlisting}[style=output]
atteso 100000, ottenuto 94808
atteso 100000, ottenuto 99120
atteso 100000, ottenuto 95983
\end{lstlisting}
(Tre esecuzioni diverse: i valori cambiano a ogni lancio.) Il bug è insidioso perché il programma compila,
termina e produce un numero plausibile: l'errore è nella \emph{perdita}
silenziosa di aggiornamenti. Si risolve con \cod{sync.Mutex} (o con un
channel come sopra).
\end{attenzione}

\begin{nota}[Un mutex non è rientrante]
Se una goroutine chiama \cod{Lock()} mentre già detiene lo stesso mutex, si
blocca per sempre su se stessa (deadlock): a differenza dei
\cod{synchronized}/\cod{ReentrantLock} di Java, \cod{sync.Mutex} non è
rientrante.
\end{nota}

\subsection{\texttt{sync.RWMutex}}

Il \textbf{mutex lettori/scrittore} (\cod{sync.RWMutex}) consente \emph{più
lettori oppure un solo scrittore}. Migliora le prestazioni quando le letture
sono frequenti e le scritture rare.

\begin{itemize}
  \item \textbf{Lettori}: \cod{RLock()}\,/\,\cod{RUnlock()} $\to$ accesso
        \emph{condiviso};
  \item \textbf{Scrittore}: \cod{Lock()}\,/\,\cod{Unlock()} $\to$ accesso
        \emph{esclusivo}.
\end{itemize}

\begin{lstlisting}[style=gocode,caption={\texttt{06\_concurrency/rwmutex.go}}]
// rwmutex.go - sync.RWMutex: molti lettori oppure un solo scrittore
// (slide 63).
package main

import (
	"fmt"
	"sync"
)

// Store è una cache chiave-valore letta spesso e scritta di rado.
type Store struct {
	mu   sync.RWMutex
	data map[string]int
}

// Get usa un lock di LETTURA: più Get possono procedere insieme.
func (s *Store) Get(k string) (int, bool) {
	s.mu.RLock()
	defer s.mu.RUnlock()
	v, ok := s.data[k]
	return v, ok
}

// Set usa il lock di SCRITTURA: accesso esclusivo (blocca lettori e altri
// scrittori).
func (s *Store) Set(k string, v int) {
	s.mu.Lock()
	defer s.mu.Unlock()
	s.data[k] = v
}

func main() {
	s := &Store{data: make(map[string]int)}
	s.Set("a", 1)

	var wg sync.WaitGroup
	// 10 lettori concorrenti.
	for i := 0; i < 10; i++ {
		wg.Add(1)
		go func() {
			defer wg.Done()
			if v, ok := s.Get("a"); !ok || v < 1 {
				fmt.Println("valore inatteso")
			}
		}()
	}
	// 1 scrittore concorrente.
	wg.Add(1)
	go func() {
		defer wg.Done()
		s.Set("b", 2)
	}()
	wg.Wait()

	v, _ := s.Get("b")
	fmt.Println("b =", v)
}
\end{lstlisting}
\begin{lstlisting}[style=output]
b = 2
\end{lstlisting}

\cod{Get} usa il lock di lettura: dieci lettori possono procedere insieme.
\cod{Set} prende il lock di scrittura, che esclude lettori e altri scrittori. Il
risultato finale (\cod{b = 2}) è deterministico perché \cod{wg.Wait()} garantisce
che lo scrittore abbia terminato.

\begin{confronto}{C, Python, Java}
Equivalenti: in \textbf{C} \cod{pthread\_mutex\_t} e
\cod{pthread\_rwlock\_t}; in \textbf{Java} \cod{synchronized}, \cod{ReentrantLock},
\cod{ReentrantReadWriteLock}; in \textbf{Python} \cod{threading.Lock} e
\cod{RLock}. La differenza di filosofia: in Go si preferisce, quando
possibile, comunicare con channel anziché proteggere memoria condivisa con lock;
i mutex restano lo strumento giusto per proteggere uno stato semplice (un
contatore, una cache).
\end{confronto}

\section{Sincronizzazione con \texttt{WaitGroup}}\label{sec:waitgroup}

\cod{sync.WaitGroup} permette a un gruppo di goroutine di \emph{attendere
collettivamente} un evento prima di proseguire in modo indipendente. È come un
\textbf{contatore sicuro per la concorrenza} con tre operazioni: \cod{Add},
\cod{Done} e \cod{Wait}.

Due usi tipici, indicati dalle slide:

\begin{enumerate}
  \item \textbf{Pulizia}: assicurarsi che \emph{tutte} le goroutine (\cod{main}
        compresa) terminino prima della fine del programma;
  \item \textbf{Pattern BSP} (\emph{Bulk Synchronous Parallel}): le goroutine
        lavorano in modo indipendente per un po', poi attendono a una
        \emph{barriera} prima di procedere; lo scambio di dati può avvenire
        alla barriera.
\end{enumerate}

\subsection{Uso per la pulizia}

\begin{lstlisting}[style=gocode,caption={\texttt{06\_concurrency/waitgroup.go}}]
// waitgroup.go - sync.WaitGroup per attendere un gruppo di goroutine
// (slide 64).
package main

import (
	"fmt"
	"sync"
)

func worker(id int, results []int, wg *sync.WaitGroup) {
	defer wg.Done()
	// Ogni goroutine scrive in una posizione DIVERSA della slice: non c'è
	// condivisione dello stesso elemento, quindi non serve un lock.
	results[id] = id * id
}

func main() {
	results := make([]int, 5)
	var wg sync.WaitGroup

	for i := range results {
		wg.Add(1) // un'unità in più da attendere
		go worker(i, results, &wg)
	}
	wg.Wait() // attende che tutti abbiano chiamato Done

	// Dopo Wait tutte le scritture sono visibili a main.
	fmt.Println(results) // [0 1 4 9 16]
}
\end{lstlisting}
\begin{lstlisting}[style=output]
[0 1 4 9 16]
\end{lstlisting}

Ogni goroutine scrive nella \emph{propria} posizione di \cod{results}: non c'è
condivisione dello stesso elemento, quindi non servono lock. Dopo
\cod{wg.Wait()} le scritture di tutte le goroutine sono visibili a \cod{main}:
\cod{Wait} stabilisce la sincronizzazione necessaria.

\subsection{Uso come barriera: il pattern BSP}

\begin{lstlisting}[style=gocode,caption={\texttt{06\_concurrency/bsp\_barrier.go}}]
// bsp_barrier.go - Bulk Synchronous Parallel con WaitGroup come barriera
// (slide 64).
//
// Pattern BSP: le goroutine lavorano indipendentemente (fase di calcolo),
// poi attendono tutte a una barriera; solo dopo si passa alla fase
// successiva, in cui i dati prodotti sono visibili a tutti.
package main

import (
	"fmt"
	"sync"
)

func main() {
	const n = 4
	partial := make([]int, n)

	// Fase 1: ogni goroutine calcola un risultato parziale.
	var wg sync.WaitGroup
	for i := 0; i < n; i++ {
		wg.Add(1)
		go func() {
			defer wg.Done()
			partial[i] = (i + 1) * 10
		}()
	}
	wg.Wait() // BARRIERA

	// Scambio dati alla barriera: ogni goroutine della fase 2 legge
	// i risultati parziali di tutte le altre.
	total := 0
	for _, p := range partial {
		total += p
	}
	fmt.Println("totale dopo la fase 1:", total) // 100

	// Fase 2: ogni goroutine usa il totale globale.
	shares := make([]float64, n)
	for i := 0; i < n; i++ {
		wg.Add(1)
		go func() {
			defer wg.Done()
			shares[i] = float64(partial[i]) / float64(total)
		}()
	}
	wg.Wait() // seconda barriera
	fmt.Println("quote:", shares)
}
\end{lstlisting}
\begin{lstlisting}[style=output]
totale dopo la fase 1: 100
quote: [0.1 0.2 0.3 0.4]
\end{lstlisting}

Nella fase~1 ogni goroutine calcola un risultato parziale; la prima
\cod{wg.Wait()} è la \textbf{barriera}: nessuno passa alla fase~2 prima che tutti
abbiano finito. Alla barriera i dati sono scambiati: \cod{main} somma i parziali
($10+20+30+40=100$). Nella fase~2 ogni goroutine usa il totale globale per
calcolare la propria quota. Il \cod{WaitGroup} è riutilizzabile: il contatore
torna a zero e si può ricominciare con nuove chiamate a \cod{Add}.

\begin{confronto}{C, Python, Java}
\cod{WaitGroup} corrisponde a \cod{pthread\_join} su un insieme di thread, a
\cod{Thread.join()} (Python e Java) o, più da vicino, al
\cod{CountDownLatch} di Java. Per la barriera, \cod{CyclicBarrier} (Java) o
\cod{pthread\_barrier\_t} (C).
\end{confronto}

\section{Priorità e scheduling}\label{sec:priority}

Le goroutine \textbf{non hanno una schedulazione a priorità}: lo scheduler di Go
punta su equità e semplicità. Per imporre una priorità occorre realizzarla
\emph{a livello applicativo}, con una complessità aggiuntiva. Due approcci per
simularla, indicati dalle slide:

\begin{enumerate}
  \item usare \textbf{più channel} per classi di priorità diverse (alta,
        bassa);
  \item usare una \textbf{coda con priorità} (\emph{heap}) per ordinare i task
        per priorità.
\end{enumerate}

L'esempio seguente realizza il primo approccio.

\begin{lstlisting}[style=gocode,caption={\texttt{06\_concurrency/priority.go}}]
// priority.go - Simulare priorità con più channel (slide 65).
//
// Lo scheduler di Go non ha priorità: la si implementa a livello di
// applicazione. Qui due channel (alta/bassa priorità); il consumatore
// controlla SEMPRE prima il channel ad alta priorità con un select non
// bloccante e passa a quello a bassa priorità solo se il primo è vuoto.
package main

import "fmt"

func main() {
	high := make(chan string, 3)
	low := make(chan string, 3)

	// I task a bassa priorità vengono inseriti per primi, di proposito.
	low <- "low-1"
	low <- "low-2"
	high <- "high-1"
	high <- "high-2"

	const total = 4
	for done := 0; done < total; done++ {
		select {
		case t := <-high: // controllato per primo
			fmt.Println("eseguo", t)
		default:
			// Nessun task ad alta priorità: si serve il basso.
			select {
			case t := <-low:
				fmt.Println("eseguo", t)
			default:
			}
		}
	}
}
\end{lstlisting}
\begin{lstlisting}[style=output]
eseguo high-1
eseguo high-2
eseguo low-1
eseguo low-2
\end{lstlisting}

\subsubsection*{Analisi del codice}

I task a bassa priorità sono inseriti per primi, di proposito. Il consumatore, a
ogni giro, esegue prima un \cod{select} \emph{non bloccante} sul channel ad alta
priorità; solo se questo è vuoto (ramo \cod{default}) tenta il channel a bassa
priorità con un secondo \cod{select}. Il risultato è che tutti i task \cod{high}
vengono serviti prima di qualunque \cod{low}, indipendentemente dall'ordine di
inserimento. Si noti che l'ordine \emph{non è garantito dallo scheduler} ma
dalla struttura del codice applicativo, come ci si aspetta.

\section{Deadlock}\label{sec:deadlock}

Un \textbf{deadlock} si verifica quando tutte le goroutine sono bloccate in attesa
di eventi che nessuna di loro potrà produrre. L'esempio più semplice:
inviare su un channel senza buffer dalla stessa goroutine che dovrebbe
riceverne il valore.

\begin{attenzione}[Errore comune: errori nella gestione dei channel]
Il programma seguente termina \emph{volutamente} con un errore fatale. L'invio
\cod{ch <- 1} blocca finché qualcuno riceve, ma l'unica goroutine esistente è
proprio quella bloccata: nessuno riceverà mai. Il runtime di Go, quando
rileva che \emph{tutte} le goroutine sono bloccate, abortisce il programma.

\begin{lstlisting}[style=gocode,caption={\texttt{06\_concurrency/pitfall\_deadlock.go}}]
// pitfall_deadlock.go - ERRORE COMUNE: deadlock su channel senza buffer.
//
// ATTENZIONE: questo programma TERMINA VOLUTAMENTE con un errore fatale:
//
//	fatal error: all goroutines are asleep - deadlock!
//
// L'invio su un channel senza buffer blocca finché qualcuno riceve, ma
// l'unica goroutine esistente è quella bloccata: nessuno riceverà mai.
// Il runtime lo rileva e abortisce il programma (exit code 2).
package main

import "fmt"

func main() {
	ch := make(chan int)

	// ERRORE: invio bloccante nella stessa goroutine che dovrebbe ricevere.
	ch <- 1
	fmt.Println(<-ch) // mai raggiunto

	// Correzioni possibili: (1) channel con buffer make(chan int, 1);
	// (2) inviare da un'altra goroutine: go func() { ch <- 1 }().
}
\end{lstlisting}
\begin{lstlisting}[style=output]
fatal error: all goroutines are asleep - deadlock!

goroutine 1 [chan send]:
main.main()
        .../pitfall_deadlock.go:18 +0x36
exit status 2
\end{lstlisting}
Rimedi: un channel con buffer (\cod{make(chan int, 1)}) oppure inviare da
un'altra goroutine (\cod{go func() \{ ch <- 1 \}()}). Si noti che il
rilevamento automatico funziona solo quando \emph{tutte} le goroutine sono
bloccate: se ne esiste una ancora attiva, una goroutine bloccata resta
semplicemente ferma (\emph{goroutine leak}) senza alcun messaggio.
\end{attenzione}

\section{Riepilogo degli errori comuni nella concorrenza}

\begin{center}
\small
\begin{tabular}{@{}p{0.34\textwidth}p{0.58\textwidth}@{}}
\toprule
\textbf{Errore} & \textbf{Conseguenza e rimedio}\\
\midrule
\cod{main} termina prima delle goroutine & Le goroutine sono abbandonate.
Usare \cod{sync.WaitGroup} o un channel \cod{done}.\\[0.4em]
Accesso concorrente non protetto & Data race, aggiornamenti persi. Usare
\cod{sync.Mutex} o comunicare via channel.\\[0.4em]
\cod{wg.Add} dentro la goroutine & \cod{Wait} può ritornare troppo presto.
Chiamare \cod{Add} prima di \cod{go}.\\[0.4em]
Invio su channel chiuso & Panic. Chiudere solo dal mittente, una sola volta.\\[0.4em]
\cod{range} su channel mai chiuso & Blocco permanente. Il mittente deve fare
\cod{close}.\\[0.4em]
Invio/ricezione senza controparte & Deadlock. Usare buffer, un'altra
goroutine o \cod{select} con \cod{default}.\\[0.4em]
Assumere un ordine tra goroutine & Output variabile. L'ordine si ottiene
solo con sincronizzazione esplicita.\\
\bottomrule
\end{tabular}
\end{center}

% ===========================================================================
\chapter{Test e benchmark}\label{ch:test}
% ===========================================================================

\section{Il package \texttt{testing}}

Il package \cod{testing} della libreria standard fornisce gli strumenti per
scrivere \textbf{unit test}. Le convenzioni, interamente imposte dal toolchain:

\begin{itemize}
  \item i test stanno, di norma, in file il cui nome termina con
        \texttt{\_test.go} (per esempio \texttt{abs\_test.go}); questi file
        sono compilati solo da \cod{go test} e non finiscono nel programma;
  \item una funzione di test si chiama \cod{TestXxx}, dove \cod{Xxx} richiama il
        nome della funzione sotto test, e riceve un parametro
        \cod{*testing.T};
  \item i test si eseguono con \cod{go test}.
\end{itemize}

\begin{lstlisting}[style=gocode,numbers=none,caption={Il test della slide}]
func TestAbs(t *testing.T) {
	got := Abs(-1)
	if got != 1 {
		t.Errorf("Abs(-1) = %d; want 1", got)
	}
}
\end{lstlisting}

L'oggetto \cod{t} serve a segnalare i fallimenti: \cod{t.Errorf} marca il test
come fallito ma lo lascia proseguire, \cod{t.Fatalf} lo interrompe subito. Non
esistono \emph{assert}: si usa un normale \cod{if} e si segnala l'errore.

\subsection{Esempio pratico}

La cartella \texttt{07\_testing\_and\_tools} è un package ordinario (\cod{package
tools}), non un programma \cod{main}: contiene la funzione sotto test e i file
di test.

\begin{lstlisting}[style=gocode,caption={\texttt{07\_testing\_and\_tools/abs.go}}]
// Package tools contiene la funzione sotto test per gli esempi di
// testing e benchmarking (slide 71-72).
//
// Non è un package main: non si esegue con "go run", ma con
//
//	go test ./07_testing_and_tools/
//	go test -bench=. ./07_testing_and_tools/
package tools

// Abs restituisce il valore assoluto di n.
func Abs(n int) int {
	if n < 0 {
		return -n
	}
	return n
}

// FillAppend costruisce una slice di n elementi partendo da una slice
// vuota e usando append (può dover riallocare più volte).
func FillAppend(n int) []int {
	var s []int
	for i := 0; i < n; i++ {
		s = append(s, i)
	}
	return s
}

// FillMake costruisce la stessa slice allocando in anticipo la capacità
// con make: una sola allocazione.
func FillMake(n int) []int {
	s := make([]int, 0, n)
	for i := 0; i < n; i++ {
		s = append(s, i)
	}
	return s
}
\end{lstlisting}
\begin{lstlisting}[style=gocode,caption={\texttt{07\_testing\_and\_tools/abs\_test.go}}]
package tools

import "testing"

// I test stanno in file *_test.go; le funzioni si chiamano TestXxx e
// ricevono un *testing.T.
func TestAbs(t *testing.T) {
	got := Abs(-1)
	if got != 1 {
		// Errorf segna il test come fallito ma lo lascia proseguire.
		t.Errorf("Abs(-1) = %d; want 1", got)
	}
}

func TestAbsPositive(t *testing.T) {
	if got := Abs(5); got != 5 {
		t.Errorf("Abs(5) = %d; want 5", got)
	}
}

// Verifica che le due costruzioni producano lo stesso risultato.
func TestFillEquivalent(t *testing.T) {
	a, b := FillAppend(100), FillMake(100)
	if len(a) != len(b) {
		t.Fatalf("lunghezze diverse: %d vs %d", len(a), len(b))
	}
	for i := range a {
		if a[i] != b[i] {
			t.Fatalf("diverse in posizione %d", i)
		}
	}
}
\end{lstlisting}

Si esegue dalla cartella \texttt{examples/}:

\begin{lstlisting}[style=shell]
go test -v ./07_testing_and_tools/
\end{lstlisting}

\begin{lstlisting}[style=output]
=== RUN   TestAbs
--- PASS: TestAbs (0.00s)
=== RUN   TestAbsPositive
--- PASS: TestAbsPositive (0.00s)
=== RUN   TestFillEquivalent
--- PASS: TestFillEquivalent (0.00s)
PASS
ok  	goguide/07_testing_and_tools	(cached)
\end{lstlisting}

\subsubsection*{Analisi del codice}

\cod{go test} compila il package insieme ai file \cod{\_test.go}, esegue ogni
funzione \cod{TestXxx} e riporta \cod{PASS} o \cod{FAIL}. Con \cod{-v} si
vede ogni test; senza, solo il risultato finale (\cod{ok}). Si noti l'uso di
\cod{t.Fatalf} in \cod{TestFillEquivalent}: se le lunghezze differiscono è inutile
confrontare gli elementi, quindi si interrompe subito.

\begin{confronto}{C, Python, Java}
Go incorpora il framework di test nel toolchain: non servono dipendenze
esterne come \emph{JUnit} (Java), \emph{pytest}/\emph{unittest} (Python) o
\emph{CUnit}/\emph{Unity} (C). Le convenzioni sui nomi sostituiscono le
annotazioni e la configurazione.
\end{confronto}

\section{Benchmark}

Go fornisce anche strumenti per misurare le \textbf{prestazioni} del codice con i
\emph{benchmark}. I benchmark stanno nei file \texttt{\_test.go} e il loro nome
inizia con \cod{Benchmark}; ricevono un \cod{*testing.B}.

\begin{lstlisting}[style=gocode,numbers=none,caption={Il benchmark della slide}]
func BenchmarkXxx(b *testing.B) {
	for i := 0; i < b.N; i++ {
		Xxx(...)
	}
}
\end{lstlisting}

Il benchmark esegue la funzione \cod{b.N} volte per raccogliere dati sulle
prestazioni; il runner \emph{regola} \cod{b.N}, ripetendo l'esecuzione più volte
finché la misura è accurata. Si lanciano con \cod{go test -bench=.}.

\subsection{Esempio: \texttt{make} contro \texttt{append}}

Le slide propongono di confrontare la costruzione di una slice con \cod{append}
partendo da una slice vuota (che può dover riallocare più volte) con la
preallocazione della capacità tramite \cod{make}. Le funzioni misurate sono
\cod{FillAppend} e \cod{FillMake} nel file \texttt{abs.go} qui sopra.

\begin{lstlisting}[style=gocode,caption={\texttt{07\_testing\_and\_tools/benchmark\_test.go}}]
package tools

import "testing"

// sink è una variabile a livello di package: assegnarvi il risultato
// impedisce al compilatore di scartare o allocare sullo stack la slice
// (altrimenti, con inlining e dimensione costante, il benchmark di make
// risulterebbe senza allocazioni e il confronto sarebbe falsato).
var sink []int

// Benchmark: il runner sceglie b.N in modo adattivo, ripetendo il corpo
// finché la misura è stabile.
//
//	go test -bench=. ./07_testing_and_tools/

func BenchmarkFillAppend(b *testing.B) {
	for i := 0; i < b.N; i++ {
		sink = FillAppend(1000)
	}
}

func BenchmarkFillMake(b *testing.B) {
	for i := 0; i < b.N; i++ {
		sink = FillMake(1000)
	}
}
\end{lstlisting}

Esecuzione (macchina usata per la verifica: CPU Intel Core i7-13700H, Windows):

\begin{lstlisting}[style=shell]
go test -run NONE -bench "." -benchmem ./07_testing_and_tools/
\end{lstlisting}

\begin{lstlisting}[style=output,basicstyle=\ttfamily\footnotesize]
BenchmarkFillAppend-20   220239   5957 ns/op   25208 B/op   12 allocs/op
BenchmarkFillMake-20     513416   3160 ns/op    8192 B/op    1 allocs/op
\end{lstlisting}
(Le colonne, separate da tabulazioni nell'output originale, sono qui
allineate con spazi; i valori sono quelli misurati: iterazioni, tempo per
operazione, byte e numero di allocazioni per operazione.)

\subsubsection*{Analisi del risultato}

\begin{itemize}
  \item I tempi assoluti (\cod{ns/op}) \textbf{dipendono dalla macchina} e variano
        da esecuzione a esecuzione: i numeri sopra sono una misura reale ma non
        riproducibile identica altrove. Ciò che si ripete è il
        \emph{confronto}: la versione con \cod{make} è più veloce.
  \item La differenza si spiega con le allocazioni: \cod{append} su una slice
        vuota riallocava 12 volte (cresce a salti, copiando ogni volta
        l'array), occupando 25\,208 byte allocati in totale; \cod{make} con
        capacità nota alloca una sola volta, 8\,192 byte, esattamente lo spazio
        di 1000 \cod{int} da 8 byte.
  \item \textbf{Perché \cod{sink}?} Nella prima versione del benchmark il risultato
        veniva scartato: il compilatore, che inline-a funzioni piccole e conosce
        la dimensione costante, poteva allocare la slice sullo stack e il
        benchmark di \cod{make} riportava \emph{zero} allocazioni, falsando il
        confronto. Assegnando il risultato a una variabile di package
        (\cod{sink}) la slice «scappa» sullo heap e la misura è equa. È una
        lezione generale: un benchmark misura anche le ottimizzazioni del
        compilatore, quindi va scritto in modo da impedire che scartino il
        lavoro da misurare.
\end{itemize}

\begin{attenzione}[Errori comuni nei test e nei benchmark]
\begin{itemize}
  \item Il file non termina con \texttt{\_test.go}, oppure la funzione non
        inizia con \cod{Test}/\cod{Benchmark} (con la maiuscola): \cod{go test}
        la ignora silenziosamente e si crede di avere test che non esistono.
  \item Dimenticare di usare \cod{b.N} come numero di iterazioni nel ciclo:
        il benchmark misurerebbe una sola esecuzione.
  \item Confrontare tempi misurati su macchine diverse o con altro carico in
        esecuzione.
\end{itemize}
\end{attenzione}

\section{Argomenti citati ma non sviluppati}

L'ultima parte delle slide elenca argomenti \emph{ulteriori}, la cui trattazione
completa è rimandata ad altro materiale: \textbf{RPC in Go} (le slide rimandano
alle slide del corso sulla comunicazione e sulle RPC) e, come esempio dei
«molti altri aspetti da esplorare», il \textbf{supporto HTTP del package
\cod{net/http}}. Questa guida li menziona senza svilupparli, per restare nel
perimetro delle slide; di \cod{net/http} compare solo il tipo
\cod{http.ResponseWriter}\,/\,\cod{*http.Request} nell'esempio del middleware
(Sezione~\ref{sec:funzioni-middleware}). Moduli, funzioni variadiche, gestione
degli errori, strumenti e test, anch'essi elencati come «argomenti
aggiuntivi», sono invece trattati nei capitoli~\ref{ch:struttura},
\ref{ch:funzioni}, \ref{ch:oop} e in questo.

% ===========================================================================
\chapter*{Conclusione e riepilogo}
\addcontentsline{toc}{chapter}{Conclusione e riepilogo}
\markboth{Conclusione}{Conclusione}
% ===========================================================================

Go è un linguaggio piccolo, con poche regole applicate in modo rigido. Se si
guardano insieme i capitoli, le scelte di progetto si ricompongono in un disegno
coerente: \emph{semplicità} (un solo ciclo, nessuna classe, nessuna aritmetica sui
puntatori), \emph{sicurezza} (tipizzazione statica, garbage collector,
controllo dei limiti, valori zero) e \emph{scalabilità} (goroutine e channel
come cittadini di prima classe).

\section*{Riepilogo dei concetti principali}

\begin{center}
\small
\renewcommand{\arraystretch}{1.25}
\begin{tabular}{@{}p{0.26\textwidth}p{0.68\textwidth}@{}}
\toprule
\textbf{Concetto} & \textbf{Punto chiave da ricordare}\\
\midrule
Package e visibilità & Maiuscola iniziale $\Rightarrow$ esportato. Esecuzione da \cod{main.main}.\\
Toolchain & \cod{go run}, \cod{go build}, \cod{go install}; i moduli con \cod{go.mod}; formattazione con \cod{gofmt}.\\
Variabili & \cod{var} ovunque, \cod{:=} solo nelle funzioni; valore zero per ogni tipo.\\
Tipi & Conversioni sempre esplicite; inferenza dal lato destro.\\
Funzioni & Risultati multipli, risultati con nome, variadiche, anonime, closure che catturano variabili.\\
Controllo di flusso & Un solo \cod{for}; \cod{if} con istruzione iniziale; \cod{switch} senza fallthrough; \cod{defer} in ordine LIFO con argomenti valutati subito.\\
Puntatori & \cod{\&} e \cod{*}, nessuna aritmetica, \cod{nil}, sicuro restituire l'indirizzo di una locale.\\
Struct, array & Valori, copiati nelle assegnazioni; la lunghezza dell'array fa parte del tipo.\\
Slice & Vista \emph{(puntatore, len, cap)} su un array; \cod{append} va \emph{riassegnato}.\\
Map & Riferimento a tabella hash; chiave assente $\Rightarrow$ valore zero; usare \cod{v, ok}; ordine non specificato.\\
Metodi e interfacce & Receiver per valore o puntatore; interfacce \emph{implicite} basate sul comportamento.\\
Errori & Valori di tipo \cod{error}, ultimo risultato; \cod{if err != nil}.\\
Goroutine & Thread leggeri gestiti dal runtime; \cod{main} non le attende; nessun ordine garantito.\\
Channel & Code tipizzate e thread-safe; senza buffer sincronizzano; \cod{close} solo dal mittente; \cod{range} fino alla chiusura.\\
\cod{select} & Attesa su più channel; scelta casuale tra i pronti; \cod{default} per non bloccare.\\
Timer & Un canale \cod{C} che scatta dopo una durata; \cod{Stop()} per annullare.\\
Sincronizzazione & \cod{sync.Mutex}, \cod{sync.RWMutex}, \cod{sync.WaitGroup} (anche come barriera BSP).\\
Test e benchmark & File \texttt{\_test.go}, \cod{TestXxx(*testing.T)}, \cod{BenchmarkXxx(*testing.B)}; \cod{go test}.\\
\bottomrule
\end{tabular}
\end{center}

\section*{Per proseguire}

Le slide chiudono ricordando che questa è \emph{solo un'introduzione}: restano da
esplorare, tra l'altro, le RPC in Go e il supporto HTTP del package
\cod{net/http}. Il modo più efficace per consolidare quanto appreso è
modificare gli esempi della cartella \texttt{examples/}: cambiare le capacità
delle slice e osservare \cod{append}, togliere un \cod{close} e provocare un
deadlock, rimuovere il mutex dal contatore e osservare la data race.

% ===========================================================================
\appendix
\chapter{Mappa tra slide, guida ed esempi}
% ===========================================================================

La tabella associa i numeri delle slide (quelli stampati sulle slide)
ai capitoli della guida e ai file in \texttt{examples/}. Le slide 1--2 contengono
solo cenni storici e classifiche (omessi) e la slide 73 solo riferimenti
bibliografici e web (omessa).

\begin{center}
\footnotesize
\renewcommand{\arraystretch}{1.15}
\begin{longtable}{@{}L{1.4cm}L{4.9cm}L{1.2cm}L{6.1cm}@{}}
\toprule
\textbf{Slide} & \textbf{Argomento} & \textbf{Cap.} & \textbf{Esempi}\\
\midrule
\endhead
1--5 & Obiettivi, eredità da C, novità, sistemi distribuiti & \ref{ch:go} & ---\\
6 & Go per applicazioni cloud-native (Go Cloud) & \ref{ch:go} & frammento non eseguibile\\
7 & Editor, plugin, \cod{gopls} & \ref{ch:go} & ---\\
8--9 & Hello world, package, \cod{main}, \cod{fmt} & \ref{ch:struttura} & \texttt{helloworld.go}\\
10 & Toolchain: \cod{go run/build/install} & \ref{ch:struttura} & ---\\
11 & Package e visibilità & \ref{ch:struttura} & \texttt{imports.go}\\
12 & Import & \ref{ch:struttura} & \texttt{imports.go}\\
13--15 & Funzioni, risultati multipli, errori & \ref{ch:funzioni} & \texttt{functions.go}, \texttt{swap.go}, \texttt{result\_and\_error.go}\\
16 & Valori di ritorno con nome & \ref{ch:funzioni} & \texttt{named\_returns.go}\\
17 & Variabili & \ref{ch:tipi} & \texttt{variables.go}\\
18 & Tipi, conversioni, inferenza & \ref{ch:tipi} & \texttt{types.go}\\
19--22 & Controllo di flusso, \cod{for}, \cod{echo} & \ref{ch:flusso} & \texttt{for\_loops.go}, \texttt{echo.go}\\
23--24 & \cod{if}/\cod{else} & \ref{ch:flusso} & \texttt{if\_else.go}\\
25 & \cod{switch} & \ref{ch:flusso} & \texttt{switch.go}\\
26 & \cod{defer} & \ref{ch:flusso} & \texttt{defer.go}, \texttt{defer\_lifo.go}\\
27 & Puntatori & \ref{ch:dati} & \texttt{pointers.go}\\
28 & Struct e array & \ref{ch:dati} & \texttt{structs.go}, \texttt{arrays.go}\\
29--31 & Slice, \cod{len}/\cod{cap}, \cod{make} & \ref{ch:dati} & \texttt{slices.go}, \texttt{making\_slice.go}\\
32--33 & \cod{append} & \ref{ch:dati} & \texttt{append.go}, \texttt{slice\_pitfalls.go}\\
34 & Map & \ref{ch:dati} & \texttt{maps.go}\\
35--36 & \cod{range} & \ref{ch:dati} & \texttt{range.go}, \texttt{range2.go}\\
37--39 & Funzioni anonime, closure & \ref{ch:funzioni} & \texttt{anonymous.go}, \texttt{closure.go}, \texttt{closure\_fib.go}\\
40 & Closure come middleware & \ref{ch:funzioni} & \texttt{middleware.go}\\
41 & Metodi & \ref{ch:oop} & \texttt{methods.go}\\
42--45 & Interfacce & \ref{ch:oop} & \texttt{interface.go}, \texttt{interface\_slice.go}\\
46--48 & Concorrenza, goroutine, scheduler & \ref{ch:conc} & \texttt{goroutine\_termination.go}, \texttt{goroutine\_wait.go}, \texttt{gomaxprocs.go}\\
49--52 & Channel & \ref{ch:conc} & \texttt{channel\_sync.go}, \texttt{channel\_sum.go}, \texttt{channel\_fib.go}\\
53 & Channel con/senza buffer & \ref{ch:conc} & \texttt{channel\_buffered.go}\\
54--55 & \cod{close}, \cod{range} su channel & \ref{ch:conc} & \texttt{channel\_close.go}, \texttt{channel\_range.go}, \texttt{pitfall\_send\_closed.go}\\
56--59 & \cod{select}, \cod{default} & \ref{ch:conc} & \texttt{channel\_fib\_select.go}, \texttt{select\_multiplex.go}, \texttt{select\_default.go}\\
60 & Timer & \ref{ch:conc} & \texttt{timers.go}\\
61--63 & Mutex con channel, \cod{sync.Mutex}, \cod{RWMutex} & \ref{ch:conc} & \texttt{mutex\_using\_channel.go}, \texttt{mutex.go}, \texttt{race\_counter.go}, \texttt{rwmutex.go}\\
64 & \cod{WaitGroup}, BSP & \ref{ch:conc} & \texttt{waitgroup.go}, \texttt{bsp\_barrier.go}\\
65 & Priorità & \ref{ch:conc} & \texttt{priority.go}\\
66 & Argomenti aggiuntivi (RPC, HTTP: citati) & \ref{ch:test} & ---\\
67 & Moduli & \ref{ch:struttura} & \texttt{go.mod}\\
68 & Funzioni variadiche & \ref{ch:funzioni} & \texttt{variadic.go}\\
69 & Gestione degli errori & \ref{ch:oop} & \texttt{errors.go}\\
70 & Regole del compilatore, \cod{gofmt}, \cod{goimports}, \cod{godoc} & \ref{ch:struttura} & ---\\
71--72 & Test e benchmark & \ref{ch:test} & \texttt{abs.go}, \texttt{abs\_test.go}, \texttt{benchmark\_test.go}\\
\bottomrule
\end{longtable}
\end{center}

Esempi aggiuntivi non presenti nelle slide, scritti per chiarire concetti già
trattati: \texttt{interface\_slice.go}, \texttt{channel\_sync.go},
\texttt{select\_multiplex.go}, \texttt{bsp\_barrier.go}, \texttt{priority.go},
\texttt{race\_counter.go}, \texttt{slice\_pitfalls.go}, \texttt{gomaxprocs.go},
\texttt{defer\_lifo.go}, \texttt{anonymous.go}, \texttt{pitfall\_deadlock.go},
\texttt{pitfall\_send\_closed.go}, oltre agli esempi sugli errori comuni
inseriti negli altri file.

% ===========================================================================
\chapter{Note della verifica incrociata}
% ===========================================================================

\section{Metodo}

La verifica ha confrontato quattro elementi: le slide originali, il testo di
questa guida, i file \texttt{.go} e gli output. In concreto:

\begin{itemize}
  \item tutte le 73 slide sono state lette prima di scrivere, e ogni argomento
        tecnico è stato associato a una sezione (Appendice~A);
  \item i listati della guida sono inclusi \emph{dai file reali}
        (\texttt{examples/}), quindi la corrispondenza è esatta per costruzione;
  \item gli output mostrati sono stati ottenuti eseguendo i programmi con Go
        1.27.3; quelli non deterministici sono stati eseguiti più volte e
        indicati come tali;
  \item ogni file è stato formattato con \cod{gofmt} e controllato con
        \cod{go vet}.
\end{itemize}

\section{Errori e imprecisioni trovati nelle slide}

Per ciascun punto si distingue il \textbf{problema di trascrizione} (errore
materiale nella slide) dalla \textbf{correzione o precisazione tecnica}
introdotta dalla guida.

\begin{center}
\footnotesize
\renewcommand{\arraystretch}{1.2}
\begin{longtable}{@{}L{0.5cm}L{1.2cm}L{5.5cm}L{6.4cm}@{}}
\toprule
\textbf{\#} & \textbf{Slide} & \textbf{Osservazione} & \textbf{Trattamento nella guida}\\
\midrule
\endhead
1 & 58 & \emph{Trascrizione.} \cod{fibonacci(c, quit)} usa la variabile \cod{c}, mai dichiarata (si dichiara \cod{ch}); non compila (\texttt{undefined: c}). &
Corretto in \cod{fibonacci(ch, quit)} in \texttt{channel\_fib\_select.go}; segnalato nel Capitolo~\ref{ch:conc}.\\
2 & 56 & \emph{Trascrizione.} \cod{case mgs1 := <-ch1} è un refuso per \cod{msg1}. &
Usato \cod{msg1}; segnalato nel Capitolo~\ref{ch:conc}.\\
3 & 48 & \emph{Precisazione tecnica.} «Set \cod{runtime.GOMAXPROCS} environment variable» mescola due cose: la variabile d'ambiente si chiama \cod{GOMAXPROCS}; \cod{runtime.GOMAXPROCS(n)} è la funzione. &
Distinte le due forme (Sez.~\ref{sec:goroutine}); nessuna modifica al concetto.\\
4 & 11 & \emph{Aggiornamento.} Il messaggio d'errore per \cod{math.pi} oggi è \texttt{undefined: math.pi (but have Pi)}; la slide riporta la forma più breve. &
Riportato il messaggio ottenuto con Go 1.27.3.\\
5 & 36 & \emph{Precisazione.} L'output \cod{a -> apple}, \cod{b -> banana} non è garantito: l'ordine di iterazione di una map è non specificato. &
L'esempio itera su un elenco di chiavi ordinato a mano; spiegato nel Capitolo~\ref{ch:dati}.\\
6 & 33 & \emph{Precisazione.} «Di norma il doppio della capacità precedente»: nell'esecuzione reale \cod{append(s, 2, 3, 4)} su una slice di capacità 2 dà capacità \textbf{6}, non 5 né 4. &
Discusso in \texttt{append.go} e nel Capitolo~\ref{ch:dati}: la regola è un dettaglio d'implementazione.\\
7 & 50, 53 & \emph{Precisazione.} «Invio e ricezione bloccano finché l'altro lato è pronto» vale solo per channel senza buffer. &
Chiarito con un box nella Sezione~\ref{sec:buffered}.\\
8 & 15 & \emph{Dipendenza dall'ambiente.} \cod{os.ReadFile("example.txt")} presuppone un file esistente. &
Esempio reso autonomo creando il file in una directory temporanea.\\
9 & 40 & \emph{Frammento incompleto.} \cod{timed} è mostrato senza \cod{hello}, imports e server. &
Reso eseguibile con \cod{httptest}, senza rete.\\
10 & 62 & \emph{Segnalazione di \cod{go vet}.} Il codice compila e funziona, ma vet segnala \texttt{return copies lock value} su \cod{NewLock}. &
Codice mantenuto come nella slide; il warning è spiegato (Sez.~\ref{sec:mutex}).\\
11 & 30 & \emph{Precisazione.} «Compile or run-time error if array length is exceeded»: l'errore è a compile time solo per indici \emph{costanti} su array; altrimenti è un panic a runtime. &
Entrambi i casi mostrati con i messaggi reali (Cap.~\ref{ch:dati}).\\
12 & 48 & \emph{Dato non verificabile qui.} Le stime «circa $5\times$ più veloci, $500\times$ meno memoria» non sono state misurate. &
Riportate come \emph{dato delle slide}, non verificato.\\
\bottomrule
\end{longtable}
\end{center}

\section{Contenuto delle slide deliberatamente omesso}

Come richiesto, non compaiono: la storia del linguaggio (anno di concezione e
autori, slide 1), le classifiche di popolarità (slide 2), l'elenco dei
riferimenti bibliografici e web (slide 73) e i collegamenti a siti esterni. Il
frammento Go Cloud (slide 6) è riportato come testo ma non ha un file
eseguibile perché dipende da una libreria esterna e da servizi cloud.

\section{Verifica degli esempi}

\begin{itemize}
  \item Tutti i programmi della cartella \texttt{examples/} compilano ed
        eseguono con Go 1.27.3 (Windows, amd64).
  \item Due programmi terminano \emph{volutamente} con un errore, e sono
        marcati nel codice e nel testo: \texttt{pitfall\_deadlock.go} (deadlock,
        exit code 2) e \texttt{pitfall\_send\_closed.go} (panic, exit code 2).
  \item Gli esempi con output non deterministico sono:
        \texttt{goroutine\_termination.go}, \texttt{goroutine\_wait.go},
        \texttt{select\_multiplex.go}, \texttt{race\_counter.go},
        l'intercalarsi delle righe in \texttt{channel\_close.go} e i tempi di
        \texttt{middleware.go} e dei benchmark.
  \item Non sono stati eseguiti con il \emph{race detector} (\cod{-race}), che
        richiede il supporto di CGO non disponibile nell'ambiente di verifica.
        La data race di \texttt{race\_counter.go} è dimostrata dal risultato
        errato, non dallo strumento.
\end{itemize}

\end{document}

